\part{Runic Poetry}

\bookStart{Runic Poetry}

When completed, this section will contain an introduction to runes, the three rune poems, along with most poetry written in runes.  Note that magical inscriptions will be edited under Galders, and the Runic version of the \emph{Dream of the Rood} under Poetry on Christian Subjects.

Metrically, most of the poetry is in the alliterative long line (Norse \Fornyrdislag).  A few fragments from Jutland are in \Ljodahattr; two from Sweden, and several from Norway are in \Drottkvett.

\bookStart{Introduction to Runes and Rune Poems}

\section{Runes and other European alphabets}

The \emph{runes} are the earliest recorded indigenous script for writing the Germanic languages, and as such are of great cultural, linguistic and historical significance.  There survive runic inscriptions from all three Germanic branches (East, West, and North Germanic), and the script remained in unbroken use by Germanic speakers for a period spanning almost twenty centuries.  The earliest reliably dated inscription is the stone from Hole, Norway, most likely from the second century \ce\ \parencites{ZilmerVasshus2023}{SolheimEtAl2025}, and as late as the beginning of the 20th century, peasants in Dalarna, Sweden, were still using a genuine, generationally transmitted form of runes to write short messages \parencites{Jans2015}{Fridell2025}.

Like the Latin or Greek scripts, the runes are an alphabet, with each symbol representing a phoneme (a vowel or consonant), rather than a whole syllable or word.  As seen by their general structure, the runes must ultimately derive from the Greek alphabet, which in turn derives from the Phoenician \emph{abjad} (a writing system with no dedicated vowel-symbols), in turn an adaptation of the Egyptian Hieroglyphs.  The runes cannot be directly derived from the Greek, however, since the runes ᚠ \textbf{f}, ᚱ \textbf{r}, and ᚺ \textbf{h} show distinct innovations shared with the Latin alphabet, a representative of Italic writing; on the other hand, Greek digamma (\textgreek{Ϝ}) stood for \emph{w}, not \emph{f}, Greek rho (\textgreek{Ρ}) lacked the stroke seen in Latin R and Runic ᚱ, and Greek eta (\textgreek{Η}) stood for long \emph{ē}, not the consonant \emph{h}.  Needless to say, the existence of numerous Runic vowels corresponding to Greek and Latin precludes a direct Phoenician origin for the runes.%textrunic

Although the runes are the oldest writing system common to the Germanic peoples, the oldest Germanic words are not written in Runic at all, but rather in the Etruscan alphabet.  The inscription in question is found on the so-called Negau B helmet, one of numerous bronze helmets found in 1812 in Negau, Styria (present-day Negova in northeastern Slovenia), likely dating to the 4th century \bce.  The full inscription, the original of which is written right-to-left, reads in transliteration \emph{hariχastiteiva \textbackslash\textbackslash\textbackslash ip}.  The final sequence \emph{ip}, preceded by three diagonal strokes, is not understood, but the first 9 letters are nigh universally understood to reflect the Proto-Germanic dithematic name \emph{*Hari-gastiz} ‘Army-Guest’, formed from \emph{*harjaz} ‘host, army’ + \emph{*gastiz} ‘guest’.  The remaining 5 letters, \emph{teiva}, have been convincingly explained as an archaic form of an unattested PGmc.~m.~\emph{an}-stem noun \emph{*tīwô} ‘priest’, formed from \emph{*tīwaz} ‘god, tew’ in the same way as \emph{*gudjô} ‘priest’ from \emph{*gudą} ‘god’.  The importance of the Negau B helmet for the present discussion is that it strongly supports a North Italic origin for the runes.  The script on the helmet itself, however, can hardly be the direct precursor to Runic; its letter \emph{r} resembles the Greek rho, lacking the diagonal stroke found in Latin and later Runic.
TODO: cite \textcite{Markey2001}.

It is clear that the runes cannot be borrowed from the Greek or Etruscan, but could they have been adopted directly from the Romans?  Although plausible on its face, a few archaic traits of early runic writing seem to point against such an origin.  In particular, the earliest runic inscriptions are frequently written right-to-left or in boustrophedon, and the shapes of ᚨ \textbf{a} and certain forms of ᛊ \textbf{s} (written with numerous strokes, e.g., on the Hole stone, \cite{ZilmerVasshus2023}) and ᚺ \textbf{h} (often written with several diagonal branches, e.g., ᚻ) seem closer to Old Italic alphabets than to Latin A, S, or H.  The most likely explanation, then, is probably that the direct source for the runes was an Italic or Alpine script with heavy influence from standard Latin, used some time during the centuries immediately before the birth of Christ.

\section{Original innovations}

If the source for the runes is most likely to be sought in Italy or the Alps, it is next important to discuss its relation to the runes.  We here ought to distinguish between direct borrowing and inspiration.  What makes the runes particularly notable is that they do not seem to have been a mere borrowing of a complete script, only gradually adapted to fit the phonology of the target language (as would happen later, when the Latin alphabet was adapted to the Germanic languages through the gradual integration of runic letters, diacritics, and the like).  Instead, already from the oldest runic inscriptions, we find an alphabet that is fully adapted to the phonology of the early Germanic language.  This includes characters for the dental fricative \emph{þ} and the voiced velar nasal \emph{ŋ}, and, even more notably from a comparative paleographic perspective, a dedicated symbol for the semivowel \emph{j} distinct from for the vowel \emph{i}.

The single most unique trait of the runic script, however, is the order and names of its letters, which distinguish the runes from all other Phoenician-derived abjads and alphabets, whether Semitic, Greek, or Italic.  The first letters of the Phoenician abjad are \texthebrew{𐤀𐤁𐤂𐤃}, and this was followed by the Greek alphabet with \textgreek{ΑΒΓΔ} and the Italic scripts like Latin and Etruscan with ABCD.  The first letters of the Runic alphabets, on the other hand, are ᚠᚢᚦᚨᚱᚲ \textbf{fuþark}, which is why these scripts are often referred to as \emph{futharks}.  The runes also do away with the originally Semitic letter-names evinced in the Greek Alpha, Beta, Gamma, Delta (< Phoenician \emph{alp} ‘ox’, \emph{bet} ‘house’, \emph{gaml} ‘throwing stick’, \emph{dalt} ‘door’).  Instead, each rune has a native Germanic name, typically a concrete noun, chosen so that its initial (for ᛜ \textbf{ŋ} its medial, for ᛉ \textbf{z} its final) sound should correspond to that rune’s sound-value.  For instance, the various Old Germanic names attested for the rune ᚠ \textbf{f} always reflect descendants of the Proto-Germanic noun \emph{*fehu} ‘cattle’.  The principle that the name of each rune should begin with the sound-value which it stands for, what runologists call the \emph{acrophonic principle}, has important consequences in the later development of runic writing as the Germanic languages grew apart.

All in all, with its phonemically informed structure, its division into three groups of eight runes each, its unique order, and its deliberately chosen names, the original \emph{Elder Futhark} gives a highly intentional impression.  It seems likely that its inventor, having seen the advantages of literacy, consciously transformed his non-Germanic source script to fit his native Proto-Germanic language after his own phonemic analysis.  In doing so, he established a tradition that would last for almost two millennia amongst his ethnic kin.  It is of note that the runes were extremely seldom used to write a non-Germanic language and, even then, only within a Germanic cultural context.  Old English speakers would at times write short Irish or Latin texts in runes, but the Irish themselves never adopted runes, nor did the Slavs, Finns, or Saami.  If this script spread widely, it only did so \emph{alongside} the rapid expansion of the Germanic peoples and their languages, but it was never to be divorced from its foundational linguistic context.

\section{Order and structure}

Before turning to the rune poems and the later developments of the futhark, let us briefly look at the order and structure of the Elder Futhark, the ancestral Runic script, attested from the 2nd century onwards.

The Elder Futhark originally consists of 24 runes: 6 vowels, 2 semi-vowels, and 16 consonants.  The earliest attestation of the complete rune-row in sequence is on the Kylver stone (G 88) from Gotland, archaeologically dated to ca.~400 \ce.%TODO: cite Spurkland2005
The flat limestone slab was found text-down as the lid to a grave; its full inscription reads \textbf{fuþarkgwhnijpïzstbemlŋdo ¶ sueus}, where \textbf{sueus} is an unexplained string written outside of the futhark-sequence.

A somewhat younger complete rune-row is found on two Swedish rune-inscribed bracteates (Ög 178 Vadstena and Nä 10 Mariedamm, IK 377.1--2) dated to ca.~500 \ce.  Both bracteates are struck from the same die and contain the full inscription \textbf{tuwatuwa : fuþarkgw : hnijïpzs : tbemlŋod}.  The sequence \textbf{tuwatuwa} is unexplained, but is clearly not part of the futhark.  Another full sequence is found on the contemporary Swedish Grumpan bracteate (Vg 207, IK 260), which reads \textbf{fuþarkgw ········ hnijïp{-}{-} ···· tbemlŋod}, where the dots serve to fill out the full space of the bracteate’s rim.  Apart from the loss of the two runes following the \textbf{p}, IK 260 contains the full rune-row in the same order as IK 377.1--2; like them, IK 260 gives the runes \textbf{ïp} and \textbf{od} in reverse order from G 88, but otherwise adheres to that order.  A heavily corrupt rune-row is found on the contemporary Lindkær bracteate (DR IK110, IK 110): \textbf{fuþarkgwhnelat-(t)suao-(u)}.

The two complete bracteate-sequences (IK 377.1--2 and IK 260) are especially important since they clearly divide the 24 runes into three groups of eight, a categorisation retained in the later Younger Futhark.  These groups are eventually named \emph{ę́ttir} (‘groups of eight’, by folk etymology reinterpreted ‘family’) in Old Icelandic sources, but the name must be far older, since already by the 9th century the Younger Futhark only contained 16 runes, and the three groups were thus no longer groups of eight.

The names of the Elder Futhark, as far as we can reconstruct them, show some intentionality.  In particularly, the first four runes seem to form semantic antonyms: \emph{*fehu} ‘fee, cattle’ : \emph{*ūruz} ‘aurochs’; \emph{*þurisaz} ‘thurse, dæmon’ : \emph{*ansuz} ‘os, god’.  After this, the pattern seems to break down, as if the namer of the alphabet struggled to find suitable words, but in the second eighth we find the sequence \textbf{hnij}: \emph{*hagalaz} ‘hail’, \emph{*naudiz} ‘need, famine’, \emph{īsaz} ‘ice’ : \emph{*jērą} ‘year, good harvest’, where three nouns describing the struggle of winter are followed by one describing the joys of spring.

\section{Rune poems}

If the order and names of the letters of the futharks stayed relatively consistent throughout the many centuries and countries in which they were used, especially in the case of the first six letters, this faithfulness should probably be ascribed to the \emph{rune poems}.  These poems, poetic lists of the rune-names with a short explanation of each, would have been passed down orally as mnemonic devices to aid early Germanic learners of writing, who were doubtless far more accustomed to learn poems by heart than written symbols, and who would have profited much from such poetry in their gradual ascent to literacy.

In fact, with the simple aid of a correctly written \emph{futhark}-amulet (and many such survive) and a memorised rune poem of some 24 lines (laughably low in an oral-mnemonic culture), even an illiterate Iron Age tribesman could, with some effort, decipher any Elder Futhark inscription.  He need simply recite the poem while moving his hand over the runes one by one until he found the match between the shape and the word he was pronouncing.  The original use for the runes would likely have been as a military technology used for coded messages, and in order to learn them in the first place one thus needed to be taught both the rune-poem by someone who already knew it and be given a futhark-amulet.  Little wonder, then, that the runes were soon attributed to the gods themselves (cf.~Vg 63 Noleby (under Runic Poetry), \textlink{Havamal}[80], 138--146).

There survive three such rune poems, from three countries: England, Norway, and Iceland.  The Old English Rune Poem (\textlink{EnglishRunePoem}) documents the Anglo-Saxon (or properly, Anglo-Frisian) \emph{Futhorc}, while the Norwegian and Icelandic Rune Poems (\textlink{NorwegianRunePoem}, \textlink{IcelandicRunePoem}) document the Scandinavian \emph{Younger Futhark} (YF).  It is useful to discuss the differences between these two scripts and their innovations from the Elder Futhark (EF) here.

\section{The later futharks}

When compared to the Common Germanic Elder Futhark, the two daughter scripts took entirely opposite paths as their languages grew apart.  The Futhorc is the more conservative elder daughter, for while it has appended several letters for new vowels to the end of the rune row, it otherwise retains all of the runes from the Elder Futhark (albeit with some innovations in their shapes and sound-values).

The rebellious Scandinavian Younger Futhark, on the other hand, has done away with numerous runes, namely the voiced velar nasal \textbf{ŋ}, the plosives \textbf{d, g, p}, the semi-vowel \textbf{w}, the vowels \textbf{o} and \textbf{e}, and the obscure hook-shaped vowel \textbf{ï}.  That much of this simplification was probably intentional, rather than the result of neglect or language change, is suggested by the following observations:

Firstly, while the transition to Old Norse that took place at roughly the same time as the transition to the Yougner Futhark was accompanied by rapid language change, this does not explain the loss of characters.  Several of the lost runes stood for sounds that did not undergo any major sound shifts in the North Germanic languages in the relevant time period.  In fact, all modern Scandinavian dialects still distinguish between the initial consonants in the descendants of ON \emph{dagʀ} ‘day’ and \emph{Týʀ} ‘\inx[P]{Tew}’, and these consonants generally have the same articulation as their equivalent sounds in Modern English or Low German, suggesting that they are very conservative sounds.  On the other hand, many of the runes whose names did undergo large changes, e.g., EF \textbf{j} (Proto-Norse \emph{*jāra} ‘year, good harvest’) > YF \textbf{a} (Old Norse \emph{ár} ‘id.’), were retained, albeit with new sound-values based on the acrophonic principle.

Secondly, we find clear proof in two early YF inscriptions that the names and sound values of some of the lost runes were still remembered and passed down even after the adoption of the simplified Younger Futhark, which precludes a gradual and organic dropping of these letters in response to sound change.  On the Swedish Rök stone (Ög 136), which is mostly composed in the Younger Futhark, runes of the Elder Futhark are employed in a peculiar cipher which works as follows (cf. \cite{FridellWilliams2021}): for every YF rune representing two or more phonemes, where one of said phonemes had its own distinct rune in the EF system whose sound-value was assimilated by the surviving rune, one of the lost runes is substituted for that surviving rune.  This results in a contrafactual form of the YF, where the EF runes that were lost in the YF replace the runes that replaced them.
For instance, the sounds /t/ and /d/ are represented in the EF by the runes \textbf{t} and \textbf{d}, respectively.  In the YF the old \textbf{d}-rune was lost, and so \textbf{t} fills the double duty of /t/ and /d/.  In the cipher, however, it is the old \textbf{d}-rune which fills the double duty of /t/ and /d/ instead.  A similar instance of two-scriptedness is found on the Ingelsta stone (Ög 43), where the old \textbf{d} rune is used ideographically in a YF inscription in order to stand for its name \emph{*dagʀ} ‘day’ (cf.~OE \emph{dæġ} ‘id.’), which is also attested as a male given name.

Third, there is virtually no regional variation in which runes disappear in the transition from EF to YF.  There is some variation in the shapes of the runes, e.g., early Danish YF inscriptions retain the old shapes ᚺ \textbf{h} and ᛗ \textbf{m} for younger ᚼ and ᛘ; but there is no region which, say, simplifies only the plosive consonants \textbf{t \char`~\ d, k \char`~\ g, b \char`~\ p} > \textbf{t, k, b}, but retains the written distinction between \textbf{o} and \textbf{u} or the voiced velar nasal \textbf{ŋ}---rather, the plosives and additional vowels all go away at once.

These observations point away from neglect or a natural development of the script, and instead suggest deliberate reform.  Since we lack historical sources, the motivations behind such a reform can only be guessed at, but making the script simpler may have been intended to increase literacy by making it easier to learn and faster to write.  If this were the case, it was certainly successful, for the transition to the simplified Younger Futhark brings with it a huge increase in inscriptions in Scandinavia, a greater complexity in the length and types of inscriptions, an interest in various cipher-techniques (especially seen on the Rök stone, which makes use of four), and a new tradition of inscribed stones in Denmark, where they were previously unknown.

This new system also quickly gives rise to several even more simplified scripts, like the “short-stave” runes found on the Rök stone, or the “staveless” runes known from northern Sweden.  Both of these variants make the runes even easier to write or carve on materials like wood, wax and bone; the runes also take up less space---very useful for carvers working with limited surfaces.

\section{Comparison of names}

In any case, and importantly for us when we deal with the rune poems, the names of the runes generally seem to have survived the linguistic and orthographic developments into both Old Norse and Old English.  In order to arrive at a likely original list of names for the runes of the Elder Futhark, and in order to contextualise the surviving rune poems, there is some reason to list the attestations of the runic names outside of the rune poems here.

\subsection{Gothic names}

Before dealing with the two futharks, I would like to begin with the names for the letters of the Gothic alphabet.  Although only preserved once, and in quite a poor state, these names appear to have been adapted from the form of the EF used by the Goths (as most famously attested on the Pietroassa ring; \cite[280]{Antonsen2002}[173--174]{MacLeod2006}).  Although the Gothic alphabet is most similar to the Greek, runic influence on it can also be seen in the existence of a separate symbol for the sound \emph{j}, seemingly apparently adapted directly from the runic ᛃ \textbf{j}.

The names for the Gothic letters survive only in Codex Vindobonensis 795 (fol.~20v), a manuscript consisting of treatises by Alcuin compiled around 798 \ce.  On the reverse of the same folio is found a list of the names of the Anglo-Frisian Futhorc; see below.

TODO: Tabulate.  \emph{aza, bercna, geuua, daaz, eyz, quertra, ezec, haal, thyth, iiz, chozma, laaz, manna, noicz, gaar, uraz, pertra, reda, sugil, tyz, uuinne, fe, enguz, uuaer, utal}.

\subsection{The Anglo-Frisian Futhorc}

The names of the Anglo-Frisian futhorc are preserved in four manuscript sources separate from the Old English Rune Poem (\textlink{EnglishRunePoem}).  They are Codex Sangallensis 270 (p.~52), Codex Vindobonensis 795 (fol.~20r), Cotton Domitian A. IX (fol.~11v), and Oxford St.~John’s College MS 17 (fol.~5v).

The runes in the following table are listed by their order in \textlink{EnglishRunePoem}; their order in the relevant manuscript is indicated by the number preceding the name (e.g., “1. \emph{feh}” under Oxford St. John's College MS 17 means that the \textbf{f}-rune appears under the name \emph{feh} as the first rune in that manuscript’s list).  Distinctions of shapes are not noted.

TODO: Tabulate.

Cod. Sang. 270: \url{https://www.e-codices.unifr.ch/en/csg/0270/51/0/}

Cod. Vind. 795: \url{https://viewer.onb.ac.at/1002EC92}

Cotton Domitian A. IX: \url{https://iiif.bl.uk/uv/#?manifest=https://bl.digirati.io/iiif/ark:/81055/vdc_100060430962.0x000001}

St. John’s College 17: \url{https://digital.bodleian.ox.ac.uk/objects/66a78997-ab65-4059-a9d3-d08a0bba067c/surfaces/a047e779-5715-4047-896c-5f92217a379c/}

\subsection{Younger Futhark}

The earliest sources for the rune-names of the Younger Futhark are in Frankish manuscripts, which preserve them as a sort of Northern curiosity.  Of particular note are the \emph{Abecedarium Nordmannicum} and the runes in the Codex Leidensis.

The \emph{Abecedarium Nordmannicum} contains a strange mix of Old High German, Old Saxon, and Old Norse.

TODO: Codex Leidensis \parencite{Myrvoll2026}.

These names show several archaic linguistic forms, as have been discussed by \textcite{Myrvoll2026}.  Of particular note is the preservation of Proto-Norse short unstressed \emph{-u} (< PGmc.~\emph{*-ų} or \emph{*-ō}) following short syllables and its analogical substitution following long syllables; this is seen in the names TODO.

14 of these names may be readily identified with their forms in the later Icelandic rune poem.  The two that remain are \textbf{ląukʀ} and \textbf{ir}.  On the surface, \textbf{ląukr} suffers from double ambiguity: the sequence \textbf{ąu} can at this time stand either for \emph{au} or \emph{ǫ}, and the rune \textbf{k} can stand for both \emph{k} and \emph{g}.  \textbf{ląukʀ} may thus stand for either \emph{*lǫgʀ} ‘water, lake’, its name in the later Icelandic and Norwegian Rune Poems, along with the Old English, or for the otherwise unparallelled \emph{*laukʀ} ‘leek’.
Yet in spite of the comparative arguments, the latter is in fact the only plausible option, for two reasons: first, the transliteration transcribes the name \emph{laucr}, not \emph{**laugr} or \emph{**logr}; second, the \emph{*-u-} inherited from Proto-Norse \emph{*laguʀ} ‘water, lake’ would undoubtedly have been preserved at this stage of the language, as seen by the names \textbf{raiþu} and \textbf{sulu}.%TODO.

\textbf{ir} is more difficult; from its order in the sequence it appears to be an Old Danish form of the West Norse \emph{ýr} ‘yew-tree’.  This is the old \textbf{ʀ}-rune, but we would then expect its name to contain the sound \textbf{ʀ} at the end, not \textbf{r}.  This discrepancy should probably be understood as a confusion between \textbf{r} and \textbf{ʀ} on the part of the Danish informant, likely alongside the spelling of the \textbf{a}-rune as \textbf{ąs}; that the informant was influenced by Continental Germanic scribal norms is in any case clear by his spelling of the name \emph{þurs} as \textbf{þhurs}, clearly due to influence from the digraph <th>.

TODO: Table of rune-names.

Of the 16 runes found in both the English and Icelandic (which appears to be more conservative than the Norwegian) Rune Poems, 10 have etymologically identical names (\textbf{f, r, h, n, i, j, s, b, m,} and \textbf{l}).  3 of the 6 divergent names (\textbf{þ, a,} and \textbf{t}) are words with clear Heathen associations in the Icelandic poem---Thurse, Os, and Tew---and so may have been changed deliberately after the conversion of England, rather than lost in the process of oral transmission.  Two more---\textbf{u} and \textbf{k}---have names which agree in form but not in meaning.  Thus it is only for the old \textbf{z}-rune where there is complete disagreement about the original name.  This is easily understood, however, since the sound which that rune designated, \emph{*z \char`~\ *ʀ}, was lost in both Old English and Old West Norse long before the rune poems were written.

\bookStart{The Old English Rune Poem}
\setBookCode{EnglishRunePoem}

\begin{flushright}%
\textbf{Dating:} C8th--C10th%TODO

\textbf{Meter:} Germanic alliterative meter
\end{flushright}%

\section{Introduction}

The \textbf{Old English Rune Poem} (\EnglishRunePoem) is the oldest of the three old rune poems.  It describes the 29 runes of the Anglo-Saxon Futhorc.

\subsection{Textual preservation and transmission}

\EnglishRunePoem\ only survives in the printed edition of \textcite[\nopp vol. 1, pp. 135]{Hickes1703}, based on the original manuscript which was certainly destroyed in the Cottonian fire of 1731.

Hickes introduces the poem with the following words (original emphasis kept): \emph{E Codice MS. Bibliothecæ} Cottonianæ, \emph{cujus nota, Otho, B. 10.} ‘From a manuscript codex of the Cotton Library, whose shelfmark is Otho B. 10.’

TODO: Note on the visual layout.  Cite Males article on Rune poems.

\subsection{Meter and stanza}

The poem is divided into stanzas, one for each rune.  Unlike the two Scandinavian rune-poems, the stanzaic structure of \EnglishRunePoem\ is somewhat uneven, and the stanzas vary in length between 2 and 5 lines.  The most common stanza-length is 3 lines (19/29: 1--8, 11--14, 16--17, 23--24, 26--28), followed by 4 lines (7/29: 15, 18--22, 25); rarer lengths are 2 (9--10) and 5 lines (29).

\sectionline

\section{Text}

\bvg\bva%[ᚠ]%
(feoh) byþ \alst{f}rófur \hld\ \alst{f}ira ge·hwylcum. &
Sceal ðeah \alst{m}anna ge·hwylc \hld\ \alst{m}iclun hyt dǽlan &
gif he wile for \alst{d}rihtne \hld\ \alst{d}ómes hleotan.\eva

\bvb TODO: TRANSLATION.\evb\evg


\bvg\bva%[ᚢ]%
(ur) byþ \alst{â}n-mód \hld\ and \alst{o}fer-hyrned, &
\alst{f}ela-\alst{f}récne deor, \hld\ \alst{f}eohteþ mid hornum, &
\alst{m}ǽre \alst{m}ór-stapa; \hld\ þæt is \alst{m}ódig wuht.\eva

\bvb TODO: TRANSLATION.\evb\evg


\bvg\bva%[ᚦ]%
(ðorn) byþ \alst{ð}earle scearp; \hld\ \alst{ð}egna ge·hwylcum &
\alst{a}n-feng ys \alst{y}fyl, \hld\ \alst{u}n-ge·metun reþe &
\alst{m}anna ge·hwylcun \hld\ ðe him \alst{m}id resteð.\eva

\bvb TODO: TRANSLATION.\evb\evg


\bvg\bva%[ᚩ]%
(os) byþ \alst{o}rd-fruma \hld\ \alst{æ}lcre spræce, &
\alst{w}ís-dómes \alst{w}raþu \hld\ and \alst{w}itena frofur, &
and \alst{eo}rla ge·hwam \hld\ \alst{ea}d-nys and to·hiht.\eva

\bvb TODO: TRANSLATION.\evb\evg


\bvg\bva%[ᚱ]%
(rad) byþ on \alst{r}ecyde \hld\ \alst{r}inca ge·hwylcum &
\alst{s}efte, and \alst{s}wiþ-hwæt \hld\ ðam ðe \alst{s}itteþ on ufan &
\alst{m}eare \alst{m}ægen-heardum \hld\ ofer \alst{m}íl-paþas.\eva

\bvb TODO: TRANSLATION.\evb\evg


\bvg\bva%[ᚳ]%
(cen) byþ \alst{c}wicera ge·hwam \hld\ \alst{c}u̇þ on fyre, &
\alst{b}lac and \alst{b}eorht-líc, \hld\ \alst{b}yrneþ oftust &
ðær hí \alst{æ}þelingas \hld\ \alst{i}nne restaþ.\eva

\bvb TODO: TRANSLATION.\evb\evg


\bvg\bva%[ᚷ]%
(gyfu) \alst{g}umena byþ \hld\ \alst{g}leng and herenys, &
\alst{w}raþu and \alst{w}yrþ-scype, \hld\ and \alst{w}ræcna ge·hwam &
\alst{a}r and \alst{æ}twist \hld\ ðe byþ \alst{o}þra leas.\eva

\bvb TODO: TRANSLATION.\evb\evg


\bvg\bva%[ᚹ]%
(wen) ne bruceþ \hld\ ðe can \alst{w}éana lýt, &
\alst{s}âres and \alst{s}orge, \hld\ and him \alst{s}ylfa hæfþ &
\alst{b}lǽd and \alst{b}lysse \hld\ and eac \alst{b}yrga ge·niht.\eva

\bvb TODO: TRANSLATION.\evb\evg


\bvg\bva%[ᚻ]%
(hægl) byþ \alst{h}wítust corna; \hld\ hwyrft hit of \alst{h}eofones lyfte, &
\alst{w}ealcaþ hit \alst{w}indes scura, \hld\ weorþeþ hit to \alst{w}ætere syððan.\eva

\bvb TODO: TRANSLATION.\evb\evg


\bvg\bva%[ᚾ]%
(nyd) byþ \alst{n}earu on breostan, \hld\ weorþeþ hi ðeah oft \alst{n}iþa bearnum &
to \alst{h}elpe and to \alst{h}æle ge·hwæþre, \hld\ gif hí his \alst{h}lystaþ æror.\eva

\bvb TODO: TRANSLATION.\evb\evg


\bvg\bva%[ᛁ]%
(is) byþ \alst{o}fer-ceald, \hld\ \alst{u}n-ge·metum slidor, &
\alst{g}lisnaþ \alst{g}læs-hluttur, \hld\ \alst{g}immum ge·licust, &
\alst{f}lor \alst{f}orste ge·woruht, \hld\ \alst{f}æger an-sýne.\eva

\bvb TODO: TRANSLATION.\evb\evg


\bvg\bva%[ᛄ]%
(ger) byþ \alst{g}umena hiht, \hld\ ðon \alst{G}od læteþ, &
\alst{h}âlig \alst{h}eofones cyning, \hld\ \alst{h}rusan syllan &
\alst{b}eorhte \alst{b}leda \hld\ \alst{b}eornum and ðearfum.\eva

\bvb TODO: TRANSLATION.\evb\evg


\bvg\bva%[ᛇ]%
(eoh) byþ \alst{u}tan \hld\ \alst{u}n-smeþe treow, &
\alst{h}eard, \alst{h}rusan fæst, \hld\ \alst{h}yrde fyres, &
\alst{w}yrt-rumun under·\alst{w}reþyd, \hld\ \alst{w}ynan on éþle.\eva

\bvb TODO: TRANSLATION.\evb\evg


\bvg\bva%[ᛈ]%
(peorð) byþ symble \hld\ \alst{p}lega and hlehter &
{[...]} \alst{w}lancum \hld\ ðar \alst{w}igan sittaþ &
on \alst{b}eor-sele \hld\ \alst{b}líþe æt·somne. \eva

\bvb TODO: TRANSLATION.\evb\evg


\bvg\bva%[ᛉ]%
(eolhx)-secg \alst{ea}rd hæfþ \hld\ \alst{o}ftust on fenne, &
\alst{w}exeð on \alst{w}ature, \hld\ \alst{w}undaþ grimme, &
\alst{b}lode \alst{b}reneð \hld\ \alst{b}eorna ge·hwylcne &
ðe him \alst{æ}nigne \hld\ \alst{o}n-feng ge·deð.\eva

\bvb TODO: TRANSLATION.\evb\evg


\bvg\bva%[ᛋ]%
(sigel) \alst{s}é-mannum \hld\ \alst{s}ymble biþ on hihte, &
ðonn hi hine \alst{f}eriaþ \hld\ ofer \alst{f}isces beþ, &
oþ hí \alst{b}rim-hengest \hld\ \alst{b}ringeþ to lande.\eva

\bvb TODO: TRANSLATION.\evb\evg


\bvg\bva%[ᛏ]%
(tir) biþ \alst{t}âcna sum, \hld\ healdeð \alst{t}rywa wel &
wiþ \alst{æ}þelingas, \hld\ \alst{â} biþ on færylde, &
ofer \alst{n}ihta ge·\alst{n}ipu \hld\ \alst{n}æfre swiceþ.\eva

\bvb TODO: TRANSLATION.\evb\evg


\bvg\bva%[ᛒ]%
(beorc) byþ \alst{b}leda leas, \hld\ \alst{b}ereþ efne swa ðeah &
\alst{t}ânas b·útan \alst{t}udder, \hld\ biþ on \alst{t}elgum wlitig, &
\alst{h}eah on \alst{h}elme \hld\ \alst{h}rysted fægere, &
ge·\alst{l}oden \alst{l}eafum, \hld\ \alst{l}yfte ge·tenge.\eva

\bvb TODO: TRANSLATION.\evb\evg


\bvg\bva%[ᛖ]%
(eh) byþ for \alst{eo}rlum \hld\ \alst{æ}þelinga wyn, &
\alst{h}ors \alst{h}ófum wlanc, \hld\ ðær him \alst{h}æleþe ymb, &
\alst{w}elege on \alst{w}icgum, \hld\ \alst{w}rixlaþ spræce, &
and biþ \alst{u}n-styllum \hld\ \alst{æ}fre frofur.\eva

\bvb TODO: TRANSLATION.\evb\evg


\bvg\bva%[ᛗ]%
(man) byþ on \alst{m}yrgþe \hld\ his \alst{m}agan leof; &
sceal þeah \alst{â}nra ge·hwylc \hld\ \alst{o}ðrum swícan, &
for ðam \alst{d}ryhten wyle \hld\ \alst{d}óme síne &
þæt \alst{ea}rme flæsc \hld\ \alst{eo}rþan be·tæcan.\eva

\bvb TODO: TRANSLATION.\evb\evg


\bvg\bva%[ᛚ]%
(lagu) byþ \alst{l}eodum \hld\ \alst{l}ang-sum ge·þuht, &
gif hí sculun \alst{n}eþun \hld\ on \alst{n}acan tealtum, &
and hi \alst{s}æyþa \hld\ \alst{s}wýþe bregaþ, &
and se \alst{b}rim-hengest \hld\ \alst{b}ridles ne gymeð.\eva

\bvb TODO: TRANSLATION.\evb\evg


\bvg\bva%[ᛝ]%
(ing) wæs \alst{ǽ}rest \hld\ mid Éast-Dęnum &
ge·\alst{s}ewen \alst{s}ęcgun, \hld\ oþ he \alst{s}iððan éast &
ofer \alst{w}ǽg ge·\alst{w}ât, \hld\ \alst{w}æn æfter rann; &
ðus \alst{h}eardingas \hld\ ðone \alst{h}æle nęmdun.\eva

\bvb TODO: TRANSLATION.\evb\evg


\bvg\bva%[ᛟ]%
(eþel) byþ \alst{o}fer-leof \hld\ \alst{æ}g·hwylcum men, &
gif he mot ðǽr \alst{r}ihtes \hld\ and ge·\alst{r}ysena on &
\alst{b}rúcan on \alst{b}lode \hld\ \alst{b}leadum oftast.\eva

\bvb TODO: TRANSLATION.\evb\evg


\bvg\bva%[ᛞ]%
(dæg) byþ \alst{d}rihtnes sond, \hld\ \alst{d}eore mannum, &
\alst{m}ære \alst{m}etodes leoht, \hld\ \alst{m}yrgþ and to·hiht &
\alst{ea}dgum and \alst{ea}rmum, \hld\ \alst{ea}llum brice.\eva

\bvb TODO: TRANSLATION.\evb\evg


\bvg\bva%[ᚪ]%
(ac) byþ on \alst{eo}rþan \hld\ \alst{ę}lda bearnum &
\alst{f}læsces \alst{f}ódor, \hld\ \alst{f}ereþ ge·lome &
ofer \alst{g}anotes bæþ; \hld\ \alst{g}âr-sęcg fandaþ &
hwæþer \alst{â}c hæbbe \hld\ \alst{æ}þele treowe.\eva

\bvb TODO: TRANSLATION.\evb\evg


\bvg\bva%[ᚫ]%
(æsc) biþ \alst{o}fer-heah, \hld\ \alst{ę}ldum dýre, &
\alst{st}iþ on \alst{st}aþule, \hld\ \alst{st}ęde rihte hylt, &
ðeah him \alst{f}eohtan on \hld\ \alst{f}ir\emph{h}as mǫnige.\eva

\bvb TODO: TRANSLATION.\evb\evg


\bvg\bva%[ᚣ]%
(yr) byþ \alst{æ}þelinga \hld\ and \alst{eo}rla ge·hwæs &
\alst{w}yn and \alst{w}yrþ-mynd, \hld\ byþ on \alst{w}icge fæger, &
\alst{f}æst-lic on \alst{f}ær-elde, \hld\ \alst{f}yrd-geatewa sum.\eva

\bvb TODO: TRANSLATION.\evb\evg


\bvg\bva%[ᛡ]%
(iar, ior) byþ \alst{éa}-fixa, \hld\ and ðeah \alst{á} bruceþ &
\alst{f}ódres on \alst{f}oldan, \hld\ hafaþ \alst{f}ægerne eard, &
\alst{w}ætre be·\alst{w}orpen, \hld\ ðær he \alst{w}ynnum leofaþ.\eva

\bvb TODO: TRANSLATION.\evb\evg


\bvg\bva%[ᛠ]%
(ear) byþ \alst{e}gle \hld\ \alst{e}orla ge·hwylcun, &
ðonn \alst{f}æst-lice \hld\ \alst{f}læsc on·ginneþ, &
\alst{h}raw colian, \hld\ \alst{h}rusan ceosan &
\alst{b}lac to ge·\alst{b}eddan; \hld\ \alst{b}leda ge·dreosaþ, &
\alst{w}ynna ge·\alst{w}itaþ, \hld\ \alst{w}era ge·swicaþ.\eva

\bvb TODO: TRANSLATION.\evb\evg

\sectionline

\bookStart{The Icelandic Rune Poem}
\setBookCode{IcelandicRunePoem}

\begin{flushright}%
\textbf{Dating:} Mediaeval.%TODO

\textbf{Meter:} Unclear.
\end{flushright}%

\section{Introduction}

The \textbf{Icelandic Rune Poem} (\IcelandicRunePoem) is one of two West Norse rune poems describing the Younger Futhark.

\subsection{Textual preservation and transmission}

\IcelandicRunePoem\ is preserved in several manuscripts, none mediaeval.  The oldest is AM 687 d 4to (henceforth 687d), dating to ca.~1500 \ce, where the poem takes up the leaf 1v.  In this ms.~each rune takes up one line.  The names of the runes are not given, but the name is written out by the rune, so that the first stanza begins \emph{ᚠ er frę́nda róg}, the second \emph{ᚢ er skýgia grátur}, and so forth.  687d/1v is not physically fragmented, but much of the ink has faded and large parts of thep oem are therefore illegible.

TODO.  It is only attested in late manuscripts which often have major disagreements with each other.

\subsection{Meter and stanza}

The meter of \IcelandicRunePoem\ is a divergent form of \Ljodahattr.  In particular, the a-verse always contains three lifts (certainly owing to the obligatory rune-name), and the c-verse may end in a dip following a heavy lift.
The poem is strictly structured.  Each rune is given two lines, the first consisting of two half-verses (the a- and b-verse), the second being a form of c-verse.  These two lines will henceforth be referred to as stanzas.

\subsection{Style and contents}

Like in the two other rune poems, the 16 runes in \IcelandicRunePoem\ are described by reference to the object to which their name refers (\emph{fé} ‘wealth’ for ᚠ, etc.).  Like its Scandinavian sister-poem \textlink{NorwegianRunePoem}, \IcelandicRunePoem\ uses kennings to describe these objects.  Numerous such kennings are shared, sometimes word-for-word, between the two.
Like in the two other rune poems, the very first word in each stanza is the respective rune’s name, which is a noun.  The name is then followed by exactly three kennings, one in each verse.


TODO: reference Mikael Males article about rune poems


\sectionline

\section{Text}

\bvg\bva%
\alst{F}é es \edtrans{\alst{f}rę́nda róg}{strife of kinsmen}{\Bfootnote{Self-explanatory; wealth causes conflict between heirs.  The same expression is found in \textlink{NorwegianRunePoem}.}} \hld\ ok \edtrans{\alst{f}lǿðar viti}{beacon of the flood}{\Bfootnote{A formulaic type of \textsc{gold}-kenning (TODO: \parencite{Meissner1921}); since gold does not rust, it is the fire which shines when submerged in water.}} &
\ind ok \edtrans{\alst{g}raf-sęiðs \alst{g}ata}{grave-saithe’s \ken{serpent’s} street}{\Bfootnote{Another formulaic kenning.  The serpent or Wyrm that lies upon its hill of treasures is a prominent image in Old Germanic legend; cf.~particularly Fathomer (\textlink{Reginsmal}, \textlink{Fafnismal}, and \VolsungaSaga) and the unnamed dragon in \textlink{Beowulf}.}}.\eva

\bvb Wealth is strife of kinsmen and beacon of the flood \\
\ind and grave-saithe’s \ken{serpent’s} street.\evb\evg


\bvg\bva%
Úr es \alst{sk}ýja grátr \hld\ ok \alst{sk}ára þvęrrir &
\ind ok \alst{h}irðis \alst{h}atr.\eva

\bvb Drizzle is weeping of clouds and ... \\
\ind and shepherd’s hatred.\evb\evg


\bvg\bva%
Þurs es \alst{k}venna \alst{k}vǫl \hld\ ok \alst{k}letta í·búi &
\ind ok \alst{v}arð-rúnar \alst{v}err.\eva

\bvb Thurse is women’s torment and indweller of hills \\
\ind and husband of the cairn-whisperess \ken{trollwoman}.\evb\evg


\bvg\bva%
\alst{Ǫ̇}ss es \alst{a}ldinn gautr \hld\ ok \alst{Ǫ̇}s·garðs jǫfurr, &
\ind ok \alst{V}al-hallar \alst{v}ísi.\eva

\bvb Os is ancient Geat, and Osyard’s chief, \\
\ind and Walhall’s overseer.\evb\evg


\bvg\bva%
Ręið es \alst{s}itjandi \alst{s}ę́la \hld\ ok \alst{s}núðig fęrð &
\ind ok \alst{jó}s \alst{ę}rfiði.\eva

\bvb Chariot is sitting bliss and twirling journey \\
\ind and horse’s heavy work.\evb\evg


\bvg\bva%
Kaun es \alst{b}arna \alst{b}ǫl \hld\ ok \alst{b}ar-dagi &
\ind ok \alst{h}old-fúa \alst{h}ús.\eva

\bvb Boil is children’s curse and TODO \\
\ind and house of flesh-rot.\evb\evg


\bvg\bva%
Hagall es \alst{k}alda \alst{k}orn \hld\ ok \alst{k}nappa drífa &
\ind ok \alst{s}náka \alst{s}ótt.\eva

\bvb Hail is cold kernel and storm of beads \\
\ind and sickness of snakes.\evb\evg


\bvg\bva%
Nauð es \alst{þ}ýjar \alst{þ}rǫ́ \hld\ ok \alst{þ}ungr kostr &
\ind ok \alst{v}ás-samlig \alst{v}erk.\eva

\bvb Need is maidservant’s yearning and scant choice \\
\ind and working in wet-cold weather.\evb\evg


\bvg\bva%
\alst{Í}ss es \alst{á}ar bǫrkr \hld\ ok \alst{u}nnar þękja &
\ind ok \alst{f}ęigra manna \alst{f}ár.\eva

\bvb Ice is river’s bark and wave’s roof \\
\ind and fey men’s danger.\evb\evg


\bvg\bva%
Ár es \alst{g}umna \alst{g}óði \hld\ ok \alst{g}ótt sumar &
\ind \emph{ok} \alst{a}l-gróinn \alst{a}kr.\eva

\bvb Year is men’s boon and good summer \\
\ind (and) all-grown acre.\evb\evg


\bvg\bva%
Sól es \alst{sk}ýja \alst{sk}jǫldr \hld\ ok \alst{sk}ínandi rǫðull &
\ind ok \alst{í}sa \alst{a}ldr-tregi.\eva

\bvb Sun is the shield of clouds and shining wheel \\
\ind and ice-sheets’ life-sorrow.\evb\evg


\bvg\bva%
Týr es \alst{ęi}n-hęndr \alst{ǫ̇}ss \hld\ ok \alst{u}lfs lęifar &
\ind ok \alst{h}ofa \alst{h}ilmir.\eva

\bvb Tew is the one-handed Os and the wolf’s leftovers \\
\ind and lord of hoves.\evb\evg


\bvg\bva%
Bjarkan es \alst{l}aufgat \alst{l}im \hld\ ok \alst{l}ítit tré &
\ind ok \alst{u}ng-samligr \alst{v}iðr.\eva

\bvb Birch is leafy branch and little tree \\
\ind and youthful wood.\evb\evg


\bvg\bva%
\alst{M}aðr es \alst{m}anns gaman \hld\ ok \alst{m}oldar auki &
\ind ok \alst{sk}ipa \alst{sk}ręytir.\eva

\bvb Man is man’s joy and the product of dust \\
\ind and adorner of ships.\evb\evg


\bvg\bva%
Lǫgr es \alst{v}ellanda \alst{v}atn \hld\ ok \alst{v}íðr kętill &
\ind ok \alst{g}lǫmmungr \alst{g}rund.\eva

\bvb Liquid is boiling water and wide kettle \\
\ind and TODO.\evb\evg


\bvg\bva%
Ýr es \alst{b}ęndr bogi \hld\ ok \alst{b}rot-gjarnt járn &
\ind ok \alst{f}ęnju \alst{f}lęygir.\eva

\bvb Yew is a bent bow and easily broken iron \\
\ind and arrow’s hurler.\evb\evg

\sectionline

\bookStart{The Norwegian Rune Poem}
\setBookCode{NorwegianRunePoem}

\begin{flushright}%
\textbf{Dating:} Mediaeval.%TODO

\textbf{Meter:} Unclear.
\end{flushright}%

\section{Introduction}

The \textbf{Norwegian rune poem} (\NorwegianRunePoem) is one of two West Norse rune poems describing the 16 runes of the Younger Futhark.

\subsection{Textual preservation and transmission}

TODO.

\subsection{Style and contents}

\NorwegianRunePoem\ is closely related to \IcelandicRunePoem, which also describes the Younger Futhark.  With the exception of runes 2 (ᚢ, here \emph{úr} ‘spark’) and 4 (ᚬ, here \emph{óss} ‘river-mouth’), the names of the runes are identical between them, as are many of the poetic expressions.

The description of all runes consists of two clauses or sentences, generally one in each verse.  The first clause gives the name of the rune and a short description of it, while the second clause is more tangential and oftentimes seems to have nothing to do with the name.  A clever solution to this problem is that the second clause connects to the visual shape of the rune (TODO: cite).  For some of the runes, this turns out to be a rewarding line of inquiry, and potential links between the shape of a rune and the second clause will be explored in the notes.

\subsection{Meter}

Each rune is described by a couplet of 2 verses linked by both end-rhyme and alliteration.  The meter is noticeably trochaic, and the end-rhyme of each verse always falls on a trochée.  In 15 out of 16 lines, the a-verse has 2 alliterations.  The b-verse always has exactly 1 alliteration, falling on its very first word.

\subsection{Language}

The language is unmistakably that of mediaeval Norway.  As can be seen from the rhymes and alliteration the following uniquely Norwegian sound changes have occurred:
\begin{itemize}
  \item \emph{hl, hn, hr} > \emph{l, n, r} (2 \emph{lęypr} < \emph{hlęypr}; 8 \emph{nęppa} < \emph{hnęppa}; 5 \emph{rossum} < \emph{hrossum}).
  \item \emph{rst} > \emph{st} (5 \emph{vęsta} < \emph{vęrsta})
\end{itemize}

\sectionline

\newpage

\section{Text}

\bvg\bva%[ᚠ]%
\alst{F}é vęldr \alst{f}rę́nda rógi; \hld\ \edtrans{\alst{f}ǿðisk ulfr í skógi}{the wolf rears itself in the wood}{\Bfootnote{Almost certainly visually alluding to the similarity between the rune and the upwards curving branches of the pine tree; the keyword is \emph{skógr} ‘wood, forest’.}}.\eva

\bvb ‘Wealth’ causes the strife of kinsmen; the wolf rears itself in the wood.\evb\evg


\bvg\bva%[ᚢ]%
\edtrans{\alst{Ú}r}{Spark}{\Bfootnote{Properly ‘drizzle’, but here used metaphorically for the sparks or tiny particles of hot metal that flow out of burning iron when struck.}} er af \alst{ę}ldu \alst{ja}rni; \hld\ \edtrans{\alst{o}pt lęypr ręinn á hjarni}{oft leaps the reindeer on the snow}{\Bfootnote{Likely visually alluding to the similarity between the rune and the hoof-print of a reindeer in the snow; the keyword is \emph{ręinn} ‘reindeer’.}}.\eva

\bvb ‘Spark’ comes from glowing iron; oft leaps the reindeer on the snow.\evb\evg


\bvg\bva%[ᚦ]%
Þurs vęldr \alst{k}vinna \alst{k}villu; \hld\ \edtrans{\alst{k}átr verðr fár af illu}{few are gladdened by evil}{\Bfootnote{Lit. ‘few a man becomes glad from evil’.  Visually alluding to the similarity between the rune and a smiling face; the keyword is \emph{kátr} ‘glad, happy’.}}.\eva

\bvb ‘Thurse’ causes women’s ailment; few are gladdened by evil.\evb\evg


\bvg\bva%[ᚬ]%
Óss er \alst{f}lęstra \alst{f}ęrða \hld\ \alst{f}ǫr, \edtrans{en skalpr er sverða}{but the scabbard-mouth of swords}{\Bfootnote{Perhaps visually alluding to the similarity between the rune and a sword in its scabbard, although that is uncertain.}}.\eva

\bvb ‘River-mouth’ is the route of most journeys, but the scabbard-mouth of swords.\evb\evg


\bvg\bva%[ᚱ]%
\alst{R}ęið kveða \alst{r}ossum vęsta; \hld\ \alst{R}ęginn sló sverð’it bęsta.\eva

\bvb ‘Chariot’ they say is worst for the horses; Rein forged the best sword.\evb\evg


\bvg\bva%[ᚴ]%
Kaun er \alst{b}arna \alst{b}ǫlvan; \hld\ \alst{b}ǫl gørvir ná fǫlvan.\eva

\bvb ‘Ulcer’ is the curse of children; a bale makes a corpse pale.\evb\evg


\bvg\bva%[ᚼ]%
Hagall er \alst{k}aldastr \alst{k}orna; \hld\ \edtrans{\alst{K}ristr skóp hęim’inn forna}{Christ created the world of yore}{\Bfootnote{Visually alluding to the similarity between the rune and the chi-rho; the keyword is \emph{Kristr} ‘Christ’.}}.\eva

\bvb ‘Hail’ is coldest of kernels; Christ created the world of yore.\evb\evg


\bvg\bva%[ᚾ]%
\alst{N}auðr gørir \alst{n}ęppa kosti; \hld\ \edtrans{\alst{n}øktan kęlr í frosti}{the naked man freezes in the cold}{\Bfootnote{Perhaps visually alluding to the similarity between the rune and the human body in the form of a naked stick-figure; the keyword is \emph{nøktr} ‘naked’.}}.\eva

\bvb ‘Need’ gives scant choice; the naked man freezes in the cold.\evb\evg


\bvg\bva%[ᛁ]%
Ís kǫllum \alst{b}rú \alst{b}ręiða; \hld\ \edtrans{\alst{b}lindan þarf at lęiða}{the blind man must be lead}{\Bfootnote{Perhaps visually alluding to the similarity between the rune and a blind man’s cane or walking stick; the keyword is \emph{blindr} ‘blind’.}}.\eva

\bvb ‘Ice’ we call a broad bridge; the blind man must be lead.\evb\evg


\bvg\bva%[ᛅ]%
Ár er \alst{g}umna \alst{g}óði; \hld\ \edtrans{\alst{g}et’k at ǫrr var Fróði}{I recall that Frood was generous}{\Bfootnote{Frood, or as he is often known, Frith-Frood (ON \emph{frið-Fróði}), was a legendarily generous Danish king during whose reign good harvests and peace (\emph{ár ok friðr} ‘year and frith’) abounded.  He is thus mentioned due to his connection with the rune’s name, not due to any visual elements of the rune.}}.\eva

\bvb ‘Year’ is men’s boon; I recall that Frood was generous.\evb\evg


\bvg\bva%[ᛋ]%
Sól er \alst{l}anda \alst{l}jómi; \hld\ \edtrans{\alst{l}úti’k hęlgum dómi}{I bow to the holy shrine}{\Bfootnote{Visually alluding to the similarity between the rune and the shape of a kneeling man; the keyword is \emph{lúti’k} ‘I bow’.}}.\eva

\bvb ‘Sun’ is the light of the lands; I bow to the holy shrine.\evb\evg


\bvg\bva%[ᛏ]%
Týr er \alst{ę}in-ęndr \alst{á}sa; \hld\ \edtrans{\alst{o}pt verðr smiðr blása}{often the smith must blow}{\Bfootnote{Perhaps visually alluding to the similarity between the shape of the rune and a hammer; the keyword is \emph{smiðr} ‘smith’.}}.\eva

\bvb ‘Tew’ is the one-handed of the Eese; often the smith must blow.\evb\evg


\bvg\bva%[ᛒ]%
Bjarkan er \alst{l}auf-grǿnstr \alst{l}ima; \hld\ \alst{L}oki bar flę́rða tíma.\eva

\bvb ‘Birch’ is the leaf-greenest of branches; Lock was successful in cunning.\evb\evg


\bvg\bva%[ᛘ]%
\alst{M}aðr er \alst{m}oldar auki; \hld\ \edtrans{\alst{m}ikil er gręip á hauki}{mighty is the grip on the hawk}{\Bfootnote{Visually alluding to the shape of the bird’s talon; the keyword is \emph{haukr} ‘hawk’.}}.\eva

\bvb ‘Man’ is the product of dust; mighty is the grip on the hawk.\evb\evg


\bvg\bva%[ᛚ]%
Lǫgr er þar er \alst{f}ęllr ór \alst{f}jalli \hld\ \alst{f}oss; en gull eru nossir.\eva

\bvb ‘Water’ is when from the fell falls a torrent, but jewels are gold.\evb\evg


\bvg\bva%[ᛦ]%
Ýr er \alst{v}etr-grǿnstr \alst{v}iða; \hld\ \alst{v}ę́nt er, er brennr, at sviða. \eva

\bvb ‘Yew’ is winter-greenest of trees; ’tis expected, when it burns, to get singed.\evb\evg

\sectionline



% \bookStart{Runic Poetry in the Elder Futhark}

% \input{books/runic/Elder Futhark.tex}

\bookStart{Runic Poetry from Sweden and Gotland}

%Runic poetry from Sweden and Gotland.

\section{Introduction}

The inscriptions are sorted according to their runological signa as found in the Scandinavian Runic-text Data Base (SRDB).  For most inscriptions the signum derives from the standard corpus edition \emph{Sveriges runinskrifter} (‘Runic inscriptions of Sweden’) and consists of two elements, the abbreviated traditional province and the number given in the corpus edition.  For example Sm 39 is the inscription given number 39 in the Småland volume of \emph{Sveriges runinskrifter}.

Inscriptions discovered after the publication of the corpus edition describing the relevant province have instead of the number an abbreviated bibliographical reference to the \emph{editio princeps}.  For example Sö Fv1948;289 is a inscription from Södermanland first published in the 1948 volume of \emph{Fornvännen}, page 289.

The province of Scania, although part of Sweden since 1658, was in the Wiking Age and the mediaeval period an integral part of Denmark and as such its metrical inscriptions are found under Runic poetry from Denmark below.

\subsection{Abbreviations of Swedish provinces}

\begin{itemize}
  \item G = Gotland
  \item Nä = Närke
  \item Sm = Småland
  \item Sö = Södermanland
  \item U = Uppland
  \item Vg = Västergötland
  \item Vs = Västmanland
  \item Ög = Östergötland
  \item Öl = Öland
\end{itemize}

\sectionline
\newpage


\section{G 113 (Ardre)}

\begin{flushright}%
\textbf{Dating:} C12th

\textbf{Location:} Ardre church

\textbf{Meter:} \Fornyrdislag
\end{flushright}%

\subsection{Introduction}

A picture stone.  The palatal \emph{ʀ} has in several places been restored for the stone’s \textbf{r}.

\subsection{Runes}

TODO: Runes%textrunic

\subsection{Transliteration}

§ A \textbf{÷ utar + ak + kaiʀuatr + ak + aiuatr + þaʀ + setu + stain + ebtir + likna(t) + faþur ÷ sen +} \\
§ B \textbf{÷ raþialbr + ak + kaiʀaiau(t)- + þaiʀ kiarþu + merki + kuþ + ubtir + man + saaran ÷ likraibr + risti + runaʀ ÷}

\subsection{Text}

\bpg\bpa[0]%
\startsub \flagstanza{\textbf{A}}%
Ótt·arr ok Gaiʀ·hvatr ok Ai·hvatr þa\emph{i}ʀ sęttu stain ęptiʀ Líkn·\emph{hv}at, faður si̇nn.\endsub\epa

\bpb Oughter and Garwhat and Eanwhat, they set the stone after Likenwhat, their father.\epb\epg


\bvg\bva[B][2]%
Ráð·þialfʀ ok Gaiʀ·\emph{n}iaut\emph{r}, &
þaiʀ giǫrðu \alst{m}ęrki góð \hld\ ęptiʀ \alst{m}ann snaran. &
Líkn·raifʀ rísti ru̇naʀ.\eva

\bvb Redthelf and Garneat, \\
they made good landmarks after a swift man. \\
Likenraf carved the runes.\evb\evg

\sectionline\clearpage


\section{G 188}

\begin{flushright}%
\textbf{Dating:} C11th

\textbf{Meter:} \Fornyrdislag
\end{flushright}%

\subsection{Introduction}

A damaged and unusually difficult stone to interpret.  Most of the inscription survives only in older depictions which are likely not entirely faithful.

\subsection{Runes}

ᚱᛆᚦᚽᛏ ᛭ ᛘᛁᚦᚱᚽᛏᚢ \\
ᛒᛁᚾ ᛭ ᛋᛁᛅᛚᚢᚼ

TODO: Runes%textrunic

\subsection{Transliteration}

\textbf{[... auk × botaiþr × þaun litu × raisa s... ... haito × mirki mirlakt eftir man koþ-n · kiarn uaʀ × hann arla × a × kuþ troa × rumshafþ ×] raþ(e)t × miþ r(e)tu m[iug × tyrlak × ta · ta × uam · la · reksþe... ...](b)in × sialu h-... [lit × kiarþi stain × uataruek...]}

\subsection{Text}

\bpg\bpa[0]%
\startsub \dots\ auk Bót·aiðr, \edtrans{þaun}{they}{\Bfootnote{The neuter pl. form indicates that the raisers are a mixed-sex group.}} létu raisa s\emph{tain} \textbf{†haito†}\endsub \\
\alst{m}ęrki \alst{m}ér-lekt \hld\ ęptiʀ \alst{m}ann góðan; \\
\alst{g}iarn vaʀ hann árla \hld\ ȧ \alst{G}uð tróa. \\
\textbf{†rumshafþ†} ráðit með réttu \\
miuk dýr-lek \textbf{†ta:ta:uam:la:reksþe†} \\
\startsub ... \emph{hial}pin siálu h\emph{ans}. \\
\textbf{†lit†} giarði stain \textbf{†uataruek†}\endsub\epa

\bpb \dots\ and Bootheath they had the stain raised \dots \\
A glorious landmark after a good man. \\
He was early eager to believe in God; \\
\dots\ rightly interpreted, \\
very glorious \dots \\
\dots\ help his soul. \\
\dots\ made the stone \dots\epb\epg

\sectionline\clearpage


\section{G 203 (Hogräns)}

\begin{flushright}%
\textbf{Dating:} C12th

\textbf{Location:} Hogräns church

\textbf{Meter:} \Fornyrdislag
\end{flushright}%

\subsection{Introduction}

TODO.

\subsection{Runes}

TODO: Runes%textrunic

\subsection{Transliteration}

\textbf{sigmutr let rasa sain eftiʀ bruþr : sina : auk : bro : kierua : eftiʀ : sikbiern : santa mikal hie[lbi] ... ...ans auk : at : botraif auk at sigraif : auk : at aibiern : faþur þaiʀa : altr : auk bikui han : i by : sunarst kaiʀuiþr lekþi ormaluʀ nemʀ : in[t]i uʀ
¶ sikmutr [{}-{}-fiʀ :] sliku : unit kuml
¶ karmanum : þet aʀ [:] ... kun : hier : mun : stanta stain : a[t] : merki bietr a : bierki in bro furiʀ
¶ roþ(b)iern risti run(i)ʀ [þ]esa kaiʀl-ifʀ sumaʀ aʀ karla kan}

\subsection{Text}

\bpg\bpa[0]%
\startsub Sig·mundr lét raisa stain eptiʀ brýðr si̇́na, auk bró gierva eptiʀ Sig·biern---Sankta Mikál hielpi \emph{siál h}ans---auk at Bót·raif auk at Sig·raif auk at Ai·biern, faður þaiʀa aldra,\endsub\epa

\bpb Syemund had this stone raised after his brothers and the bridge made after Syebarn---may Saint Michael help his soul---and after Bootraf and after Syeraf and after Eanbarn, the father of them all,\epb\epg

\bvg\bva[][2]%
auk \alst{b}yggvi hann \hld\ i̇́ \alst{b}ý sunnarst.\eva

\bvb and he lived on the southernmost farm.\evb\evg

\bpg\bpa[0]%
Gaiʀviðr lęgði orm-áluʀ; némʀ innti ýʀ.\epa

\bpb Garwith laid the serpent-tracks; TODO.\epb\epg

\bvg\bva[][4]%
\alst{S}ig·mundr \emph{hefiʀ} \hld\ \alst{s}·líku unnit &
\alst{k}uml \alst{k}arl-mannum. \hld\ Þet aʀ †\alst{k}e...† kunn. &
Hier mun \alst{st}anda \hld\ \alst{st}ainn at merki, &
\alst{b}ie\emph{r}tr á \alst{b}iergi, \hld\ en \alst{b}ró fyriʀ; &
Róð·biern \alst{r}ísti \hld\ \alst{r}úniʀ \emph{þ}essaʀ, &
\alst{G}aiʀ·l\emph{a}ifʀ sumaʀ, \hld\ aʀ \alst{g}arla kann.\eva

\bvb Syemund has suchwise accomplished \\
a monument for men; that is known to ... \\
Here the stone will stand as a mark, \\
bright on the hill and the bridge below. \\
Rothbarn carved these runes, \\
{[and]} Garlaf, who knows clearly, some.\evb\evg

\sectionline\clearpage


\section{G 207}

\begin{flushright}%
\textbf{Dating:} C12th

\textbf{Meter:} \Fornyrdislag
\end{flushright}%

\subsection{Introduction}

A damaged stone. TODO.  The same people also commissioned the likewise damaged G 208 below.

\subsection{Runes}

TODO: Runes%textrunic

\subsection{Transliteration}

\textbf{butmuntr : auk : butraifʀ : auk : kunu[-r : þaiʀ : raistu : stain ...arþi : karþ] : auk : sunarla : sat : miþ : skinum : auk : han : entaþis : at : ulfshala : þa : [¶ han : hil(k)(i)...]}

\subsection{Text}

\bpg\bpa[0]%
Bót·mundr auk Bót·raifʀ auk Gunn·\emph{h}v\emph{at}r, þaiʀ raistu stain ... \textbf{†arþi†} garð, \\
auk \alst{s}unnar-la \hld\ \alst{s}at mið skinnum; \\
auk hann \alst{e}ndaðis \hld\ at \alst{U}lfs-hala, \\
þȧ hann \textbf{†hil(k)(i)†} ...\epa

\bpb Bootmund and Bootraf and Guthwhat, they raised the stone ... the enclosure, \\
and in the south sat with skins, \\
and he met his end by Wolfshale \\
when he \dots\epb\epg

\sectionline\clearpage


\section{G 208}

\begin{flushright}%
\textbf{Dating:} C12th

\textbf{Meter:} \Fornyrdislag
\end{flushright}%

\subsection{Introduction}

A damaged stone. TODO.  The same people also commissioned the likewise damaged G 207 above.

\subsection{Runes}

TODO: Runes%textrunic

\subsection{Transliteration}

\textbf{[b]u[t]muntr : auk : butraifʀ : auk : kunua-(r) [: -]... ...tu : stain : þ... ... ...[at : faþur : si]n : kuþ : h(i)albi : (s)elu : hans : auk : kus : muþiʀ : (b)etr : þen : uiʀ : biþia : kunin · ¶ -(a)... ...(a)-(n) (:)}

\subsection{Text}

\bpg\bpa[0]%
\startsub Bot·mundr ok Bót·raifʀ ok Gunn\emph{h}va\emph{t}r \emph{þaiʀ rais}tu stain þ\emph{enna eptiʀ} ... \emph{-hv}at, faður si̇nn.\endsub \\
\alst{G}uð hialpi sélu hans \hld\ ok \alst{G}uðs móðiʀ \\
\alst{b}etr þen víʀ \hld\ \alst{b}iðja kunnin ...\epa

\bpb Bootmund and Bootraf and Guthwhat, they raised this stone in memory of \dots-what, their father. \\
God help his soul and God’s mother, \\
better than we were able to pray. \dots\epb\epg

\sectionline\clearpage


\section{G 343}

\begin{flushright}%
\textbf{Dating:} C11th

\textbf{Meter:} \Fornyrdislag
\end{flushright}%

\subsection{Introduction}

TODO.

\subsection{Runes}

TODO: Runes%textrunic

\subsection{Transliteration}

§ A \textbf{...n : raisti : kubl : eftiʀ : hailkaiʀ : fa... ... ¶ (k)-... ... -aulu : hans : eu : miþan : uaralt : uakiʀ : ligʀ : merki : hier : yfiʀ : mani : þaim : aʀ : erfiki : iftiʀ -erþi ¶ ... auk : þorlaifʀ : þau/þaiʀ : ristu stain} \\
§ B \textbf{f-{}-{}-{}-{}-{}-{}-...}

\subsection{Text}

\bpg\bpa[0]%
\dots\ raisti kumbl eftiʀ Hail·gaiʀ, fa\emph{ður} \dots\ \emph{s}áulu hans. \\
\ind Ęy, miðan \alst{v}er-ald \alst{v}akiʀ, \\
liggʀ \alst{m}ęrki hięr \hld\ yfiʀ \alst{m}anni þaim, \\
aʀ \alst{ę}rfingi \hld\ \alst{ę}ptiʀ \emph{gi}ęrði. \\
\dots\ auk Þȯr·laifʀ, þau rístu stain.\epa

\bpb TODO: trans\epb\epg

\sectionline\clearpage


\section{G 370}

\begin{flushright}%
\textbf{Dating:} C11th

\textbf{Meter:} \Fornyrdislag
\end{flushright}%

\subsection{Introduction}

TODO.

\subsection{Runes}

TODO: Runes%textrunic

\subsection{Transliteration}

\textbf{uatar : auk hilkaiʀ : raistu : stain iftir hailka f-þur : sin : hn : uahʀ -istr : farin miþ uikikum}

\subsection{Text}

\bpg\bpa[0]%
\startsub \emph{H}vatarr ok H\emph{a}il·gaiʀ raistu stain ęptiʀ Hailga, f\emph{a}ður si̇nn.\endsub \\
Hann vaʀ \emph{\alst{v}}estr farinn \hld\ með \alst{v}íkingum.\epa

\bpb TODO: trans\epb\epg

\sectionline\clearpage


\section{Nä 29}

\begin{flushright}%
\textbf{Dating:} C11th

\textbf{Meter:} \Fornyrdislag
\end{flushright}%

\subsection{Introduction}

TODO.

\subsection{Runes}

TODO: Runes%textrunic

\subsection{Transliteration}

\textbf{bofriþr : let resa eftir biurn : bruþur : si(n) : han uar farin fultrekila}

\subsection{Text}

\bpg\bpa[0]%
Bó·friðr lét rêsa ęptiʀ Biorn, bróður si̇nn. \\
Hann var \alst{f}arinn \hld\ \alst{f}ull-dręngi⸗la.\epa

\bpb TODO: trans\epb\epg

\sectionline\clearpage


\section{Sm 5}

\begin{flushright}%
\textbf{Dating:} C11th

\textbf{Meter:} \Fornyrdislag
\end{flushright}%

\subsection{Introduction}

TODO.
The carver seems to struggle with consonant clusters.  Thus he spells 2b \emph{mestr} as \textbf{mesr} and 3b \emph{aldri} as \textbf{ali}(!).  Line 2 is found near-identically in Ög 77 and Sm 37 and line 3 in Vg 187.

\subsection{Runes}

TODO: Runes%textrunic

\subsection{Transliteration}

§ A \textbf{: kotr : sati : sten : þana : eftʀ : ketil :} \\
§ B \textbf{: sun : sin : han : faʀ :} \\
§ C \textbf{: mana : mesr o:niþikʀ : eʀ a : eklati : ali : tunþi}

\subsection{Text}

\bvg\bva[]%
Gø̂tr satti stên þęnna eptiʀ Kętil sun si̇nn. &
Hann vaʀ \alst{m}anna \hld\ \edtrans{\alst{m}estr ȯ·níðingʀ}{greatest un-nithing}{\Bfootnote{An instance of \emph{litotes}; Kettle was a very virtuous man.}} &
eʀ ȧ \alst{E}ng-landi \hld\ \edtext{\alst{a}l\emph{dr}i}{\Afootnote{\emph{\textbf{ali} Sm 5}}} tẏnði.\eva

\bvb Geat set this stone after Kettle, his son. \\
He was of men the greatest un-nithing, \\
who in England lost his life.\evb\evg

\sectionline\clearpage


\section{Sm 16}

\begin{flushright}%
\textbf{Dating:} C11th

\textbf{Meter:} \Fornyrdislag
\end{flushright}%

\subsection{Introduction}

A not very tall stone, inscribed on two sides.  Side A (1 \emph{Hró·stęinn} \dots\ 3 \emph{kęnni⸗ligt}) is marked with a large cross while side B (3 \emph{ęptir} \dots\ 5 \emph{ru̇n}a) is entirely unadorned.

The inscription uses the stung runes \emph{e} and \emph{y} and writes the vowel \emph{o} with the old \emph{ą}-rune.

\subsection{Runes}

TODO: Runes%textrunic

\subsection{Transliteration}

§ A \textbf{rostein · auk · eilifʀ · aki : auk · hakun : reisþu · þeiʀ · sueinaʀ · eftiʀ sin · ¶ faþur · kubl ¶ keni·likt ·} \\
§ B \textbf{ftiʀ · kala · tauþan : þy : mun · ko... ... -m kitit · uerþa · meþ · sin · lifiʀ · auk · stafiʀ · run}

\subsection{Text}

\bvg\bva[]%
Hró·stęinn auk \alst{Ęi}·lífʀ, \hld\ \alst{Ȧ}ki auk Hǫ́·kun, &
ręistu þęiʀ \alst{s}vęinaʀ \hld\ ęptiʀ \alst{s}i̇nn fǫður &
\alst{k}umbl \alst{k}ęnni⸗ligt \hld\ ęptiʀ \alst{K}ala dauðan. &
Þý mun \alst{g}ó\emph{ðs manns} \hld\ \emph{u}m \alst{g}etit verða, &
með \alst{s\emph{t}}\emph{ę}inn lifiʀ \hld\ ok \alst{st}afiʀ ru̇n\emph{a}.\eva

\bvb Rothstan and Anlif, Onick and Hathkin, \\
these lads raised after their father \\
a remarkable monument, after the dead Cale. \\
Therefore will the good man be spoken of \\
while the stone lives and the staves of the runes.\evb\evg

\sectionline\clearpage


\section{Sm 29}

\begin{flushright}%
\textbf{Dating:} C11th

\textbf{Meter:} \Fornyrdislag
\end{flushright}%

\subsection{Introduction}

TODO.

\subsection{Runes}

TODO: Runes%textrunic

\subsection{Transliteration}

\textbf{...r · rsþi · stin · eftiʀ · þur(k)iʀ · faþur · sin · saʀ · (e)(t)aþis · ą · -klanti ·}

\subsection{Text}

\bpg\bpa[0]%
...r rêsþi stên ęptiʀ Þȯr·gêʀ, fǫður si̇nn; \\
sáʀ \alst{ę}ndaðis \hld\ ȧ \emph{\alst{Ę}}nglandi.\epa

\bpb ... raised the stone after Thurgar, his father. \\
He met his end in England.\epb\epg

\sectionline\clearpage


\section{Sm 35 (Replösa)}

\begin{flushright}%
\textbf{Dating:} C11th

\textbf{Meter:} \Fornyrdislag
\end{flushright}%

\subsection{Introduction}

A stone still standing at its original location.  To its south stands a stone ship consisting of 12 stones.

\subsection{Runes}

TODO: Runes%textrunic

\subsection{Transliteration}

\textbf{: kutraþr : karþi : kubl : þisi : iftiʀ : astraþ : faþur : sin : þan : frita : ak : þih:na : bistan : iʀ a : fin:iþi : forþum : uf| |faʀi :}

\subsection{Text}

\bvg\bva[]%
Gaut·ráðr gęrði kumbl þęssi ęptiʀ Ǫ̇st·ráð, fǫður si̇nn. &
\alst{þ}ann frę̇nda \hld\ ok \alst{þ}egna bętstan,
eʀ ȧ \edtrans{\alst{F}inn·ęiði}{the Finnheath}{\Bfootnote{Finnveden, the folkland of Smalland in which the stone is located.}} \hld\ \alst{f}orðum of váʀi.\eva

\bvb Geatred made these monuments after Estred, his father, \\
the best of kinsmen and of thanes \\
who erstwhile lived upon the Finnheath.\evb\evg

\sectionline\clearpage


\section{Sm 37 (Rörbro)}

\begin{flushright}%
\textbf{Dating:} C11th

\textbf{Meter:} \Fornyrdislag
\end{flushright}%

\subsection{Introduction}

Sm 37 is a stone still standing at its original location, Rörbro in Nöttja parish in the Finnheath.  The inscription is tied to the non-metrical Sm 36, which stands just 1.5 km south of Sm 37, also at its original location.

Sm 36 is raised by an Eanwind (\textbf{uintʀ}) and Swain (\emph{Svęinn}) in memory of another Answer (\textbf{ąsmr}), but we are here met with a problem, since it is obviously not possible for two men to posthumously raise runestones for each other.  Assuming that Swain and Eanwind of Sm 36 were the sons of Answer (something which, though likely, is not stated explicitly), there are two likely explanations, both of which require a son named after his grandfather: (1) Answer I fathered Eanwind and Swain, who raised Sm 36 first.  Eanwind fathered Answer II, who later raised Sm 37 after him.  (2) Eanwind I was the father of Answer, who raised Sm 37 first in memory of him.  Answer in turn fathered Eanwind II and Swain, who later raised Sm 36 in memory of their father.

\subsection{Runes}

TODO: Runes%textrunic

\subsection{Transliteration}

\textbf{[:] a[sur] · karþi : kubl : þesi : eftiʀ : uit : faþur : sin · [h]an uaʀ : mana : mestr : uniþikʀ [:] uaʀ : i[n]tr : mataʀ : uk umun hats : kuþr · þi(k)[n] kus tru : kuþa : hafþi :}

\subsection{Text}

\bvg\bva[]%
Ǫssurr gęrði kumbl þęssi ęptiʀ Ęy\emph{n}d, fǫður si̇nn. &
Hann vaʀ \alst{m}anna \hld\ \alst{m}ęstr ȯ·níðingʀ, &
vaʀ \alst{y}ndr mataʀ \hld\ ok \alst{ȯ}·mun\emph{ʀ} hatrs. &
\alst{G}óðr þegn \alst{G}uðs-tró \hld\ \alst{g}óða hafði.\eva

\bvb Answer made these monuments after Eanwind, his father. \\
He was of men the greatest un-nithing, \\
was giving with food and late to hatred. \\
The good thane had good faith in God.\evb\evg

\sectionline\clearpage


\section{Sm 39}

\begin{flushright}%
\textbf{Dating:} C11th

\textbf{Meter:} \Fornyrdislag
\end{flushright}%

\subsection{Introduction}

A standing stone inscribed on two sides, one of which has a large cross.  The expression is formulaic; cf.~Sm 44, Sö 130, U 703, U 739, and U 805.  For “\inx[C]{good of meat}” cf.~\textlink{Havamal}[39].  The first line is not poetic.

\subsection{Runes}

TODO: Runes%textrunic

\subsection{Transliteration}

§ A \textbf{: kuni : sati : stin : þana : iftiʀ} \\
§ B \textbf{: suna : faþr : sin : milan : u... ... mataʀ kuþa... :}

\subsection{Text}

\bvg\bva[]%
Gunni satti stên þęnna eptiʀ Su̇na, fǫður si̇nn, &
\alst{m}il\emph{d}an o\emph{rða \hld\ auk} \alst{m}ataʀ góða\emph{n}.\eva

\bvb Guthe set this stone after Sown, his father, \\
generous in words and good of meat.\evb\evg

\sectionline\clearpage


\section{Sm 44}

\begin{flushright}%
\textbf{Dating:} C11th

\textbf{Meter:} \Fornyrdislag
\end{flushright}%

\subsection{Introduction}

A heavily damaged stone.  The inscription is more visible in Bautil on which the reading has to rely.
The expression is formulaic; cf.~Sm 39, Sö 130, U 703, U 739, and U 805.

\subsection{Runes}

TODO: Runes%textrunic

\subsection{Transliteration}

§ A \textbf{[: uim]ut[r] : [sa]ti : sten : þana ...} \\
§ B \textbf{[b]ruþr : s[in ...(i)n mi(l)... ... sina] ...} \\
§ A \textbf{m[a]tʀ : kuþan : [i] urþ[l]uf:i : at[ra : m...r]}

\subsection{Text}

\bvg\bva[]%
Vé·mundr satti stên þęnna \emph{ęptir} bróður si̇nn \edtrans{...\textbf{in}}{...}{\Bfootnote{The name of the brother cannot be restored, but the gap only leaves room for at most two runes, and so the main candidate names are Swayn (\textbf{*suin} \emph{Svênn}) and Stan (\textbf{*stin} \emph{Stênn}).}} &
\alst{m}il\emph{dan við} si̇nna \hld\ \emph{ok} \alst{m}ataʀ góðan, &
i̇́ \alst{o}rð-lofi \hld\ \alst{a}\emph{ll}ra m\emph{iklu}.\eva

\bvb Wighmund set this stone after his brother, ..., \\
generous with his kin and good of meat, \\
held in great praise by all.\evb\evg

\sectionline\clearpage


\section{Sm 45}

\begin{flushright}%
\textbf{Dating:} C11th

\textbf{Meter:} \Fornyrdislag
\end{flushright}%

\subsection{Introduction}

TODO.

\subsection{Runes}

TODO: Runes%textrunic

\subsection{Transliteration}

\textbf{· uest-n · karþ- · (k)ub- · þesi · efteʀ · esburn · bruþur · sin · uitrik · þesi · a uiki·muti}

\subsection{Text}

\bpg\bpa[0]%
\startsub Vé·stęinn gęrð\emph{i} kumb\emph{l} þęssi ęptiʀ Ę̇s·biorn, bróður si̇nn,\endsub \\
\alst{v}itring þęssi \hld\ ȧ \alst{v}ega-móti.\epa

\bpb Wighstan made these monuments after Osbarn, his brother, \\
this memento at the crossroads.\epb\epg

\sectionline\clearpage


\section{Sm 46}

\begin{flushright}%
\textbf{Dating:} C11th

\textbf{Meter:} \Fornyrdislag
\end{flushright}%

\subsection{Introduction}

TODO.

The stone is lost.

\subsection{Runes}

TODO: Runes%textrunic

\subsection{Transliteration}

\textbf{[...-nui krþi : kubl : þesi : iftiʀ suin : sun : sin : im ÷ itaþisk ou·tr i krikum]}

\subsection{Text}

\bpg\bpa[0]%
\startsub ...n·ví gęrði kumbl þęssi ęptiʀ Svên, sun si̇nn,\endsub \\
e\emph{ʀ} \alst{ę}ndaðisk \hld\ \alst{au}str i̇́ Grikkjum.\epa

\bpb ...wye made these monuments after Swayn, her son, \\
who met his end to the east in Greece.\epb\epg

\sectionline\clearpage


\section{Sm 60}

\begin{flushright}%
\textbf{Dating:} C11th

\textbf{Meter:} \Fornyrdislag
\end{flushright}%

\subsection{Introduction}

TODO.

\subsection{Runes}

TODO: Runes%textrunic

\subsection{Transliteration}

§ A \textbf{· sen · uk u·starki · karþu ·} \\
§ B \textbf{· kubl · þesi · efeiʀ · kuþmut · faþur · sin · uitrk · þasi · ą ueha·muti ·}

\subsection{Text}

\bpg\bpa[0]%
\startsub S\emph{v}ênn ok Ȯ·starki gęrðu kumbl þęssi ęptiʀ Guð·mund, fǫður si̇nn,\endsub \\
\alst{v}itring þási \hld\ ȧ \alst{v}ega-móti. \epa

\bpb Swayn and Unstark made these monuments after Godmund, their father, \\
this memento at the crossroads.\epb\epg

\sectionline\clearpage


\section{Sö 9 (Lifsinge)}

\begin{flushright}%
\textbf{Dating:} C11th

\textbf{Meter:} \Fornyrdislag
\end{flushright}%

\subsection{Introduction}

TODO.

\subsection{Runes}

TODO: Runes%textrunic

\subsection{Transliteration}

\textbf{barkuiþr × auk × þu : helka × raistu × stain × þansi : at · ulf : sun · sin · han × entaþis + miþ : ikuari +
¶ kuþ + hialbi + salu ulfs ×}

\subsection{Text}

\bpg\bpa[0]%
\startsub Bark·viðr ok þau Hęlga ręistu stęin þannsi at Ulf, sun si̇nn.\endsub \\
Hann \alst{ę}ndaðis \hld\ með \alst{I}ng·vari. \\
Guð hialpi sǫ́lu Ulfs.\epa

\bpb TODO\epb\epg

\sectionline\clearpage


\section{Sö 14 (Gåsinge)}

\begin{flushright}%
\textbf{Dating:} C11th

\textbf{Meter:} \Fornyrdislag
\end{flushright}%

\subsection{Introduction}

TODO.

The stone mentions a western expedition led by Cnut, most likely King Canute \parencite{Wessen1960}.  Other stones that mention this expedition are the non-metrical Ög 111 (see Ög 68 below), U 194, U 241, and U 344.

TODO: specific page reference in Wessén 1960.

\subsection{Runes}

TODO: Runes%textrunic

\subsection{Transliteration}

\textbf{rakna · raisti · stain · þansi · at · suin · buta · sit · auk · sifa · auk · r-knburk · at · sit · faþur · kuþ · hil[b]i · at · [hat]s · uit · iak · þet · uaʀ · sui- · uestr · miþ · kuti}

\subsection{Text}

\bpg\bpa[0]%
\startsub Ragna ręisti stęin þannsi at Svęin, bónda si̇nn, auk Sefa auk Ragn·borg at si̇nn fǫður. Guð hialpi and hans.\endsub \\
\alst{V}\emph{ę}it iak þęt vaʀ Svęinn \hld\ \alst{v}estr með K\emph{n}úti. \epa

\bpb Rainie raised this stone after Swayn, her husband, and Seve and Rainbury after their father. God help his spirit. \\
I know that Swayn was in the west with Cnut.\epb\epg

\sectionline\clearpage


\section{Sö 33}

\begin{flushright}%
\textbf{Dating:} C11th

\textbf{Meter:} \Fornyrdislag
\end{flushright}%

\subsection{Introduction}

TODO.

\subsection{Runes}

TODO: Runes%textrunic

\subsection{Transliteration}

\textbf{+ gnubha ÷ liþ : raisa : stain : þinsa : hibtiʀ : kulaif : bruþur sin ¶ han : antaþis : austr : at þikum}

\subsection{Text}

\bpg\bpa[0]%
Gnúpa lét ręisa stęin þęnnsa ęptiʀ Gullęif, bróður si̇nn. \\
Hann \alst{ę}ndaðis \hld\ \alst{au}str at þingum.\epa

\bpb TODO: trans\epb\epg

\sectionline\clearpage


\section{Sö 34--35 (Tjuvstigen)}

\begin{flushright}%
\textbf{Dating:} C11th--C12th

\textbf{Meter:} \Fornyrdislag
\end{flushright}%

\subsection{Introduction}

Sö 34 and 35 are two paired stones standing next to each other.  They are carved in the same style and the inscription on Sö 35 refers back to SÖ 34.  The last two lines of Sö 35 is not poetic.

\subsection{Runes}

TODO: Runes%textrunic

\subsection{Transliteration}

Sö 34: \textbf{styrlaugʀ · auk · hulmbʀ · staina · raistu · at · bryþr · sina · brau(t)u · nesta · þaiʀ · entaþus · i · austruiki · þurkil · auk sturbiarn þiaknaʀ · kuþiʀ} \\
Sö 35: \textbf{lit · igikeʀ · anan · raisa · stain · at · suni · sina · su[n·]a · kiarþi · kuþ · hialbi · ant · þaira × þuriʀ · hiu ·}

\subsection{Text}

\subsubsection{Sö 34}

\bvg\bva[]%
\alst{St}yr·laugʀ ok Holmbʀ \hld\ \alst{st}ęina ręistu &
at \alst{b}rǿðr si̇́na, \hld\ \edtrans{\alst{b}rautu nę́sta}{nearest to the road}{\Bfootnote{Namely the forest path where the stones still stand.  Note the similar wording in \textlink{Havamal}[72]/3--4.}}. &
Þęir \alst{ę}ndaðus \hld\ i̇́ \alst{au}str-vegi, &
\alst{Þ}ȯr·kęll ok Styr·biǫrn, \hld\ \alst{þ}iagnar góðir.\eva

\bvb Sturley and Holm raised the stones \\
after their brothers, nearest to the road. \\
They met their end on the Eastward way, \\
Thurkettle and Sturbarn, good thanes.\evb\evg

\subsubsection{Sö 35}

\bvg\bva[]%
Lét \alst{I}ngi·gęiʀʀ \hld\ \alst{a}nnan ręisa stęin &
at \alst{s}yni \alst{s}i̇́na, \hld\ \alst{s}ẏna giǫrði. &
Guð hialpi ǫnd þęira. &
Þóriʀ hjó.\eva

\bvb Inggar let another stone be raised \\
after his sons; he made it seen. \\
God help their spirit. \\
Thurgar hewed.\evb\evg

\sectionline\clearpage


\section{Sö 41}

\begin{flushright}%
\textbf{Dating:} C11th

\textbf{Meter:} \Fornyrdislag
\end{flushright}%

\subsection{Introduction}

TODO.

\subsection{Runes}

TODO: Runes%textrunic

\subsection{Transliteration}

\textbf{tati + iok + eftiʀ + faþur + sin × skaka × mirki| |it + mikla + man + (i)auʀn}

\subsection{Text}

\bpg\bpa[0]%
TODO: text\epa

\bpb TODO: trans\epb\epg

\sectionline\clearpage


\section{Sö 47}

\begin{flushright}%
\textbf{Dating:} C11th

\textbf{Meter:} \Fornyrdislag
\end{flushright}%

\subsection{Introduction}

TODO.

\subsection{Runes}

TODO: Runes%textrunic

\subsection{Transliteration}

\textbf{...(ʀ)-kʀ : kiarþi : kuml : þat:si : eftiʀ : ąsmunt : sun : sin + han : is : krafin : ą · ku... <rauʀ> : <uart/fart> : at : ryʀ:iks : sun}

\subsection{Text}

\bpg\bpa[0]%
TODO: text\epa

\bpb TODO: trans\epb\epg

\sectionline\clearpage


\section{Sö 54}

\begin{flushright}%
\textbf{Dating:} C11th

\textbf{Meter:} \Fornyrdislag
\end{flushright}%

\subsection{Introduction}

TODO.

\subsection{Runes}

TODO: Runes%textrunic

\subsection{Transliteration}

§ A \textbf{þurstain : auk : aystain : auk : natfari : raistu : stain : at : finuiþ : auk : ulif : þurkil buruþr sin ¶ uaru · aliʀ · uikiks · suniʀ ¶ latburniʀ · man · litu · rita · stain} \\
§ B \textbf{stainkil · rist · runaʀ}

\subsection{Text}

\bpg\bpa[0]%
\startsub Þȯr·stęinn auk Ęy·stęinn auk Nátt·fari ręistu stęin at Finn·við auk Ȯ·lęif, Þȯr·kęl, brǿðr si̇́n\emph{a}.\endsub \\
\alst{V}ǫ́ru alliʀ \hld\ \alst{V}íkings syniʀ; \\
\alst{l}and-burniʀ męnn \hld\ \alst{l}étu rétta stęin. \\
Stęin·kęll ręist ru̇naʀ. \epa

\bpb Thurstan and Eanstan and Nightfarer raised the stone after Finnwood and Anlaf [and] Thurkettle, their brothers. \\
They were all the sons of Wiking; \\
land-born men had the stone erected. \\
Stankettle carved the runes.\epb\epg

\sectionline\clearpage


\section{Sö 55 (Bjudby)}

\begin{flushright}%
\textbf{Dating:} C11th

\textbf{Location:} Bjudby, Blacksta parish, Oppunda hundred

\textbf{Meter:} \Fornyrdislag
\end{flushright}%

\subsection{Introduction}

TODO: INTRODUCTION.

The language is later than most other England stones.  The os-rune is used for \textipa{/o/}, palatal \emph{ʀ} has dropped after all consonants, and stung \emph{e} is used.

\subsection{Runes}

TODO: Runes%textrunic

\subsection{Transliteration}

\textbf{þorstain (l)(i)... ...sa : stain : þena : ... sik : sialfan : auk : sun : sin : hefni : uaʀ til : enklans : ukr : trenkr : farin : uarþ : þa · haima : at : harmi tauþr ¶
kuþ hialbi : sialu : þaima ¶
bruni : auk : sloþi : þaiʀ ...(u) stan þena}

\subsection{Text}

\bpg\bpa[0]%
\startsub Þȯr·stęinn \emph{lét ręi}sa stęin þęnna \emph{ęptiʀ} sik sialfan auk sun si̇nn Hęfni.\endsub \\
Vaʀ til \alst{Ę}ng·lands \hld\ \alst{u}ngr dręngr farinn, \\
varð þȧ \alst{h}ęima \hld\ at \alst{h}armi dauðr. \\
Guð hialpi sialu þęi\emph{ʀ}a. \\
Brúni auk Slóði þęiʀ \emph{ríst}u stęin þęnna.\epa

\bpb Thurstan had this stone raised after himself and his son Hefner. \\
To England had the young dreng travelled, \\
then at home he died tragically. \\
God help their souls. \\
Brown and Sloth they carved this stone.\epb\epg

\sectionline\clearpage


\section{Sö 56 (Fyrby)}

\begin{flushright}%
\textbf{Dating:} C11th--C12th

\textbf{Meter:} \Fornyrdislag
\end{flushright}%

\subsection{Introduction}

TODO: INTRODUCTION.

\subsection{Runes}

TODO: Runes%textrunic

\subsection{Transliteration}

\textbf{iak · uait : hastain : þa : hulmstain : bryþr · menr : rynasta : a : miþkarþi : setu : stain : auk : stafa : marga eftiʀ · fraystain · faþur · sin ·}

\subsection{Text}

\bvg\bva[]%
Iak vęit \alst{H}á·stęin \hld\ þá \alst{H}olm·stęin brǿðr &
\alst{m}ęnnr rẏnasta \hld\ ȧ \alst{M}ið·garði &
sęttu \alst{st}ęin \hld\ auk \alst{st}afa marga &
eptir \alst{F}ręy·stęin \hld\ \alst{f}ǫður si̇nn.\eva

\bvb I know Highstan and Holmstan, the brothers, \\
the most rune-cunning men in Middenyard. \\
They set the stone and many staves \\
after Freastan, their father.\evb\evg

\sectionline\clearpage


\section{Sö 61}

\begin{flushright}%
\textbf{Dating:} C11th

\textbf{Meter:} \Fornyrdislag
\end{flushright}%

\subsection{Introduction}

TODO.

\subsection{Runes}

TODO: Runes%textrunic

\subsection{Transliteration}

\textbf{þorstain : lit · þina : rita : stain : e(f)ila : stetr : eftiʀ : þorbiarn : salui · auk : simiþr : at · sen boroþur}

\subsection{Text}

\bvg\bva[]%
\alst{Þ}ȯr·stęinn lét \hld\ \alst{þ}ęnna rétta &
\alst{st}ęin; ę́fila \hld\ \alst{st}ęndr ęptiʀ Þȯr·biorn. &
\alst{S}ǫlvi auk \alst{S}miðr \hld\ at \alst{s}i̇nn bróður. \eva

\bvb Thurstan let erect this stone; \\
fittingly it stands after Thurbarn. \\
Sallow and Smith [raised it] after their brother.\evb\evg

\sectionline\clearpage


\section{Sö 65 (Djulefors)}

\begin{flushright}%
\textbf{Dating:} C11th--C12th

\textbf{Meter:} \Fornyrdislag\ with hendings in the b-verses
\end{flushright}%

\subsection{Introduction}

A standing stone inscribed on one side with a large cross.  Already on the earliest depictions the stone was damaged, but an even larger part has now gone missing.  Other stones that mention \inx[P]{Longbeardland} (Lombardy, northern Italy) are the paired U 133 and U 141; it may also be mentioned on Sö Fv1954;22, although only the first two letters \textbf{la} survive.

The meter is highly unusual for runic Swedish poetry by its reliance on hendings (in line 2 an ethel-hending \emph{arð- : barð-}, in line 3 a shot-hending \emph{land- : ęnd-}).  Although unparallelled in Sweden, the meter is documented by Snorre (TODO).

\subsection{Runes}

TODO: Runes%textrunic

\subsection{Transliteration}

\textbf{[inka : raisti : stain : þansi : at : ulai](f) : sin : [a...k] : han : austarla : arþi : barþi : auk : ą : lakbarþilanti : [anlaþis +]}

\subsection{Text}

\bvg\bva[]%
Inga ręisti stęin þannsi at Ȯ·lęif si̇nn \textbf{a...k}. &
Hann \alst{au}star⸗la \hld\ \edtrans{\alst{a}rði barði}{ploughed with the prow}{\Bfootnote{I.e. ‘sailed’.  The same verse is found in an anonymous \Kviduhattr-line from the Third Grammatical Treatise (Anon (\emph{TGT}) 14 in \Skp\ 3), which reads: \emph{sá’s af Ís-landi \hld\ arði barði} ‘he who [away] from Iceland ploughed with the prow’.}} &
auk ȧ \alst{L}ang-barði- \hld\ \alst{l}andi an\emph{d}aðis.\eva

\bvb Inge raised this stone after Anlaf, her ... \\
Easterly he ploughed with the prow, \\
and in Longbeardland met his end.\evb\evg

\sectionline\clearpage


\section{Sö 85}

\begin{flushright}%
\textbf{Dating:} C11th

\textbf{Meter:} \Fornyrdislag
\end{flushright}%

\subsection{Introduction}

TODO.

\subsection{Runes}

TODO: Runes%textrunic

\subsection{Transliteration}

\textbf{: ansuar : auk : ern... ... [: faþur sin : han : enta]þis : ut i : krikum (r)uþr : -{}-{}-...unk-{}-{}-{}-{}-{}-an-{}-{}-{}-{}-}

\subsection{Text}

\bpg\bpa[0]%
Ansvarr ok Ęrn\dots, fǫður si̇nn. \\
Hann \alst{ę}ndaðis \hld\ \alst{ú}t i̇́ Grikkum. \\
\textbf{†\dots (r)uþr:\dots unk\dots an\dots †} \epa

\bpb TODO: trans\epb\epg

\sectionline\clearpage


\section{Sö 106}

\begin{flushright}%
\textbf{Dating:} C11th

\textbf{Meter:} \Fornyrdislag
\end{flushright}%

\subsection{Introduction}

The family involved in raising this stone were fairly prolific in raising runestones, for they also raised U 617 and the non-metrical Sö 101.  Based on these 3 inscriptions, their family tree may be reconstructed as follows:

\hspace{0.5cm}

\noindent \resizebox{\textwidth}{!}{
\begin{tikzpicture}[
    node distance = 1.2cm and 0.8cm,
    person/.style = {draw, rectangle, minimum width=2.2cm, minimum height=0.8cm, align=center},
    line/.style = {draw, thick},
    marriage/.style = {draw, thick, dashed}
]

\node[person] (sigridr) {\female Sig·ríðr};
\node[person] (spiut) [left=1.5cm of sigridr] {\male Spiut};
\node[person] (sygrodr) [right=1.5cm of sigridr] {\male Syg·røðr};
\node[person] (gautr) [right=0.8cm of sygrodr] {\male Gautr};
\node[person] (ginnlaug) [right=0.8cm of gautr] {\female Ginn·laug};
\node[person] (ossur) [right=1.5cm of ginnlaug] {\male Ǫssurr};

\node[person] (ormr) [above=1.5cm of sigridr] {\male Ormr};
\node[person] (holmgeirr) [above=1.5cm of gautr] {\male Holm·gęiʀʀ};
\node[person] (hakon) [above=1.5cm of ossur] {\male Hǫ́kon jarl};

\path (spiut) -- (sigridr) coordinate[midway] (mid_ss);
\node[person] (alrikr) [below=1.5cm of mid_ss] {\male Al·ríkr};

\draw[line] (ormr) -- (sigridr);
\draw[line] (hakon) -- (ossur);

\coordinate (h_down) at ($(holmgeirr.south) + (0,-0.6)$);
\draw[line] (holmgeirr.south) -- (h_down);
\draw[line] (sygrodr.north) |- (h_down) -| (ginnlaug.north);
\draw[line] (gautr.north) -- (gautr.north |- h_down);

\draw[line] (spiut) -- (sigridr) coordinate[midway] (u1);
\draw[line] (u1) -- (alrikr);

\draw[marriage] (sigridr) -- (sygrodr) node[midway, above, font=\scriptsize] {};
\draw[marriage] (ginnlaug) -- (ossur) node[midway, above, font=\scriptsize] {};

\end{tikzpicture}}

\subsection{Runes}

TODO: Runes%textrunic

\subsection{Transliteration}

\textbf{alrikʀ ¤ raisti ¤ stain × sun × siriþaʀ × at × sin faþur × sbiut × ¶ × saʀ × uisitaula × um × uaʀit : hafþi × burg × um brutna : i : auk × um barþa +× firþ × han × kar(s)aʀ + kuni + alaʀ ×}
%TODO: check ¤

\subsection{Text}

\bpg\bpa[0]%
\startsub Al·ríkʀ ręisti stęin, sunn Sig·ríðaʀ, at si̇nn fǫður Spiut;\endsub \\
sáʀ \alst{v}estar⸗la \hld\ um \alst{v}aʀit hafði, \\
\alst{b}org um \alst{b}rutna \hld\ i̇́ ok um \alst{b}arða; \\
fęrð hann \alst{k}a\emph{st}raʀ \hld\ \alst{k}unni allaʀ.\epa

\bpb TODO: trans\epb\epg

\sectionline\clearpage


\section{Sö 122 (Skresta)}

\begin{flushright}%
\textbf{Dating:} C11th

\textbf{Meter:} \Fornyrdislag
\end{flushright}%

\subsection{Introduction}

A standing stone at Skresta.  The semantic division between the words \emph{dauðan} and \emph{Ǫ̇s·gautr} is marked in the inscription by a large space.

The commissioner, Osgeat, also raised another stone in memory of Highstan, the unmetrical Sö 123.  That stone also stands at Skresta and reads: \textbf{ąskutr : rsþi : stan : þansi : at : hastin} \emph{Ǫ̇s·gautr r}ęi\emph{sþi stęin þannsi at Há·stęin} ‘Osgeat raised this stone after Highstan.’

\subsection{Runes}

TODO: Runes%textrunic

\subsection{Transliteration}

\textbf{: stain : stanr : it : histain : raisþi stalfʀ : faþiʀ : at : sun tauþan ąskutr : kiarþi : tre}

\subsection{Text}

\bvg\bva[]%
Stęinn stęndr \emph{a}t H\emph{á}·stęin; &
ręisþi \alst{s}\emph{i}alfʀ faðiʀ \hld\ at \alst{s}un dauðan. &
Ǫ̇s·gautr giarði tré.\eva

\bvb TODO: trans\evb\evg

\sectionline\clearpage


\section{Sö 126}

\begin{flushright}%
\textbf{Dating:} C11th

\textbf{Meter:} \Fornyrdislag
\end{flushright}%

\subsection{Introduction}

TODO.

\subsection{Runes}

TODO: Runes%textrunic

\subsection{Transliteration}

\textbf{hu(l)(m)(f)riþ · ilin-{}-r · [þ]aʀ · litu · hakua · stain · eftiʀ ¶ eskil · faþur · sin · han · trauh · orustu · i · austru[i]hi ¶ aþaa · fulks · krimʀ · fala · orþi}

\subsection{Text}

\bpg\bpa[0]%
\startsub Holm·fríðr, \textbf{†ilin?\dots r†}, þáʀ létu hǫggva stęin ęftiʀ Ę̇s·kęl, fǫður si̇nn.\endsub \\
Hann draug \alst{o}rrustu \hld\ i̇́ \alst{au}str-vegi, \\
áðan \alst{f}olks Grímʀ \hld\ \alst{f}alla orði.\epa

\bpb TODO: trans\epb\epg

\sectionline\clearpage


\section{Sö 130}

\begin{flushright}%
\textbf{Dating:} C11th--C12th

\textbf{Meter:} \Fornyrdislag
\end{flushright}%

\subsection{Introduction}

A standing stone. TODO.  The expression is formulaic; cf.~Sm 39, Sm 44, U 703, U 739, and U 805.

\subsection{Runes}

TODO: Runes%textrunic

\subsection{Transliteration}

§ A \textbf{fiuriʀ : kirþu : at : faþur : kuþan : tyrþ : trikela : at : tumara : miltan : urþa uk : mataʀ kuþan : þat · (u)-(h)-{}-{}-(u){}-{}-(u)(k)(þ)} \\
§ PB \textbf{h\textasciicircum a l\textasciicircum f kirþu <o>} \\
§ QB \textbf{h\textasciicircum a <h> kirþu <o/þ>...}

\subsection{Text}

\bvg\bva[]%
\alst{F}iuriʀ gęrðu \hld\ at \alst{f}ǫður góðan &
\alst{d}ýrð \alst{d}ręngi⸗la \hld\ at \alst{D}ómara &
\edtrans{\alst{m}ildan orða \hld\ ok \alst{m}ataʀ góðan}{generous in words and good of meat}{\Bfootnote{The same line is found in Sm 39.}}. &
Þat \dots\eva

\bvb Four men made after their good father, \\
honourably a mark of praise after Doomer \\
generous in words and good of meat. \\
This \dots\evb\evg

\sectionline\clearpage


\section{Sö 131}

\begin{flushright}%
\textbf{Dating:} C11th

\textbf{Meter:} \Fornyrdislag
\end{flushright}%

\subsection{Introduction}

TODO.

\subsection{Runes}

TODO: Runes%textrunic

\subsection{Transliteration}

\textbf{: sbiuti : halftan : þaiʀ : raisþu : stain : þansi : eftiʀ : skarþa : ¶ bruþur sin : fur : austr : hiþan : miþ : ikuari : ą sirk¶lanti : likʀ : sunʀ iuintaʀ}

\subsection{Text}

\bpg\bpa[0]%
\startsub Spiuti, Half·dan, þęiʀ ręisþu stęin þannsi ęptiʀ Skarða, bróður si̇nn.\endsub \\
Fór \alst{au}str heðan \hld\ með \alst{I}ng·vari; \\
ȧ \alst{S}ęrk-landi liggʀ \hld\ \alst{s}onʀ Ęy·vindaʀ.\epa

\bpb TODO: trans\epb\epg

\sectionline\clearpage


\section{Sö 136}

\begin{flushright}%
\textbf{Dating:} C11th

\textbf{Meter:} \Fornyrdislag
\end{flushright}%

\subsection{Introduction}

The stone is lost.

The Swain on the inscription has been identified with the one mentioned on Sö 165 in the same district; the difficult form \textbf{†sluia†} may perhaps be identified with the Slode (\emph{Slóði}) on Sö 2 and 141.

\subsection{Runes}

TODO: Runes%textrunic

\subsection{Transliteration}

\textbf{[: suain : iuk : sluia · þaiʀ : raisþu : -{}-{}-... ... ...-nu · at : faþur : sin : hirsi · uksniauin ian uas : unt hifni bistr]}

\subsection{Text}

\bpg\bpa[0]%
\startsub Svęinn ok \textbf{†sluia†}, þęiʀ ręisþu \textbf{†?-{}-{}.-..-nu†} at fǫður si̇nn,\endsub
\alst{h}ęrsi \emph{\alst{h}}ug-snia\emph{lla}n, \hld\ ian vas und \alst{h}ifni bętztr.\epa

\bpb TODO: trans\epb\epg

\sectionline\clearpage


\section{Sö 137 (Aspa)}

\begin{flushright}%
\textbf{Dating:} C11th

\textbf{Meter:} \Fornyrdislag
\end{flushright}%

\subsection{Introduction}

TODO.

The last line of side B (\textbf{sunr þat} ... \textbf{maʀ}) is, with the exception of the final word \textbf{maʀ}, enciphered.  The first 7 runes, \textbf{sunrþat}, are written along a single vertical stave, while \textbf{rakna} ... \textbf{i(f)} are written in staveless runes.  This line is difficult to read, employs no word separators, and has not been convincingly deciphered.

Weeper is also mentioned on the unmetrical Sö 138, likewise from Aspa, which is dedicated to his unnamed son.  In that inscription Thurie is named by her full name, Thurun (\emph{Þȯr·unn}).  It is unclear whether this Weeper is the same as the one mentioned on Sö Fv1982;235.

\subsection{Runes}

TODO: Runes%textrunic

\subsection{Transliteration}

§ A \textbf{þura : raisþi : stin : þ-{}-si at : ubi : buanti : sin} \\
§ BP \textbf{: stain : saʀ:si : stanr : at : ybi : ą þik·staþi : at ¶ : þuru : uar : han : uestarla : uakti : karla ¶ [sa þar] · sunr þaþ · traknasuatau(k)i(f)maʀ[sua]} \\
§ BQ \textbf{: stain : saʀ:si : stanr : at : ybi : ą þik·staþi : at ¶ : þuru : uar : han : uestarla : uakti : karla ¶ · sunr þat raknasuatau-i(f) maʀ}

\subsection{Text}

\bpg\bpa[0]%
TODO: text\epa

\bpb TODO: trans\epb\epg

\sectionline\clearpage


\section{Sö 145}

\begin{flushright}%
\textbf{Dating:} C11th

\textbf{Meter:} \Fornyrdislag
\end{flushright}%

\subsection{Introduction}

TODO.

\subsection{Runes}

TODO: Runes%textrunic

\subsection{Transliteration}

§ A \textbf{[: tusti : austin : þiʀ : raistu : at : tuka : sniʀ : kiarþu · at : san :: faþur :: snialan :]} \\
§ B \textbf{[: tuki : ati : ru harfan : krimulfu :: ati : hafan : iu : ata i :: uþuli]}

\subsection{Text}

\bpg\bpa[0]%
Tósti, Ęy·stęinn, þęiʀ ręistu at Tóka, \\
\alst{s}\emph{y}niʀ giǫrðu \hld\ at \alst{s}\emph{i}nn fǫður snjallan. \\
Tóki átti \emph{b}ý ha\emph{l}fan, Grím·ulf\emph{ʀ} átti ha\emph{l}fan \emph{b}ý, átti í óðuli.\epa

\bpb TODO: trans\epb\epg

\sectionline\clearpage


\section{Sö 148}

\begin{flushright}%
\textbf{Dating:} C11th

\textbf{Meter:} \Fornyrdislag
\end{flushright}%

\subsection{Introduction}

TODO.

The words \textbf{han uas antaþ austr} are carved with staveless runes.

\subsection{Runes}

TODO: Runes%textrunic

\subsection{Transliteration}

\textbf{þiuþulfʀ : bui : þaiʀ : raisþu : stain þansi : at : farulf : faþur : sin : han uas antaþ austr i kaþ(u)(m)}

\subsection{Text}

\bpg\bpa[0]%
Þiuð·ulfʀ, Bói, þęiʀ ręisþu stęin þannsi at Fǫr·ulf, fǫður si̇nn. \\
Hann vas \alst{a}ndað\emph{r} \hld\ \alst{au}str i̇́ Gǫ\emph{r}ðum.\epa

\bpb Thedwolf, Bowe---they raised this stone after Farulf, their father. \\
He met his end to the east in the Yards.\epb\epg

\sectionline\clearpage


\section{Sö 154 (Skarpåker)}

\begin{flushright}%
\textbf{Dating:} C11th

\textbf{Meter:} \Fornyrdislag
\end{flushright}%

\subsection{Introduction}

A stone.

The couplet at the end, expressing a father’s grief for his son, serves as a good example of the Wiking Age preoccupation with the End Times.  The stone is decorated with a cross, but the text has no signs of Christian influence, and the language is traditional.

\subsection{Runes}

TODO: Runes%textrunic

\subsection{Transliteration}

\textbf{kunar : raisþi : stain : þansi : at : lyþbiurn : sun : sin : iarþ sal rifna uk ubhimin}

\subsection{Text}


\bvg\bva[]%
Gunn·arr ręisti stęin þannsi at Lýð·biorn, son si̇nn. &
\edtext{\alst{Iǫ}rð \edtrans{sal}{shall}{\Bfootnote{An early example of a dialectal form of \emph{skal} ‘id.,’ cf.~dialectal Swedish \emph{sa}.}} rifna \hld\ ok \alst{u}pp-himinn.}{\Bfootnote{The verse line (with the exception of \textbf{h} and \textbf{m} in \textbf{ubhimin} \emph{upp-himinn}) is written with staveless runes and without punctuation.  Since there is clearly place for more runes on the stone the reason is unlikely to be strictly economical, especially considering that staveless runes are very rare in the Wiking Age outside of Helsingland and also used especially for poetry in other inscriptions from Södermanland (Sö 148 above).  Rather the motivation is almost certainly to encrypt the verse and thereby accord it greater sanctity, perhaps also to obscure a reference to the pagan eschatology at a time when Christianity was increasingly hegemonic.

For the expression cf.~\textlink{Voluspa}[50], Arn \emph{Hryn} 3/7--8 (in \Skp\ 2, see also note there): \emph{meiri verði þinn an þeira \hld\ þrifnuðr allr, und’s himinn rifnar.} ‘greater than theirs may all thy wealth be until heaven rends.’}\lemma{Iǫrð \dots\ upphiminn ‘Earth \dots\ Up-heaven’}\Bfootnote{A formulaic merism; see Index for other occurrences.}}\eva

\bvb Guther raised this stone after Leodbarn, his son. \\
Earth shall rend, and Up-heaven.\evb\evg

\sectionline\clearpage


\section{Sö 159}

\begin{flushright}%
\textbf{Dating:} C11th

\textbf{Meter:} \Fornyrdislag
\end{flushright}%

\subsection{Introduction}

TODO.

\subsection{Runes}

TODO: Runes%textrunic

\subsection{Transliteration}

\textbf{: ikialtr : ak : aluiʀ : raisþu : stain : þansi : at : þurbiurn : faþur : sin : han uaistr hafʀ/(a)fʀ uf/(u)f uaʀit leki ¶ rorikʀ · kumytr biu · kunlaifʀ hiuk ru-aʀ}

\subsection{Text}

\bpg\bpa[0]%
\startsub Ingi·aldr auk Ǫl·víʀ ręisþu stęin þannsi at Þȯr·biorn, fǫður si̇nn.\endsub \\
Hann \alst{v}estr hafiʀ \hld\ of \alst{v}aʀit lęngi. \\
\emph{H}rǿ·ríkʀ, Gu\emph{ð}·mundr \textbf{†biu†}. Gunn·lęifʀ hiogg ru̇\emph{n}aʀ.\epa

\bpb Ingeld and Alewigh raised this stone after Thurbarn, his father. \\
He has been in the west for a long time. \\
Rethric, Godmund \dots. Guthlaf hewed the runes.\epb\epg

\sectionline\clearpage


\section{Sö 163}

\begin{flushright}%
\textbf{Dating:} C11th

\textbf{Meter:} \Fornyrdislag
\end{flushright}%

\subsection{Introduction}

TODO.

\subsection{Runes}

TODO: Runes%textrunic

\subsection{Transliteration}

\textbf{þruʀikr : stain : at : suni : sina : sniala : trakia : for : ulaifr : i : krikium : uli : sifti :}

\subsection{Text}

\bvg\bva[]%
\edtrans{Þrý·ríkʀ stęin}{Thriric [raised] the stone}{\Bfootnote{This line is an orphan.  The verb is left out.}} &
at \alst{s}yni \alst{s}i̇́na, \hld\ \alst{s}nialla dręngja; &
fór Ȯ·lęifʀ i̇́ \alst{G}rikkjum, \hld\ \alst{g}ulli skipti.\eva

\bvb Thriric [raised] the stone \\
after his sons, valiant drengs; \\
Anlaf went to Greece, dealt out gold.\evb\evg

\sectionline\clearpage


\section{Sö 164}

\begin{flushright}%
\textbf{Dating:} C11th

\textbf{Meter:} \Fornyrdislag
\end{flushright}%

\subsection{Introduction}

TODO.

\subsection{Runes}

TODO: Runes%textrunic

\subsection{Transliteration}

§ P \textbf{kuþbirn : uti : þaiʀ r(a)isþu : stan þansi : at : kuþmar : f(a)þur : sin : stuþ : triki:l(a) : i · stafn skibi : likʀ uistarla uf huln sar tu :} \\
§ Q \textbf{kuþbirn : uti : þaiʀ r(a)isþu : stan þansi : at : kuþmar : f(a)þur : sin : stuþ : triki:l(a) : i · stafn skibi : likʀ uistarla <uf> <hul>n bar-{}-}

\subsection{Text}

\bpg\bpa[0]%
TODO: text\epa

\bpb TODO: trans\epb\epg

\sectionline\clearpage


\section{Sö 165}

\begin{flushright}%
\textbf{Dating:} C11th

\textbf{Meter:} \Fornyrdislag
\end{flushright}%

\subsection{Introduction}

TODO. The Swain on the stone may be the same person as the one mentioned in Sö 136.

\subsection{Runes}

TODO: Runes%textrunic

\subsection{Transliteration}

§ P \textbf{kuþrun : raisti : stain : at : hiþin : uaʀ : nafi suais : uaʀ : han : ¶ : i : krikum iuli skifti ¶ : kristr : hialb : ant : kristunia :} \\
§ Q \textbf{kuþrun : raisti : stain : at : hiþin : uaʀ : nafi suais : uaʀ : han : ¶ : i : krikum iuli skifti ¶ : kristr : hialb : ant : krist uni a :}

\subsection{Text}

\bpg\bpa[0]%
TODO: text\epa

\bpb TODO: trans\epb\epg

\sectionline\clearpage


\section{Sö 171}

\begin{flushright}%
\textbf{Dating:} C11th

\textbf{Meter:} \Fornyrdislag
\end{flushright}%

\subsection{Introduction}

TODO.

\subsection{Runes}

TODO: Runes%textrunic

\subsection{Transliteration}

\textbf{(i)nk(i)f(a)[s]tr · l[i](t) (h)(a)ku... st(a)...n · eftiʀ · sihuiþ · faþ-r · si[n · han · fial · i h]ul(m)[karþi · skaiþaʀ · uisi mi]þ · ski...ra}

\subsection{Text}

\bpg\bpa[0]%
TODO: text\epa

\bpb TODO: trans\epb\epg

\sectionline\clearpage


\section{Sö 173 (Tystberga)}

\begin{flushright}%
\textbf{Dating:} C11th

\textbf{Meter:} \Fornyrdislag
\end{flushright}%

\subsection{Introduction}

TODO.

The stone is closely connected to the non-metrical Sö 374, which presently stands beside it.  Sö 374 is somewhat younger, and was raised by Mishie (there \textbf{musku} gen.~sg.) after Manny.  The woman Mishie also raised the non-metrical Sö 13 with Wiking, perhaps her husband.

\subsection{Runes}

TODO: Runes%textrunic

\subsection{Transliteration}

§ A \textbf{mes:kia : a(u)[k :] (m)an(i) : litu : rasa : ku[(m)(l) : þausi : at : b]ruþur · (s)in : hr(u)þkaiʀ · auk : faþur sin hulm:stain ·} \\
§ B \textbf{· han hafþi · ystarla u(m) : uaʀit · lenki : tuu : a:ustarla : meþ : inkuari}

\subsection{Text}

\bpg\bpa[0]%
Myskja auk Manni létu rę\emph{i}sa kumbl þausi at bróður si̇nn Hróð·gęiʀ fǫður Holm·stęin. \\
Hann hafði \alst{v}estar⸗la \hld\ um \alst{v}aʀit lęngi; \\
dóu \alst{au}star⸗la \hld\ með \alst{I}ng·vari.\epa

\bpb TODO: trans\epb\epg

\sectionline\clearpage


\section{Sö 174 (Aspö)}

\begin{flushright}%
\textbf{Dating:} C11th

\textbf{Meter:} \Fornyrdislag
\end{flushright}%

\subsection{Introduction}

TODO.

\subsection{Runes}

TODO: Runes%textrunic

\subsection{Transliteration}

\textbf{[ub]lubʀ · lit · kira : kuml : likhus : auk : bru · at sun sin : biurn : uaʀ trebin : a : kut:lanti : þy : lit : fiur · sit : fluþu : kankiʀ : þaiʀ uiþ[ulkuʀ] : uiltu iki halta : guþ : hilbi : anta : hans}

\subsection{Text}

\bpg\bpa[0]%
TODO: text\epa

\bpb TODO: trans\epb\epg

\sectionline\clearpage


\section{Sö 179 (Gripsholm)}

\begin{flushright}%
\textbf{Dating:} C11th

\textbf{Meter:} \Fornyrdislag
\end{flushright}%

\subsection{Introduction}

TODO: INTRODUCTION.

The three-line stanza is a biographical addition following a typical prose memorial formula.

%TODO: The stone has been connected to Sö 278 and 279, but on what basis?

\subsection{Runes}

᛭ ᛏᚢᛚᛅ ᛬ ᛚᛁᛏ ᛬ ᚱᛅᛁᛋᛅ ᛬ ᛋᛏᛅᛁᚾ ᛬ ᚦᛁᚾᛋᛅᛏ ᛬ ᛋᚢᚾ ᛬ ᛋᛁᚾ ᛬ ᚼᛅᚱᛅᛚᛏ ᛬ ᛒᚱᚢᚦᚢᚱ ᛬ ᛁᚾᚴᚢᛅᚱᛋ ᛬ ᚦᛅᛁᛦ ᚠᚢᚱᚢ ᛬ ᛏᚱᛁᚴᛁᛚᛅ ᛬ ᚠᛁᛅᚱᛁ ᛬ ᛅᛏ ᛬ ᚴᚢᛚᛁ ᛬ ᛅᚢᚴ ᛬ ᛅ᛬ᚢᛋᛏᛅᚱᛚ ᛅᚱ᛬ᚾᛁ ᛬ ᚴᛅᚠᚢ ᛬ ᛏᚢᚢ ᛬ ᛋᚢᚾᛅᚱ᛬ᛚᛅ ᛬ ᛅᛋᛁᚱᚴ᛬ᛚᛅᚾ᛬ᛏᛁ%textrunic

\subsection{Transliteration}

\textbf{× tula : lit : raisa : stain : þinsa| |at : sun : sin : haralt : bruþur : inkuars : þaiʀ furu : trikila : fiari : at : kuli : auk : a:ustarla| |ar:ni : kafu : tuu : sunar:la : a sirk:lan:ti}

\subsection{Text}

\bpg\bpa[0]%
\startsub Tóla lét ręisa stęin þennsa at son si̇nn Har·ald, bróður Ing·vars.\endsub\epa

\bpb Toolie had this stone raised after her son Harold, brother of Ingwar.\epb\epg

\bvg\bva[][2]%
Þęiʀ \alst{f}óru dręngi⸗la \hld\ \alst{f}iarri at gulli &
auk \alst{au}star⸗la \hld\ \edtrans{\alst{a}rni gǫ́fu}{gave to the eagle}{\Bfootnote{They provided a feast for the eagle, namely with the carnage of slain foes; for eagles and ravens as eaters of corpses and drinkers of blood cf.~\textcite[118,203,207--208]{Meissner1921}.  Similar expressions naturally occur in numerous Scaldic poems from Iceland and Norway, and the lack of an object to \emph{gǫ́fu} ‘gave’ reveals that this expression must have been well known also in Sweden.}}, &
dóu \alst{s}unnar⸗la \hld\ ȧ \alst{S}ęrk-landi.\eva

\bvb They journeyed like drengs far for gold \\
and easterly gave to the eagle; \\
they died southerly in Serkland.\evb\evg

\sectionline\clearpage


\section{Sö 198 (Mervalla)}

\begin{flushright}%
\textbf{Dating:} C11th

\textbf{Meter:} \Fornyrdislag
\end{flushright}%

\subsection{Introduction}

TODO.

\subsection{Runes}

TODO: Runes%textrunic

\subsection{Transliteration}

\textbf{siriþ · lit · resa · stan · [þin](a) [·] (a)(t) · suen · sin · [b]unta · h[n] · uft · siklt · til · simk(a)(l)(a) · t(u)ru[m] · knari · um · tumisnis}

\subsection{Text}

\bvg\bva[]%
Sig·ríð\emph{r} lét rêsa stên þęnna at Svên, si̇nn bónda. &
H\emph{a}nn opt \alst{s}iglt \hld\ til \alst{S}êmgala, &
\alst{d}ýrum knarri \hld\ um \edtrans{\alst{D}ómis-nęs}{Domeness}{\Bfootnote{The old Germanic name for Cape Kolka in Latvia, corresponding to German \emph{Domesnes}.}}.\eva

\bvb Syerid had this stone raised after Swayn, her husband. \\
He had often sailed to the Semigallians \\
on his costly knarr around Domeness.\evb\evg

\sectionline\clearpage


\section{Sö 204 (Överselö church)}

\begin{flushright}%
\textbf{Dating:} C11th

\textbf{Meter:} \Fornyrdislag
\end{flushright}%

\subsection{Introduction}

The inscription is heavily damaged, but can be reconstructed on the basis of the verse in Sö 206.

\subsection{Runes}

TODO: Runes%textrunic

\subsection{Transliteration}

\textbf{: hier : skal : stenta : stei(n) ... ...kia(l)t : run(u)... ... ... (u)arfaitr : (i)(f)tiʀ : faþurs:broþur · sen ·}

\subsection{Text}

\bvg\bva[]%
Hiar skal \alst{st}anda \hld\ \alst{st}ęinn \emph{at Ingj·}ald, &
\alst{r}u̇n\emph{um \alst{r}oðinn} \hld\ \emph{\alst{r}ęisti} Var·fęitr &
eptir fǫðurs bróður si̇nn.\eva

\bpb TODO: trans\epb\epg

\sectionline\clearpage


\section{Sö 206/208}

\begin{flushright}%
\textbf{Dating:} C11th

\textbf{Meter:} \Fornyrdislag
\end{flushright}%

\subsection{Introduction}

Two stones raised by a mother and her daughter.  The plural \emph{stęina} on Sö 206 suggests that the inscription between Sö 206 and 208 is continuous.

The “-barn” named on Sö 208 may perhaps be identified with one of the two men mentioned on Sö 213.

\subsection{Runes}

TODO: Runes%textrunic

\subsection{Transliteration}

Sö 206: \textbf{: hir : skal : stenta : staena : þisiʀ : runum : ru-niʀ :·· raisti : k-{}-{}-auk : at syni : sina : auk : hielmlauk : at bryþr : sina ·} \\
Sö 208: \textbf{: kuþlauk · auk : hielmlaug : þaʀ : -{}-{}-{}-(u) · raistu : ... ...ern : auk : at : þorstain :: boanta : koþan : buki : i frayslutum : arfi : fu(l)... ... sum iʀ : sustursun inibrataʀ :}

\subsection{Text}

\bpg\bpa[Sö 206]%
Hér skal \alst{st}anda \hld\ \alst{st}ęina þessiʀ \\
\alst{r}u̇num \alst{r}u\emph{ð}nir \hld\ \alst{r}ęisti G\emph{uð·l}aug \\
at syni sína auk Hialm·laug at brǿðr sína.\epa

\bpb Here shall these stones stand \\
reddened with runes.  Guthleigh raised them \\
after her sons, and Helmleigh after her brothers.\epb\epg


\bpg\bpa[Sö 208]%
\startsub Guð·laug auk Hialm·laug þáʀ \dots\ ręistu \dots\ \emph{at} \dots\ -bięrn auk at Þȯr·stęin\endsub \\
\alst{b}óanda góðan, \hld\ \alst{b}yggi i̇ Fręys-lundum \\
arfi Fu\emph{lluga} \dots\ sum eʀ systur-son Ęnni·brattaʀ.\epa

\bpb Guthleigh and Helmleigh, those \dots\ raised \dots\ [after] -barn and after Thurstan, \\
a good farmer.  He lived in the Freaslunds, \\
the heir of Fullow \dots\ who is the sister’s son of Endbrant.\epb\epg

\sectionline\clearpage


\section{Sö 207}

\begin{flushright}%
\textbf{Dating:} C11th

\textbf{Meter:} \Fornyrdislag
\end{flushright}%

\subsection{Introduction}

TODO.

\subsection{Runes}

TODO: Runes%textrunic

\subsection{Transliteration}

\textbf{kuþr... ... (f)aþur sin · fur · hfila · hn · til · iklans · kuþ halbi · sil hns}

\subsection{Text}

\bpg\bpa[0]%
TODO: text\epa

\bpb TODO: trans\epb\epg

\sectionline\clearpage


\section{Sö 213}

\begin{flushright}%
\textbf{Dating:} C11th

\textbf{Meter:} \Fornyrdislag
\end{flushright}%

\subsection{Introduction}

TODO.

\subsection{Runes}

TODO: Runes%textrunic

\subsection{Transliteration}

\textbf{s\textasciicircum tain : hiuk : esbern : stintn : at : uitum : bat miþ : runum : raisti : kyla : at : gaiʀbern : boanta : sin :· auk · kofriþ : at : faþur : sin : han uaʀ : boanti : bestr i : kili : raþi : saʀ : kuni :}

\subsection{Text}

\bpg\bpa[0]%
TODO: text\epa

\bpb TODO: trans\epb\epg

\sectionline\clearpage


\section{Sö 320}

\begin{flushright}%
\textbf{Dating:} C11th

\textbf{Meter:} \Fornyrdislag\ or \Kviduhattr
\end{flushright}%

\subsection{Introduction}

TODO.

\subsection{Runes}

᛬ ᚴᛅᛁᛦᚢᛅᛏᚱ ᛬ ᛅᚢᚴ ᛬ ᛅᚾᚢᛏᚱ ᛬ ᛅᚢᚴ ᛬ ᚢᛏᛅᛘᚱ ᛬ ᚱᛁᛏᛅ ᛬ ᛋᛏᛅᛁᚾ ᛬ ᛅᛏ ᛬ ᛒᚤᚱᛋᛏᛅᛁᚾ ᛫ ᛒᚱᚢᚦᚢᚱ ᛬ ᛋᛁᚾ ᛬ ᛋᛅᛦ ᚢᛅᛦ ᛬ ᛅᚢᛋᛏᚱ ᛫ ᛘᛁᚦ ᛁᚴᚢᛅᚱᛁ ᛬ ᛏᚱᛁᚴ ᛬ ᛋᚾᛁᛅᛚᛅᚾ ᛬ ᛋᚢᚾ ᛬ ᛚᛁᚠᛅᚤᛅᛦ%textrunic

\subsection{Transliteration}

\textbf{: kaiʀuatr : auk : anutr : auk : utamr : rita : stain : at : byrst(a)in · bruþur : sin : saʀ uaʀ : austr · miþ ikuari : trik : snialan : sun : lifayaʀ ×}

\subsection{Text}

\bpg\bpa[0]%
Gęiʀ·hvatr auk Ǫnundr auk Ȯ·tamr rétta stęin at Býr·stęin, bróður si̇nn, \\
sáʀ vaʀ \alst{au}str \hld\ með \alst{I}ng·vari, \\
dręng \alst{s}niallan, \hld\ \alst{s}unn \emph{H}líf·ęyjaʀ.\epa

\bpb TODO: trans\epb\epg

\sectionline\clearpage


\section{Sö 335}

\begin{flushright}%
\textbf{Dating:} C11th

\textbf{Meter:} \Fornyrdislag
\end{flushright}%

\subsection{Introduction}

TODO.

\subsection{Runes}

TODO: Runes%textrunic

\subsection{Transliteration}

\textbf{uulf\textasciicircum ui : raisti : stain þana| |at bruþur sin : u:snikin saʀ furs : a:ust:arla : maþ : i:ikn:u:ari : ksibari hulmstains}

\subsection{Text}

\bpg\bpa[0]%
TODO: text\epa

\bpb TODO: trans\epb\epg

\sectionline\clearpage


\section{Sö 338}

\begin{flushright}%
\textbf{Dating:} C11th

\textbf{Meter:} \Fornyrdislag
\end{flushright}%

\subsection{Introduction}

TODO.

\subsection{Runes}

TODO: Runes%textrunic

\subsection{Transliteration}

§ A \textbf{· ketil : auk + biorn + þaiʀ + raistu + stain + þin[a] + at + þourstain : faþur + sin + anuntr + at + bruþur + sin + auk : hu[skar]laʀ + hifiʀ + iafna + ketilau at + buanta sin · ¶ bruþr uaʀu þaʀ bistra mana : a : lanti auk : i liþi : uti : h(i)(l)(t)u sini huska(r)la : ui- +} \\
§ B \textbf{han + fial + i + urustu + austr + i + garþum + lis + furugi + lanmana + bestr}

\subsection{Text}

\bpg\bpa[0]%
\startsub Kętill auk Biorn, þęiʀ ręistu stęin þęnna at Þȯr·stęin, fǫður si̇nn, Ǫnundr at bróður si̇nn auk hús-karlaʀ *ęp\emph{t}iʀ iafna, Kętil·ęy at bóanda si̇nn.\endsub \\
\alst{B}rǿðr vǫ́ʀu þęiʀ \hld\ \alst{b}ętstra manna \\
\ind ȧ \alst{l}andi auk i̇́ \alst{l}iði úti; \\
\alst{h}eldu si̇nn\emph{a} \hld\ \alst{h}ús-karla ve\emph{l}. \\
Hann fiall i̇́ \alst{o}rrustu \hld\ austr i̇́ Gǫrðum, \\
\alst{l}iðs forungi, \hld\ \alst{l}and-manna bętstr. \epa

\bpb TODO: trans. \\
They were brothers of the best men, \\
\ind on land and out in the troop.
He fell in battle to the east in the Yards, \\
the steward of the troop, the best of lended men.\epb\epg

\sectionline\clearpage


\section{Sö 340}

\begin{flushright}%
\textbf{Dating:} C11th

\textbf{Meter:} \Fornyrdislag
\end{flushright}%

\subsection{Introduction}

TODO.

\subsection{Runes}

TODO: Runes%textrunic

\subsection{Transliteration}

\textbf{[: stin : saʀ:si sąnr : at · frusta ik]irunaʀ : [sun : aik : kuþlifs : ulifʀ : raisti aftiʀ bruþur sin ¶ kus]}

\subsection{Text}

\bpg\bpa[0]%
TODO: text\epa

\bpb TODO: trans\epb\epg

\sectionline\clearpage


\section{Sö 345}

\begin{flushright}%
\textbf{Dating:} C11th

\textbf{Meter:} \Fornyrdislag
\end{flushright}%

\subsection{Introduction}

TODO.

\subsection{Runes}

TODO: Runes%textrunic

\subsection{Transliteration}

\textbf{... ...in × þinsa × at × kai(ʀ)... ... ...-n · eʀ · e[n-a]þr × ut - × kr...}

\subsection{Text}

\bpg\bpa[0]%
TODO: text\epa

\bpb TODO: trans\epb\epg

\sectionline\clearpage


\section{Sö Fv1948;289 (Aspa)}

\begin{flushright}%
\textbf{Dating:} C11th

\textbf{Meter:} \Fornyrdislag
\end{flushright}%

\subsection{Introduction}

TODO: INTRODUCTION.

The Rainold on this stone has been identified with the man of the same name mentioned on Sö Fv1948;295, although the conncetion is uncertain.

\subsection{Runes}

TODO: Runes%textrunic

\subsection{Transliteration}

\textbf{ąstriþ : lit : -ira : ku(m)... ...usi ÷ at : anunt ÷ auk : raknualt : sun : sin ÷: urþu : ta...ʀ : - (t)an...-...(k)u : ua-u : rikiʀ : ą rauniki : ak : snialastiʀ : i : suiþiuþu}

\subsection{Text}

\bpg\bpa[0]%
\startsub Ǫ̇st·ríðr lét \emph{g}ęra kum\emph{bl þa}usi at Ǫn·und ok Ragn·vald, sun si̇nn.\endsub \\
Urðu \alst{d}\emph{auði}ʀ \hld\ i̇́ \alst{D}an-\emph{mǫrk}u, \\
vǫ́\emph{ʀ}u \alst{r}íkiʀ \hld\ ȧ \alst{R}auningi \\
auk \alst{s}niallastiʀ \hld\ i̇́ \alst{S}vé·þiuðu. \epa

\bpb Estfrid had these monuments made after Anund and Rainold, her son. \\
They died in Denmark, \\
were powerful in Reening, \\
and the most valiant in Sweden.\epb\epg

\sectionline\clearpage


\section{Sö Fv1954;22}

\begin{flushright}%
\textbf{Dating:} C11th

\textbf{Meter:} \Fornyrdislag
\end{flushright}%

\subsection{Introduction}

TODO.

\subsection{Runes}

TODO: Runes%textrunic

\subsection{Transliteration}

\textbf{...i : risti : -{}-{}-... ... ...in... ... sin : han : iʀ : entaþr : i : austruiki : ut : ą : la-...}

\subsection{Text}

\bpg\bpa[0]%
TODO: text\epa

\bpb TODO: trans\epb\epg

\sectionline\clearpage


\section{U 11}

\begin{flushright}%
\textbf{Dating:} C11th

\textbf{Meter:} \Fornyrdislag
\end{flushright}%

\subsection{Introduction}

TODO.

The stone clearly bears a connection to U 16, which also mentions King Hathkin’s “rowing” (\emph{róð}).
%TODO: investigate connection to Sö 101, Sö 106, and U 617.

\subsection{Runes}

TODO: Runes%textrunic

\subsection{Transliteration}

\textbf{raþ| |þu : runaʀ : ret : lit : rista : toliʀ : bry[t]i : i roþ : kunuki : toliʀ : a(u)k : gyla : litu : ris... ...- : þaun : hion : eftiʀ ...k : merki ¶ srni... ¶ haku(n) · (b)aþ : rista}

\subsection{Text}

\bpg\bpa[0]%
\alst{R}áð þú \alst{r}u̇naʀ \hld\ \alst{r}étt lét \alst{r}ísta \\
Toliʀ bryti i̇ róð konungi. \\
Toliʀ auk Gylla létu rísta \dots\ þaun hjón eptiʀ \dots\ męrki \textbf{†srni†} \dots \\
Hǫ́·kun bað rísta.\epa

\bpb TODO: trans\epb\epg

\sectionline\clearpage


\section{U 69}

\begin{flushright}%
\textbf{Dating:} C11th

\textbf{Meter:} \Fornyrdislag
\end{flushright}%

\subsection{Introduction}

TODO.

\subsection{Runes}

TODO: Runes%textrunic

\subsection{Transliteration}

\textbf{raknilfʀ × lit × kirua × bru × þasi × iftiʀ × anunt + sun + s[i](n) [× k]uþan × kuþ -{}-{}-bi × ąns × ant × uk × salu × bitr × þan × ąn krþi × til × munu × iki × mirki × miʀi × uirþa × muþiʀ × karþi × i(f)tiʀ × sun × sin × ainika ×}

\subsection{Text}

\bpg\bpa[0]%
\startsub Ragn·ęlfʀ lét gęrva bró þási ęptiʀ Ǫnund, sun si̇nn góðan. Guð hialpi hans ǫnd ok sǫ́lu bętr þan hann gęrði til.\endsub \\
\alst{M}unu ęigi \alst{m}ęrki \hld\ \alst{m}ęiʀi verða. \\
Móðiʀ gǫrði ęptiʀ sun si̇nn ęiniga.\epa

\bpb TODO: trans\epb\epg

\sectionline\clearpage


\section{U 102}

\begin{flushright}%
\textbf{Dating:} C11th

\textbf{Meter:} \Fornyrdislag
\end{flushright}%

\subsection{Introduction}

The inscription is very hard to read.

TODO.

\subsection{Runes}

TODO: Runes%textrunic

\subsection{Transliteration}

\textbf{kal lit hkua heli þisa iftiʀ · suni sina toa · auk þau ikiþora · bʀ- (a)ina kiarþu mukit mirki furiʀ arkum mani}

\subsection{Text}

\bvg\bva[]%
Káll lét \alst{h}\emph{ǫ}ggva \hld\ \alst{h}ęlli þęssa &
ęptiʀ \alst{s}yni \hld\ \alst{s}i̇́na tvȧ, &
auk þau \alst{I}ngi·þȯra b\emph{ró} \hld\ \alst{ęi}na giǫrðu, &
\alst{m}ykit \alst{m}ęrki \hld\ fyriʀ \emph{\alst{m}}ǫrgum manni.\eva

\bvb TODO: trans\evb\evg

\sectionline\clearpage


\section{U 112}

\begin{flushright}%
\textbf{Dating:} C11th

\textbf{Meter:} \Fornyrdislag
\end{flushright}%

\subsection{Introduction}

TODO.

%The inscription has been connected to a large number of runestones due to the sharing of names: U 100, 160, 161, 225, 226, 328, 336, 343, 344, 364, Fv1968;279b.

\subsection{Runes}

TODO: Runes%textrunic

\subsection{Transliteration}

§ A \textbf{· rahnualtr · lit · rista · runar · efʀ · fastui · moþur · sina · onems · totʀ · to i · aiþi · kuþ hialbi · ant · hena ·} \\
§ B \textbf{runa · rista · lit · rahnualtr · huar a × griklanti · uas · lis · forunki ·}

\subsection{Text}

\bpg\bpa[0]%
TODO: text\epa

\bpb TODO: trans\epb\epg

\sectionline\clearpage


\section{U 114}

\begin{flushright}%
\textbf{Dating:} C11th

\textbf{Meter:} \Fornyrdislag
\end{flushright}%

\subsection{Introduction}

TODO.

\subsection{Runes}

TODO: Runes%textrunic

\subsection{Transliteration}

§ A \textbf{ikriþ ’ l[i]t · laþbo · kiara ’ auk · stain · haku[a eftir] ikim[a]r (b)(o)[ta s]in · auk · eftar · tan · auk · eftir · baka · suni · sina} \\
§ B \textbf{’ þaiʀ byku ’ i rynby ’ auk ’ bo atu ’ [kr]istr · ialbi ’ s(a)(l)(u) · þai[r- ·] (þ)it skal ’ at minum · mana ’ miþan · min li[fa]}

\subsection{Text}

\bpg\bpa[0]%
TODO: text\epa

\bpb TODO: trans\epb\epg

\sectionline\clearpage


\section{U 127/164/165/261}

\begin{flushright}%
\textbf{Dating:} C11th

\textbf{Meter:} \Fornyrdislag
\end{flushright}%

\subsection{Introduction}

These four stones all share a near-identical inscription and visual layout.  All four originally belong to the so-called Earlbanke’s bridge, a stretch of road at which U 164--5 are still standing.  U 127 and 261 have been moved to different locations.

Since the wording of the inscriptions are near-identical, I thought it simplest to include them all under the same section and note the differences between them (excepting minor divergences in punctuation) by means of the critical apparatus.

Like many other one-line verses, the poetic nature of these inscriptions is doubtful, especially since the transposition of \emph{allan} in U 127, which is not inferior to the others in terms of execution, renders the line deficient and makes one question whether Earlbanke really intended a poetic line.

Earlbanke belonged to an elite family which raised a number of other stones in southern Attundaland, north of modern Stockholm.  Other stones securely connected to Earlbanke’s family are U 101, 135--137, 140, 142--143, 147--150, and 212.

\subsection{Runes}

TODO: Runes%textrunic

\subsection{Transliteration}

U 127: \textbf{× iarla×baki × lit raisa [×] staina × þasa at sik × kuikuan ¶ × auk bru þisa karþi × fur ąnt sina × auk × ain ati tabu ala-}
U 164: \textbf{× iarlabaki × lit × raisa × stain × þisa × at sik × kuikuan × ¶ × auk bru × þisa × karþi × fur ąnt × sina × auk ain ati × alan × tabu × kuþ hialbi ąnt hans}
U 165: \textbf{ia[r]labaki × li- -{}-{}-sa × staina × þisa × at sik × kuikuan × auk bru þis(a) [× karþi fur ąnt ×] si-{}- -{}-{}- × ati + (a)lan × tabu +}
U 261: \textbf{iarlaba... lit × raisa × s[taina × þasa × a]t sik × kuikuan × ¶ auk × bru × þisa × karþi × fur ąnt [×] s[in]a × auk × ain ati alan tabu}

\subsection{Text}

\bpg\bpa[0]%
\startsub \edtext{Iarla·banki}{\Afootnote{\emph{\textbf{iarlaba...}} \textup{U 261}}} \edtext{lét ręisa}{\Afootnote{\emph{\textbf{li- -{}-{}-sa}} \textup{U 165}}} \edtext{stęina}{\Afootnote{\emph{\textbf{stain}} \textup{U 164}}} \edtext{þása}{\Afootnote{\emph{þęssa (‘\textbf{þisa}’)} \textup{U 164/165}}} at sik kvikvan, auk bró þęssa gęrði fyr ǫnd \edtext{si̇́na,\endsub \\
auk \alst{ęi}nn}{\lemma{si̇́na, auk ęinn}\Afootnote{\emph{\textbf{si-{}- -{}-{}-}} \textup{U 165}}} \alst{á}tti \hld\ \edtext{\alst{a}llan \edtrans{Tȧ·bý}{Toughby}{\Bfootnote{The ancestor of \emph{Täby}, the name of the place in modern Swedish.
Naturally, Wiking Age Toughby was merely a farmstead or village, and not the large municipality it is today.}}}{\Afootnote{rev. \emph{(‘\textbf{tabu ala-}’)} \textup{U 127}}}. \\
\edtext{Guð hialpi ǫnd hans.}{\Afootnote{om. \textup{U 164/165/261}}}\epa

\bpb Earlbanke had these stones raised after himself while alive and made this bridge for his spirit, \\
and he alone owned all Toughby. \\
God help his spirit.\epb\epg

\sectionline\clearpage


\section{U 136}

\begin{flushright}%
\textbf{Dating:} C11th

\textbf{Meter:} \Fornyrdislag
\end{flushright}%

\subsection{Introduction}

TODO.

For the family responsible for raising the stone, see the Introduction to U 127.

\subsection{Runes}

TODO: Runes%textrunic

\subsection{Transliteration}

\textbf{× astriþr × la(t) + raisa × staina × þasa × [a]t austain × buta sin × is × suti × iursalir auk antaþis ub i × kirkum}

\subsection{Text}

\bpg\bpa[0]%
TODO: text\epa

\bpb TODO: trans\epb\epg

\sectionline\clearpage


\section{U 153/155}

\begin{flushright}%
\textbf{Dating:} C11th

\textbf{Meter:} \Fornyrdislag
\end{flushright}%

\subsection{Introduction}

TODO.

These two stones are united by textual evidence and geography; they may thus be taken to form a continual inscription, even if only U 153 is actually metrical.

\subsection{Runes}

TODO: Runes%textrunic

\subsection{Transliteration}

U 153: \textbf{...[(u)](a)i- × [(a)]uk × ulf- litu × raisa × stai-(a) × e(f)tiʀ × hlftan · auk · eftiʀ × kunar × bryþr × sina × þaiʀ · antaþus × aust... ...(u)m} \\
U 155: \textbf{...n × auk × ulfr litu × rai(s)(a) sta... ... ... ...na × þaiʀ × uaru × suniʀ × arnar × auk × raknfriþar ×}

\subsection{Text}

\bpg\bpa[0]%
TODO: text\epa

\bpb TODO: trans\epb\epg

\sectionline\clearpage


\section{U 166}

\begin{flushright}%
\textbf{Dating:} C11th

\textbf{Meter:} \Fornyrdislag
\end{flushright}%

\subsection{Introduction}

TODO.

\subsection{Runes}

TODO: Runes%textrunic

\subsection{Transliteration}

\textbf{ku lit · rasa · kuþlauh · iftiʀ · sun · miuk · nutan · trik · mirki · þisa · an it · sihuastr ·}

\subsection{Text}

\bpg\bpa[0]%
TODO: text\epa

\bpb TODO: trans\epb\epg

\sectionline\clearpage


\section{U 207--8}

\begin{flushright}%
\textbf{Dating:} C11th

\textbf{Meter:} \Fornyrdislag
\end{flushright}%

\subsection{Introduction}

TODO.  The inscription is syntactically continuous between the two stones.

\subsection{Runes}

TODO: Runes%textrunic

\subsection{Transliteration}

U 207: \textbf{ulfr ’ auk ’ þurmontr ’ auk ’ kamal ’ lata ’ reis(a) þisa ’ stina ’ þar + eftʀ ’ faþur ’ sin ’}
U 208: \textbf{ristu ’ merk(i) at ’ man ’ metan ’ sunir ’ alkoþir ’ at ÷ sin ’ faþur sterkar usiti ’ rsti stina}

\subsection{Text}

\bpg\bpa[0]%
TODO: text\epa

\bpb TODO: trans\epb\epg

\sectionline\clearpage


\section{U 214--15}

\begin{flushright}%
\textbf{Dating:} C11th

\textbf{Meter:} \Fornyrdislag
\end{flushright}%

\subsection{Introduction}

Two paired stones with a single continuous inscription.  At present both are located in Vallentuna church.  U 215 should be read before 214 in the series.  The Ingbury who raised the stone is most likely the same Ingbury who raised the non-metrical U 146 for her husband and son.  U 146 is also located in Vallentuna (though not in the church) and gives the name of the father as Holmstan (\emph{Holm·stęinn}) and the son as Thurstan (\emph{Þȯr·stęinn}).

\subsection{Runes}

TODO: Runes%textrunic

\subsection{Transliteration}

U 214: \textbf{uk × inkiber × eftiʀ × buanta × sin ’ han ’ troknaþi ÷ a ’ holms ’ hafi ’ skreþ ’ knar ’ hans ’ i ’ kaf þriʀ ’ eniʀ ’ kamo ’ af} \\
U 215: \textbf{[þaniltr · uk · olfil](r) · [li](t)(u) s[t]ain × e(f)[tiʀ faþor · uk broþor sin ·]}

\subsection{Text}

\bpg\bpa[0]%
\startsub \emph{R}agn·hildr ok Ulf·hildr létu stęin ęptiʀ fǫður ok bróður si̇nn, ok Ingi·berg ęptiʀ bóanda si̇nn.\endsub \\
\alst{H}ann drunknaði \hld\ ȧ \alst{H}olms-hafi; \\
skrêð \alst{k}narr hans i̇́ \alst{k}af; \\
þríʀ \alst{ê}niʀ kǫ́mu \alst{a}f.\epa

\bpb TODO: trans\epb\epg

\sectionline\clearpage


\section{U 225--26 (Bällsta)}

\begin{flushright}%
\textbf{Dating:} C11th

\textbf{Location:} Arkils tingstad, Bällsta, Vallentuna

\textbf{Meter:} \Fornyrdislag
\end{flushright}%

\subsection{Introduction}

U 225--6 are two paired runestones still standing in their original location, namely at the thingstead mentioned in U 225.  In the 18th century the two stones were connected by a line of smaller, uninscribed stones between them, but these have since been lost.  The inscription is continuous and consists of 7 alliterative long-lines.

The family responsible for U 225--6 was prolific, and also raised U 100, 160--161, 328, 336, 343--344, and 364.

\subsection{Runes}

TODO: Runes%textrunic

\subsection{Transliteration}

U 225: \textbf{... uk · arkil · uk · kui · þiʀ · kariþu · iar · þikstaþ ... ...unu · iki mirki · maiʀi · uirþa · þan · ulfs · suniʀ · iftiʀ · kir... ...iʀ · suinaʀ · at · sin · faþur} \\
U 226: \textbf{ristu · stina · uk · staf · uan · uk · in · mikla · at · iartiknum uk kuriþi · kas at · uiri · þu mon i krati · kiatit lata ¶ kunar ik stin}

\subsection{Text}

\bpg\bpa[U 225]%
\startsub\dots\ ok Arn·kęll ok Gýi þęiʀ gęrðu ’iar þing-stað \dots\endsub \\
\emph{\alst{M}}unu ęigi \alst{m}ęrki \hld\ \alst{m}ęiʀi verða, \\
þan \alst{U}lfs syniʀ \hld\ ęptiʀ gęr\emph{ðu,} \\
\emph{\alst{s}niall}iʀ \alst{s}vęinaʀ \hld\ at \alst{s}i̇nn fǫður,\epa

\bpb\dots\ and Arnkettle and Gew they made the thingstead here, \dots \\
No greater landmarks will there be \\
than Wolf’s sons afterwards made, \\
valiant swains, for their father.\epb\epg

\bvg\bva[U 226]%
ręistu \alst{st}ęina \hld\ ok \edtrans{\alst{st}af \edtext{o\emph{f}an}{\Afootnote{\emph{\textbf{uan} U 226}; \emph{unnu} most other eds.}}}{the staff above}{\Bfootnote{Unclear; perhaps a tall wooden staff or column?}} &
ok \edtrans{\alst{i}nn mikla}{the great \dots}{\Bfootnote{The word \emph{inn} may either be read as the article \emph{inn, hinn} ‘the’ or as a f. noun of unclear sense.  The latter is preferable since the weakly stressed article cannot uphold alliteration, and since \emph{hinn mikla} ‘the great one’ is meaningless.}} \hld\ at \alst{ia}r-tęiknum, &
ok \alst{G}ý·ríði \hld\ \alst{g}ats at veri; &
þý man \edtrans{i̇́ \alst{g}ráti}{in weeping}{\Bfootnote{Or ‘in a dirge’, if \emph{grátr} refers uniquely to a lost genre of poems of lamentation, in which case the line probably refers to a longer oral poem serving as a complement to the necessarily quite limited runic poem.}} \hld\ \alst{g}iatit láta. &
Gunnarr ’iogg stęin.\eva

\bvb They raised the stones and the staff above, \\
and the great \dots\ as tokens of remembrance; \\
also Guefrid cared for her husband. \\
Thus will they have it spoken of in weeping. \\
Guther hewed the stone.\evb\evg

\sectionline\clearpage


\section{U 300 (Skånela)}

\begin{flushright}%
\textbf{Dating:} C11th

\textbf{Meter:} \Fornyrdislag
\end{flushright}%

\subsection{Introduction}

TODO.

\subsection{Runes}

TODO: Runes%textrunic

\subsection{Transliteration}

\textbf{[þyrui + lit + raisa + stain × eftiʀ + hlftain + merki + mukit + eftiʀ × man kuþan]}

\subsection{Text}

\bpg\bpa[0]%
Þýr·ví lét ręisa stęin ęptiʀ Halfda*n, \\
\alst{m}ęrki \alst{m}ykit \hld\ ęptiʀ \alst{m}ann góðan.\epa

\bpb TODO: trans\epb\epg

\sectionline\clearpage


\section{U 323 (Sälna)}

\begin{flushright}%
\textbf{Dating:} C11th

\textbf{Meter:} \Fornyrdislag
\end{flushright}%

\subsection{Introduction}

An exceptionally tall standing stone.  No decorations, but explicitly Christian.  Note the sporadic loss of initial \emph{h} and use of stung runes ᚤ, ᚵ, ᛂ \textbf{y, g, e}.

\subsection{Runes}

× ᛁᚤᛋᛏᛁᚾ × ᛅᚢᚴ × ᛁᚢᚱᚢᚾᛏᚱ × ᛅᚢᚴ × ᛒᛁᚢᚱᚾ × ᚦᛁᛦ [× ᛒᚤᚱᚤᚦᚱ × ᚱᛁᛋᚦᚢ] \dots\ \dots ᛋᛏᛁᚾ × ᛏᚱᚢᛘᛋ × ᚠ[ᛅᚦᚢᚱ] × ᛋᛁᚾ × ᚴᚢᚦ × ᛁᚼᛚᛒᛁ × ᚭᚾᛋ × ᚭᚾᛏ × ᛅᚢᚴ × ᛋᛂᛚᚢ ᛭ ᚠᚢᚱ ᛭ ᚴᛁᚠᛁ × ᚭᚾᚢᛘ × ᛋᛆᚴᛅᛦ × ᛅᚢᚴ ᛋᚢᛏᛁᛦ ×

× ᚼᛁ × ᛘᚢᚾ × ᛚᛁᚵᛁᛅ × ᛘᛁᚦ ᛭ ᛆᛚᛏᚱ ᛭ ᛚᛁᚠᛁᛦ × ᛒᚱᚢ × ᚼᚱᚦ × ᛋᛚᛅᚵᛁᚾ × ᛒᚱᛁᚦ × ᛂ[ᚠᛏ ᚴ]ᚢᚦ-{}- ᛋᚢᛂᚾᛅᛦ ᚴᛅ[ᚱᚦᚢ ×] ᛅᛏ × ᛋᛁᚾ × ᚠᛅᚦᚢᚱ × ᛘᚭ × ᛁᚵᛁ × ᛒᚱᚢᛏᛅᛦ × ᚴᚢᛘᛚ × ᛒᛂᛏᚱᛅ × ᚢᛂᚱᚦᛅ ᛭%textrunic

\subsection{Transliteration}

\textbf{× iystin × auk × iuruntr × auk × biurn × þiʀ [× byryþr × risþu] ... ...stin × trums × f[aþur] × sin × kuþ × ihlbi × ąns × ąnt × auk × selu + fur+kifi × ąnum × sakaʀ × auk × sutiʀ × ¶ \\
× hi × mun × ligia × meþ + altr + lifiʀ × bru × hrþ×slagin × briþ × e[ft × k]uþ-{}- suenaʀ k[arþu ×] at × sin × faþur × mą × igi × brutaʀ×kuml × betra × uerþa +}

\subsection{Text}

\bpg\bpa[0]%
\startsub Ęy·stên auk Iurundr auk Biorn, þêʀ brǿðr rêsþu \dots\ \emph{ęptir} \dots stên \textbf{trums} fǫður si̇nn.\endsub \\
\startsub Guð ’ialpi ’ans ǫnd auk sêlu, for·gifi ’ǫ̇num sakaʀ auk syndiʀ.\endsub\epa

\bpb Eanstan and Earwind and Barn---those brothers raised \dots\ [after] -stan \dots, their father. \\
God help his breath and soul, forgive him his crimes and sins.\epb\epg


\bvg\bva[][3]%
\emph{\alst{Ę}}i mun liggja \hld\ með \alst{a}ldr lifiʀ &
\alst{b}ró harð-slagin \hld\ \alst{b}rêð e\emph{ft g}óð\emph{an}; &
\alst{s}vênaʀ k\emph{ęrðu} \hld\ at \alst{s}i̇nn fǫður &
má êgi \alst{b}rø̂taʀ-kuml \hld\ \alst{b}ętra verða.\eva

\bvb For ever will lie while the age lives \\
the hard-hammered bridge, broad, after a good man. \\
The swains made it after their father; \\
no roadside monument can be greater.\evb\evg

\sectionline\clearpage


\section{U 336}

\begin{flushright}%
\textbf{Dating:} C11th

\textbf{Meter:} \Fornyrdislag
\end{flushright}%

\subsection{Introduction}

TODO.

\subsection{Runes}

TODO: Runes%textrunic

\subsection{Transliteration}

\textbf{[ul]fʀ × lit × risa stin × þi[n]a × iftiʀ × unim × faþurs×bruþr sin þiʀ × buku × baþiʀ × i × baristam ’}

\subsection{Text}

\bpg\bpa[0]%
\startsub Ulfʀ lét rêsa stên þęnna ęptiʀ Ȯ·nę́m, fǫðurs-bróð\emph{u}r si̇nn.\endsub \\
Þêʀ \alst{b}yggu \alst{b}áðir \hld\ i̇́ \alst{B}ári-stǫm.\epa

\bpb TODO: trans\epb\epg

\sectionline\clearpage


\section{U 349}

\begin{flushright}%
\textbf{Dating:} C11th

\textbf{Meter:} \Fornyrdislag
\end{flushright}%

\subsection{Introduction}

The stone is lost but survives as a well executed drawing.

\subsection{Runes}

TODO: Runes%textrunic

\subsection{Transliteration}

\textbf{[uikitil × uk × usur -...u × risa × stin × þina · iftiʀ × ustin × faþ... ... ...þan × ąn furs × uti × miþ × ala × skibin × kuþ × ialbi × (a)t]}

\subsection{Text}

\bpg\bpa[0]%
Ví·kętill ok Ǫssurr \dots \emph{lét}u rêsa stên þenna eptir Ø̂·stên, fǫð\emph{ur si̇nn gó}ðan. ’Ann fórs úti með alla skipan. \\
Guð ’ialpi ǫnd.\epa

\bpb TODO: trans\epb\epg

\sectionline\clearpage


\section{U 358}

\begin{flushright}%
\textbf{Dating:} C11th

\textbf{Meter:} \Fornyrdislag
\end{flushright}%

\subsection{Introduction}

It has been suggested that U 358 was originally paired with the fragmentary U Fv1979;244b (TODO: cite).

Note the loss of initial \emph{h-}.

\subsection{Runes}

TODO: Runes%textrunic

\subsection{Transliteration}

\textbf{fulkmar × lit × risa × stin × þina × iftiʀ × fulkbiarn × sun × sin × saʀ × itaþis × uk miþ krkum × kuþ × ialbi × ans × ąt uk salu}

\subsection{Text}

\bpg\bpa[0]%
\startsub Folk·marr lét rêsa stên þęnna ęptiʀ Folk·biǫrn, sun si̇nn.\endsub \\
Sáʀ \alst{ę}ndaðis \hld\ \alst{o}k með Gr\emph{i}kkum. \\
Guð ’ialpi ’ans ǫnd ok sǫ́lu.\epa

\bpb TODO: trans\epb\epg

\sectionline\clearpage


\section{U 437 (Broby)}

\begin{flushright}%
\textbf{Dating:} C11th

\textbf{Meter:} \Fornyrdislag
\end{flushright}%

\subsection{Introduction}

TODO.

\subsection{Runes}

TODO: Runes%textrunic

\subsection{Transliteration}

\textbf{kuna : a[uk su]tari : ristu : s(t)[... ... sikbiar- sati : iftiʀ : sin faþur : stin : almykin : i : s...]}

\subsection{Text}

\bpg\bpa[0]%
TODO: text\epa

\bpb TODO: trans\epb\epg

\sectionline\clearpage


\section{U 512}

\begin{flushright}%
\textbf{Dating:} C11th

\textbf{Meter:} \Fornyrdislag
\end{flushright}%

\subsection{Introduction}

TODO.

\subsection{Runes}

TODO: Runes%textrunic

\subsection{Transliteration}

§ A \textbf{fryb... ... ...i × huk f(r)... ...(f)-þr × þiʀ ristu × sti- -...(n)a × iftiʀ × kuþ-ar × ...(a)þur × sin kuþan × huk tuma + (b)...-þur × sin ×× kuþ ...i ąt × þiʀa × huk × salu × hu... ... ...uþiʀ × bitr × þan × þiʀ × -ar- × h...} \\
§ B \textbf{...ir × skal × stan\textasciicircum ta × stin × uiþ- bryku × suniʀ × at × faþuʀ satu · kuþan · kaiʀ -... ...} \\
§ C \textbf{-{}-(u)-aʀ × mirki × at × bua-... ...}

\subsection{Text}

\bpg\bpa[0]%
TODO: text\epa

\bpb TODO: trans\epb\epg

\sectionline\clearpage


\section{U 605}

\begin{flushright}%
\textbf{Dating:} C11th

\textbf{Meter:} \Fornyrdislag
\end{flushright}%

\subsection{Introduction}

A lost runestone.

\subsection{Runes}

TODO: Runes%textrunic

\subsection{Transliteration}

\textbf{[· iskirun · harþiʀ · totiʀ · lit · risti · runiʀ · ati · sik · sialfan · hn · uil · austr · fara · auk · ut · til · iursala · fair · risti · runiʀ ·]}

\subsection{Text}

\bpg\bpa[0]%
\startsub I\emph{n}gi·ru̇n, Harð\emph{a}ʀ dóttiʀ, lét ríst\emph{a} ru̇naʀ at* sik sialfa*.\endsub \\
Hȯn vill \alst{au}str fara \hld\ auk \alst{ú}t til Iór-sala. \\
F\emph{ót}r rísti ru̇niʀ.\epa

\bpb Ingrun, Hareth’s daughter, had these runes carved after herself. \\
She wants to go east and out to Jerusalem. \\
Foot carved the runes.\epb\epg

\sectionline\clearpage


\section{U 614}

\begin{flushright}%
\textbf{Dating:} C11th

\textbf{Meter:} \Fornyrdislag
\end{flushright}%

\subsection{Introduction}

TODO.

\subsection{Runes}

TODO: Runes%textrunic

\subsection{Transliteration}

\textbf{· skuli · auk · folki · lata · reisa · þinsa · stein · iftʀ · broþur · sin · husbiorn · hn us| |siok · uti · þa þiʀ · kialt · toku · a kutlanti ·}

\subsection{Text}

\bpg\bpa[0]%
TODO: text\epa

\bpb TODO: trans\epb\epg

\sectionline\clearpage


\section{U 617 (Bro)}

\begin{flushright}%
\textbf{Dating:} C11th

\textbf{Meter:} \Fornyrdislag
\end{flushright}%

\subsection{Introduction}

TODO.

A beautiful rectangular stone.  The family involved in raising this stone were quite prolific; see Sö 106 above.

\subsection{Runes}

TODO: Runes%textrunic

\subsection{Transliteration}

\textbf{kinluk × hulmkis × tutiʀ × systiʀ × sukruþaʀ × auk × þaiʀa × kaus × aun × lit × keara × bru × þesi × auk × raisa × stain × þina × eftiʀ × asur × bunta · sin × sun × hakunaʀ × iarls × saʀ × uaʀ × uikika × uaurþr × miþ × kaeti × kuþ × ialbi × ans × nu × aut × uk × salu}

\subsection{Text}

\bpg\bpa[0]%
\startsub Ginn·laug, Holm·gęiʀs dóttiʀ, systiʀ Syg·røðaʀ auk þęiʀa Gau\emph{t}s, hǫ̇n let giǫra bró þęssi auk ręisa stęin þęnna ęptiʀ Ǫssur, bȯnda si̇nn, sun \edtrans{Hǫ́·konaʀ iarls}{Earl Hathkin}{\Bfootnote{Unclear, but he is probably not the Norwegian Earl of Lade Hathkin Siwardson (c. 937--995).}}.\endsub \\
Sáʀ vaʀ \alst{v}íkinga \hld\ \alst{v}ǫrðr með Gęiti. \\
Guð hialpi hans nú ǫnd ok sǫ́lu.\epa

\bpb Yinlie, Holmgar’s daughter, the sister of Syeferth and Geat---she had this bridge made and this stone raised after Answer, her husband, the son of Earl Hathkin. \\
He was the watchman against Wikings with Goater. \\
God now help his spirit and soul.\epb\epg

\sectionline\clearpage


\section{U 644}

\begin{flushright}%
\textbf{Dating:} C11th

\textbf{Meter:} \Fornyrdislag
\end{flushright}%

\subsection{Introduction}

The stone is clearly paired with U 654, but their inscriptions are distinct enough to be treated separately.  The Andwit mentioned on these stones is likely the same man as mentioned in the non-metrical U 643.

\subsection{Runes}

TODO: Runes%textrunic

\subsection{Transliteration}

\textbf{an(u)(i)(t)r : auk · kiti : auk · kar : auk · blisi · auk · tiarfr · þir · raistu · stain þina · aftiʀ · kunlaif · fąþur : sin han : fil · austr : miþ : ikuari ¶ kuþ heabi ąntini}

\subsection{Text}

\bpg\bpa[0]%
\startsub An\emph{d·v}éttr auk Giti auk Kǫ́rr auk Blesi auk Diarfr, þęir ręistu stęin þenna ęptiʀ Gunn·lęif, fǫður si̇nn,\endsub \\
hann fell \alst{au}str \hld\ með \alst{I}ng·vari. \\
Guð hia\emph{l}pi ǫnd’inni.\epa

\bpb TODO: trans\epb\epg

\sectionline\clearpage


\section{U 654}

\begin{flushright}%
\textbf{Dating:} C11th

\textbf{Meter:} \Fornyrdislag
\end{flushright}%

\subsection{Introduction}

TODO.

\subsection{Runes}

TODO: Runes%textrunic

\subsection{Transliteration}

\textbf{+ a-{}-itr : auk · ka(r) auk : kiti : auk : -[l]isi : auk · tiarfr : ris[t]u : stain : þena : aftir : kunlaif : fąþur sin ¶ is u[a]s nus(t)(r) · m[i](þ) ikuari : tribin kuþ : hialbi : ą(t) þaira ¶ al-ikr| |raistik · runar ¶ is kuni + ual · knari stura}

\subsection{Text}

\bpg\bpa[0]%
\startsub An\emph{d·v}éttr auk Kǫ́rr auk Giti auk Blesi auk Diarfr ręistu stęin þęnna ęptir Gunn·lęif, fǫður si̇nn,\endsub \\
es vas \alst{\emph{a}u}str \hld\ með \alst{I}ng·vari drepinn. \\
Guð hialpi ǫnd þęira. \\
Al·\emph{ríkr} \alst{r}ęist’ik \alst{r}u̇naʀ \\
es \alst{k}unni vęl \hld\ \alst{k}nęrri stýra.\epa

\bpb TODO: trans\epb\epg

\sectionline\clearpage


\section{U 661}

\begin{flushright}%
\textbf{Dating:} C11th

\textbf{Meter:} \Fornyrdislag
\end{flushright}%

\subsection{Introduction}

TODO.

\subsection{Runes}

TODO: Runes%textrunic

\subsection{Transliteration}

\textbf{kairui × auk × kula × ristu × stain þina × aftir × ąnunt × fąþur sia is uas × austr × tauþr × miþ × ikuari × kuþ × hialbi ąt × ąnutar}

\subsection{Text}

\bpg\bpa[0]%
TODO: text\epa

\bpb TODO: trans\epb\epg

\sectionline\clearpage


\section{U 703 (Västra Väppeby)}

\begin{flushright}%
\textbf{Dating:} C11th

\textbf{Meter:} \Fornyrdislag
\end{flushright}%

\subsection{Introduction}

A standing stone inscribed on one side.  There is no cross present, but a large four-legged beast with a long tail.  The stone is heavily damaged but mostly readable.  The expression is formulaic; cf.~Sm 39, Sm 44, Sö 130, U 739, and U 805.  For “\inx[C]{good of meat}”, which also occurs in \textlink{Havamal} see Index.  The first line is not poetic.

The carver notably struggles with distinguishing the mirrored runes \textbf{a} and \textbf{n}.

\subsection{Runes}

TODO: Runes%textrunic

\subsection{Transliteration}

\textbf{: -sui · let [·] rnisa · stain [·] þisa · at · anulf · sun · sen · koþan ·· han · byki · her ...- · mantr · matar · -o-r · auk · mls · ri(s)ia ·}

\subsection{Text}

\bpg\bpa[0]%
\emph{Ǫ̇}s·vi lét ręisa stęin þe\emph{nn}sa at Ǫ\emph{r}n·ulf, son si̇nn góðan. Hann byggi hér \edtext{...}{\Bfootnote{In this lacuna there is room for around 5 runes; not enough for a half-line, which suggests that \emph{Hann byggi hér} is not poetic.}}, \\
\alst{m}andr \alst{m}atar \emph{g}ó\emph{ð}r \hld\ ok \alst{m}\emph{á}ls risi\emph{nn}.\epa

\bvb Oswye let raise this stone after Arnolf, her good son. He lived here ..., \\
a man good of meat and proud of speech.\evb\evg

\sectionline\clearpage


\section{U 707/710 (Husby)}

\begin{flushright}%
\textbf{Dating:} C11th

\textbf{Meter:} \Fornyrdislag
\end{flushright}%

\subsection{Introduction}

Two pieces of a single stone.  U 707 is the larger fragment while U 710, originally not identified with U 707, is smaller and contains the words \textbf{(h)...(n) · byki| |i · h(u)}.

\subsection{Runes}

TODO: Runes%textrunic

\subsection{Transliteration}

U 707/U 710: \textbf{· ysurkʀ -{}-... · stan · þina · nt · ka(u)bm...(n) · faþur · sen · koþan · (h)...(n) · byki| |i · h(u)... ¶ · bali · risti · r(u)-{}-r · þis-(r) · hier · man · stanta · stan · ne(r) (b)(r)-...-u ·}

\subsection{Text}

\bpg\bpa[0]%
TODO: text\epa

\bpb TODO: trans\epb\epg

\sectionline\clearpage


\section{U 729}

\begin{flushright}%
\textbf{Dating:} C11th

\textbf{Meter:} \Fornyrdislag
\end{flushright}%

\subsection{Introduction}

TODO.

\subsection{Runes}

TODO: Runes%textrunic

\subsection{Transliteration}

\textbf{· uiþugsi · lit · raisa · stain · þiasn · iftiʀ · seref · faþur · sen · koþan · han · byki| |i · agurstam · hier mn · stanta · stan · miþli · bua · raþi · tekr · þaʀ · ryn si · runum · þim sum · bali · risti ·}

\subsection{Text}

\bpg\bpa[0]%
TODO: text\epa

\bpb TODO: trans\epb\epg

\sectionline\clearpage


\section{U 735 (Långarnö)}

\begin{flushright}%
\textbf{Dating:} C11th

\textbf{Meter:} \Fornyrdislag
\end{flushright}%

\subsection{Introduction}

TODO.

The Wetherold on this stone has been connected to the man of the same name appearing on the unmetrical U 463, located some 27 km east of U 735.  Since both stones are at their original locations, the close distance supports a possible relation.

\subsection{Runes}

TODO: Runes%textrunic

\subsection{Transliteration}

\textbf{· ueþralti ’ lit ’ uʀ · lakarni · stan · almikin · uʀ · staþi · fyra · auk ·· arker ’ þau ’ litu ’ kubl ’ raisa · þisa · at ’ siktryk ’ sun · sen ’}

\subsection{Text}

\bpg\bpa[0]%
Veðr·aldi \alst{l}ét \hld\ úʀ \alst{L}ang-garni \\
\alst{st}ęin al-mikinn \hld\ úʀ \alst{st}aði fǿra \\
ok Ar\emph{n}·gę\emph{rðr} þau létu kumbl ręisa þessa at Sig·trygg, son si̇nn.\epa

\bpb TODO: trans\epb\epg

\sectionline\clearpage


\section{U 739 (Gådi)}

\begin{flushright}%
\textbf{Dating:} C11th

\textbf{Meter:} \Fornyrdislag
\end{flushright}%

\subsection{Introduction}

A standing stone inscribed on one side, with a large cross present.  There are no major difficulties with the reading.  The expression is formulaic; cf.~Sm 39, Sm 44, Sö 130, U 703, and U 805.  “mild of meat” appears to be a variant of “\inx[C]{good of meat}”, which also occurs in \textlink{Havamal}; see Index.  The first line is not poetic.  For other stones raised by someone in memory of themselves, see TODO.

\subsection{Runes}

TODO: Runes%textrunic

\subsection{Transliteration}

\textbf{· hulbiorn · lit · raisa · stan · at ’ sik · sialfan ¶
· hau|n| |u|naʀ · miltr · mataʀ · auk · mals·risin ·}

\subsection{Text}

\bvg\bva[]%
Hol·biǫrn lét ręisa stęin at sik sialfan. &
\edtext{Ha\emph{nn v}aʀ}{\Afootnote{\emph{\textbf{haunaʀ}} (for \emph{\textbf{*hanuaʀ}) U 739}}} \alst{m}ildr \alst{m}ataʀ \hld\ ok \alst{m}áls-risinn.\eva

\bvb Holbarn had this stone raised after himself. \\
He was mild of meat and proud of speech.\evb\evg

\sectionline\clearpage


\section{U 778}

\begin{flushright}%
\textbf{Dating:} C11th

\textbf{Meter:} \Fornyrdislag
\end{flushright}%

\subsection{Introduction}

The original location of the stone is unknown, but judging from the inscription it must have been part of a larger monument consisting of several stones.

\subsection{Runes}

TODO: Runes%textrunic

\subsection{Transliteration}

\textbf{þialfi × auk × hulmnlauk × litu × raisa × staina þisa × ala × at baka × sun sin × is ati × ain ¶ × sir × skib × auk × austr × stu[rþi ×] i × ikuars × liþ × kuþ hialbi × ąt × baka × ask(i)l × raist}

\subsection{Text}

\bpg\bpa[0]%
TODO: text\epa

\bpb TODO: trans\epb\epg

\sectionline\clearpage


\section{U 792}

\begin{flushright}%
\textbf{Dating:} C11th

\textbf{Meter:} \Fornyrdislag
\end{flushright}%

\subsection{Introduction}

TODO.

\subsection{Runes}

TODO: Runes%textrunic

\subsection{Transliteration}

\textbf{kar lit · risa · stin · þtina · at · mursa · faþur · sin · auk · kabi · at · mah sin · fu- hfila · far · aflaþi ut i · kri[k]um · arfa · sinum}

\subsection{Text}

\bpg\bpa[0]%
TODO: text\epa

\bpb TODO: trans\epb\epg

\sectionline\clearpage


\section{U 805}

\begin{flushright}%
\textbf{Dating:} C11th

\textbf{Meter:} \Fornyrdislag
\end{flushright}%

\subsection{Introduction}

The stone is lost and the interpretation must rely on old depictions which makes the reading, especially two of the personal names, uncertain.

The expression is formulaic; cf.~Sm 39, Sm 44, Sö 130, U 703, and U 739.  For “\inx[C]{good of meat}”, which also occurs in \textlink{Havamal}; see Index.  The first line is not poetic.

\subsection{Runes}

TODO: Runes%textrunic

\subsection{Transliteration}

§ P \textbf{[fulkir : lit : risa : st... -{}-r : iel : bruþur : sin : auk : kunimar : eftir : menk : faþur : sin bunta : kuþan : matar : buki| |i uikby]} \\
§ Q \textbf{[fulkir · lit : risa : st... -{}-{}- : mer : bruþur : sin : auk : kuni:mar : eftir : m·enk : faþur : sin bunta : kuþan : matar : brike -ukni ·]} \\
§ R \textbf{[fulkir : lit : risa : st... ...(r)ir : mei- : bruþur : sih auk : kuni:mar : eftir : mink : faþur : sin bunta : kuþan : matar : buki| |i : uiki]}

\subsection{Text}

\bpg\bpa[0]%
\startsub Fylkir lét ręisa st\emph{ęin epti}r \textbf{†iel†}, bróður si̇nn, ok Gunni·marr eptir \textbf{†menk†}, fǫður si̇nn,\endsub \\
\alst{b}ónda góðan matar; \hld\ \alst{b}yggi i̇́ Vík-bý.\epa

\bpb Filch had this stone raised after ..., his brother, and Guthmar after ..., his father, \\
a farmer good of meat; he lived in Wickby.\epb\epg

\sectionline\clearpage


\section{U 838}

\begin{flushright}%
\textbf{Dating:} C11th

\textbf{Meter:} \Fornyrdislag
\end{flushright}%

\subsection{Introduction}

TODO.

\subsection{Runes}

TODO: Runes%textrunic

\subsection{Transliteration}

\textbf{þufr · auk · þorfatr · þair · litu ’ raisa ’ stan · at · þorborn · faþur · sen · koþan · ¶ · hir maa ’ stanta ’ stain ’ ner ’ brautu ’ auk ’ (k)ilauk ’ riþ · kirua ’ merki ’ [at] (b)(o)a-{}-{}- · sen :}

\subsection{Text}

\bpg\bpa[0]%
TODO: text\epa

\bpb TODO: trans\epb\epg

\sectionline\clearpage


\section{U 887}

\begin{flushright}%
\textbf{Dating:} C11th

\textbf{Meter:} \Fornyrdislag
\end{flushright}%

\subsection{Introduction}

TODO.

\subsection{Runes}

TODO: Runes%textrunic

\subsection{Transliteration}

\textbf{× iokeʀ × auk × fastkeʀ × auk × au(r)(i)(k)r × litu · rita · sten × yftiʀ × borkeʀ × faþur sin frehn ¶ selfr · hiok × aurikr × yfti : sin · faþur · runi · þisa · roþi ’ sa ’ kuni}

\subsection{Text}

\bpg\bpa[0]%
\startsub Io·gêʀ auk Fast·gêr auk Ęy·ríkr létu rétta stên øpti Borg·gêʀ fǫður si̇nn frę́gan.\endsub \\
\alst{S}ialfr hiogg Ęy·ríkr \hld\ øpti \alst{s}i̇nn fǫður \\
\alst{r}u̇ni þęssa: \hld\ \alst{r}áði sá kunni.\epa

\bpb TODO: trans\epb\epg

\sectionline\clearpage


\section{U 1016}

\begin{flushright}%
\textbf{Dating:} C11th

\textbf{Meter:} \Fornyrdislag
\end{flushright}%

\subsection{Introduction}

TODO.

\subsection{Runes}

TODO: Runes%textrunic

\subsection{Transliteration}

§ P \textbf{· liutr : sturimaþr · riti : stain : þinsa : aftir : sunu · sina : sa hit : aki : sims uti furs : sturþ(i) · -(n)ari · kuam ·: hn krik·:hafnir : haima tu : ...-mu-... ...(k)(a)(r)... (i)uk (r)(u)-(a) · ...} \\
§ Q \textbf{· liutr : sturimaþr · riti : stain : þinsa : aftir : sunu · sina : sa hit : aki : sims uti furs : sturþ(i) · -(n)ari · kuam ·: hn krik · : hafnir : haima tu : ...-mu-... ...(k)(a)(r)... (i)uk (r)(u)-(a) · ...}

\subsection{Text}

\bpg\bpa[0]%
TODO: text\epa

\bpb TODO: trans\epb\epg

\sectionline\clearpage


\section{U Fv1983;228}

\begin{flushright}%
\textbf{Dating:} C11th

\textbf{Meter:} \Fornyrdislag
\end{flushright}%

\subsection{Introduction}

TODO.

\subsection{Runes}

TODO: Runes%textrunic

\subsection{Transliteration}

\textbf{× nesbiurn : (l)it : raisa : sta-n : þensa : eftiʀ : (o)laif : sun : sin : hiar : skal : stanta : stain · a : l(a)(i)...au...}

\subsection{Text}

\bpg\bpa[0]%
TODO: text\epa

\bpb TODO: trans\epb\epg

\sectionline\clearpage


\section{U Fv1986;84}

\begin{flushright}%
\textbf{Dating:} C11th

\textbf{Meter:} \Fornyrdislag
\end{flushright}%

\subsection{Introduction}

TODO.

\subsection{Runes}

TODO: Runes%textrunic

\subsection{Transliteration}

\textbf{asmu-tr ... ris-{}- · runaʀ · eftiʀ × stein · faþurs·faþur · sin · auk · faþur · siba · ok × geiʀbiarnaʀ × aok ... ulfs · eaʀ · merki · mikit · at · man · koþan ×}

\subsection{Text}

\bpg\bpa[0]%
TODO: text\epa

\bpb TODO: trans\epb\epg

\sectionline\clearpage


\section{U Fv1992;157}

\begin{flushright}%
\textbf{Dating:} C11th

\textbf{Meter:} \Fornyrdislag
\end{flushright}%

\subsection{Introduction}

TODO.

\subsection{Runes}

TODO: Runes%textrunic

\subsection{Transliteration}

\textbf{× kunar : auk biurn : auk × þurkrimr × ra-... ...tain : þina · at þurst... × bruþur sin : is uas austr : tauþr · m... ...ari × auk × karþ... ...u þisi}

\subsection{Text}

\bpg\bpa[0]%
TODO: text\epa

\bpb TODO: trans\epb\epg

\sectionline\clearpage


\section{U Fv1912;8}

\begin{flushright}%
\textbf{Dating:} C11th

\textbf{Meter:} \Drottkvett
\end{flushright}%

\subsection{Introduction}

U Fv1912;8 is a small box, in which were originally found some Wiking Age scales along with a powderized organic material, probably the pouch in which the scales were originally contained.  The inscription has the rare honour of containing the only authentically East Norse poetry in the \Drottkvett.

\subsection{Runes}

TODO: Runes%textrunic

\subsection{Transliteration}

§ A \textbf{tiarfr × fik af × simskum × mąni × skalaʀ × þis... (×) i ...(ą)t(i) × in uirmuntr × f(a)þi × runąr × þisar} \\
§ B \textbf{fuhl × ualua × slait × (f)aluąn × fąn kauk × ą nąs au-a}

\subsection{Text}

\bpg\bpa[0]%
\startsub Djarfr fekk af sęmskum manni skálaʀ þęss\emph{aʀ} i̇́ \dots-\emph{l}andi, en Vermundr fáði ru̇naʀ þęssaʀ.\endsub \\
\alst{F}ugl \edtrans{vęlva}{robber}{\Bfootnote{A hapax in the entire Old Norse corpus, the word appears to be a cognate with Gothic \emph{wilwa} ‘robber, bandit’ which translates Greek \textgreek{ἅρπαξ} \parencite[175]{Streitberg1910}.}} slęit \alst{f}ǫlvan \hld\ \alst{f}ann \edtext{-gauk ȧ nás-}{\lemma{-gauk \dots\ nás- ‘corpse-cuckoo’}\Bfootnote{Taken as a compound \emph{nás-gaukr} ‘corpse-cuckoo \ken{raven/eagle}’.}} \edtrans{au\emph{k}a}{swelling}{\Bfootnote{Viz. from the robber’s flesh.}}.\epa

\bpb Derve got these scales from a Sambian man in \dots-land, but Wermund carved these runes. \\
The fowl tore apart the death-pale robber; one observed the swelling on the corpse-cuckoo \ken{raven/eagle}.\epb\epg

\sectionline\clearpage


\section{Vg 32}

\begin{flushright}%
\textbf{Dating:} C11th

\textbf{Meter:} \Fornyrdislag
\end{flushright}%

\subsection{Introduction}

TODO.  The poetic nature of the inscription is suggested by the unusual word-order, as the verb \emph{rêsþu} ‘raised’ is bracketed by the object \emph{þęnna \dots\ stên} ‘this stone’.

\subsection{Runes}

TODO: Runes%textrunic

\subsection{Transliteration}

\textbf{þurþr × uk : þurun × þana × risþu × stin × efti × era × alkuþan : trik}

\subsection{Text}

\bvg\bva[]%
\alst{Þ}ȯrðr ok \alst{Þ}ȯr·unn \hld\ \alst{þ}ęnna rêsþu &
stên ępti \alst{Ę}rra \hld\ \alst{a}ll-góðan dręng.\eva

\bvb Thurfrith and Thurun raised this stone \\
after Er, an all-good dreng.\evb\evg

\sectionline\clearpage


\section{Vg 59}

\begin{flushright}%
\textbf{Dating:} C11th

\textbf{Meter:} \Fornyrdislag
\end{flushright}%

\subsection{Introduction}

A large standing stone.  The arrangement of the lines makes the division into three parts clear.  First the prose memorial formula, then the first two lines (Eesie’s poem) and finally the raisers’ line.

\subsection{Runes}

TODO: Runes%textrunic

\subsection{Transliteration}

\textbf{rifnikʀ : auk : kiali : auk : brunulfʀ : auk : kifulfʀ : satu : stin : þąnsi : iftiʀ : fut : faþur : sin : harþa : kuþąn : þign ¶
: sua : hifiʀ : ąsa : as : igi : mun : sum : kuin : ift : uir : siþąn : kaurua : ¶
: hialmʀ : auk : hiali : hiaku : runaʀ ·}

\subsection{Text}

\bpg\bpa[0]%
\startsub \emph{H}ręfningʀ auk Gialli auk Bryn·ulfʀ auk Gef·ulfʀ sǫttu stên þannsi ęptiʀ Fót, fǫður si̇nn, harða góðan þegn.\endsub\epa

\bpb Rafning and Yeller and Byrnulf and Yevulf set this stone after Foot, their father, a very good thane.\epb\epg

\bvg\bva[][2]%
Svá hafiʀ \alst{Ȧ}sa, \hld\ es \alst{ê}gi mun &
\alst{s}um kvę̇n ępt vęr \hld\ \alst{s}íðan gǫrva. &
\alst{H}ialmʀ ok \alst{H}ialli \hld\ \alst{h}iǫggu ru̇naʀ.\eva

\bvb So has Eesie [made a monument], like never a wife \\
henceforth after her husband will make. \\
Helm and Helle hewed the runes.\evb\evg

\sectionline\clearpage


\section{Vg 61 (Härlingstorp)}

\begin{flushright}%
\textbf{Dating:} C11th

\textbf{Meter:} \Fornyrdislag
\end{flushright}%

\subsection{Introduction}

TODO.

\subsection{Runes}

TODO: Runes%textrunic

\subsection{Transliteration}

\textbf{: tula : sati : sten : þ... ...[iʀ kʀ : sun] : sin : harþa × kuþąn : trąk : sa × ¶ uarþ : tuþr : ą : uastr:uakm : i : uikiku :}

\subsection{Text}

\bpg\bpa[0]%
TODO: text\epa

\bpb TODO: trans\epb\epg

\sectionline\clearpage


\section{Vg 187}

\begin{flushright}%
\textbf{Dating:} C11th

\textbf{Meter:} \Fornyrdislag
\end{flushright}%

\subsection{Introduction}

A large standing stone decorated with a cross.  The runes are well preserved.

\subsection{Runes}

TODO: Runes%textrunic

\subsection{Transliteration}

\textbf{+ giʀi · sati · stin · þana · eftiʀ · kuþa · bruþur · sin · eʀ · a ąk·lanti · altri · tynþi ×}

\subsection{Text}

\bpg\bpa[0]%
\startsub Gêʀi satti stên þęnna ęptiʀ Goða, bróður si̇nn,\endsub \\
eʀ ȧ \alst{Ę}ng-landi \hld\ \alst{a}ldri tẏnði.\epa

\bpb Gare set this stone after Gid, his brother, \\
who in England lost his life.\epb\epg

\sectionline\clearpage


\section{Vs 15 (Kyringe)}

\begin{flushright}%
\textbf{Dating:} C11th

\textbf{Meter:} \Fornyrdislag
\end{flushright}%

\subsection{Introduction}

A stone with two large flat sides.  The main inscription is § A, whereas the versified § B is less elaborate.

The carver Balle the Red (\emph{Balli hinn rauði}) is also responsible for Vs 24 Hassmyra, where he calls himself \emph{Rø̂ð-Balliʀ} ‘Red-Baller’.  The unique orthography, shared ornamental style, and geographical proximity of the two stones clearly confirm their common authorship.

A notable orthographic aspect is that the carver appears to employ the short-stave forms of the runes \textbf{a} (ᛆ) and \textbf{n} (ᚿ) for the less elaborate parts of the inscription, namely his signatures on both sides, and the whole of side § B.  On the other hand, the main inscription of side § A consistently uses the long-stave forms ᛅ and ᚾ.

\subsection{Runes}

§ A ᚼᛁᚵᚢᛚᚠᚱ × ᛚᛁᛏ × ᚴᛁᚱᛅ × ᛘᛂᚱᚴᛁ × ᚦᛁᛋᛅ × ᛅᛏ × ᚴᚾᚢᛏ × ᚠᛅᚦᚢᚱᛋᛅᚾ × ᚢᚴ × ᚴᚢᚦᛚᚢᚵ · ᚢ × ᚴ ᛋᛁᛋᛏᚢᚱ × ᚼᛅᚾᛁᛋ × ᛒᛅᛚᛁ × ᚱᛁᛏᛁ × ᛁᛋᛏᛅᛂᚾ \\
× ᛒᛆᛚᛁ ×× ᚴᚿᚢᛏᚱ × \\
§ B × ᛋᛏᛂᚿ ᛬ ᚼᛆᚠᛁᚱ × ᚱᛁᛏᚯᚿ × ᚦᚯᚿ × ᛋᛏᚯᚿᛏᛆ × ᛘᚯ × ᛒᛆᛚᛁᚼᛁ[ᚿ]ᚱᛆᚢᚦᛁ × ᚤᚠᛏᛁᚱ × ᛒᚱᚢᚦ[ᚢ]ᚱ × \\
ᛒᛆᛚᛁ
%textrunic

\subsection{Transliteration}

§ A \textbf{× higulfr × lit × kira × merki × þisa × at × ¶ knut × faþur san × uk × kuþlug · u×k sistur × hanis × bali × riti × istaen ¶ × bali ×× knutr ×} \\
§ B \textbf{× sten : hafir × ritąn × þąn × stąnta × mą × bali (h)i- rauþi × yftir × bruþ[u]r × ¶ bali}

\subsection{Text}

\bpg\bpa[0]%
\startsub Hęg·ulfʀ lét gęra męrki þęssa at Knút, fǫður s\emph{i}nn, ok Guð·lø̂g, \edtext{*}{\Afootnote{\textbf{\emph{u×k}} ‘and’ \textup{Vs 15}}} systur \edtext{han*s}{\Afootnote{\textbf{\emph{hanis}} \textup{Vs 15}}}. Balli rétti *stęin.\endsub \\
\edtrans{Balli. Knútr.}{Balle.  Cnut.}{\Afootnote{carved outside of the band \textup{Vs 15}}}\epa

\bpb Hawwolf had these landmarks made after Cnut, his father, and Guthleigh, his sister.  Balle erected the stone. \\
Balle.  Cnut.\epb\epg


\bvg\bva[][3]%
\alst{St}ên hafiʀ réttan, \hld\ \edtrans{þann \alst{st}anda má}{which shall stand}{\Bfootnote{Implicitly: stand through the ages.}}, &
\alst{B}alli hi\emph{nn} rauði, \hld\ øptiʀ \alst{b}róður. &
\edtrans{Balli.}{Balle.}{\Afootnote{carved outside of the band \textup{Vs 15}}}\eva

\bvb The stone has been erected, which shall stand, \\
by Balle the Red after his brother. \\
Balle.\evb\evg

\sectionline\clearpage


\section{Vs 24 (Hassmyra)}

\begin{flushright}%
\textbf{Dating:} C11th

\textbf{Meter:} \Fornyrdislag
\end{flushright}%

\subsection{Introduction}

\textbf{Vs 24} or \textbf{Hassmyrastenen} (‘the Hazemoor Stone’) is an 11th century stone belonging to the Central Swedish runestone tradition.  The stone would originally have been placed at the farm Hassmyra (cf.~l.~2 n.), but today stands at the local village outdoors museum (\emph{hembygdsgård}) in Fläckebo, around a kilometer to the west.

%TODO: physical description

All carvings on the stone are found on the single large, flat side.  In the center, taking up most of the flat side’s surface, is a large great beast (probably to be identified as a hart) locked in battle with a serpent, a common motif in the Urnes style, though likely with much older roots (cf.~\textlink{Fafnismal}[2] n.).  No Christian symbolism is found, and in combination with the name \emph{Óðin-Dísa} ‘Weden-Dise’ (see note), this gives the stone a pagan impression.

The inscription is found inside two bands.  The first fully encloses the great beast and the serpent with which it fights.  The band is not a circle, but rather has two terminating ends at the bottom of the stone.  The left end has a tail, while the right has a serpentine head; the head and tail are bound together by a sort of chain-link or fetter, but the serpent is pulling its head away from the chain-link and appears to be struggling to escape.  Outside of the first band on the left half of the stone, there is another band with no ornamentations.  The second band closely follows the edge of the stone.  Part of its text has been obliterated due to the loss of part of the stone.

The inscription itself begins in the bottom right of the first band, by the serpent’s head.  It thereafter continues through the first band until it runs out of space.  After this 7 runes (\textbf{byi raþr}) are written on the inside of the band near the serpent’s tail, apparently due to lack of space.  The inscription continues in the second band, where it begins at the bottom and runs to the top of the stone, where it terminates.

Ortohgraphically, the os-rune is used for both the oral vowels \emph{o} (\textbf{oþintisu}) and \emph{ø} (\textbf{hulmkoetr, roþbalir}) and for nasal-adjacent \emph{a} (\textbf{buonti, seno, þon}).

The stone is signed by the carver Balle the Red (\emph{Rø̂ð-Balliʀ}), otherwise known from Vs 15 Kyringe.  It is uncertain whether Balle the Red is to be identified with Balle, a prolific carver chiefly active in western Uppland.

\subsection{Runes}

ᛒᚢᚬᚿᛏᛁ × ᚴᚢᚦᚱ × ᚼᚢᛚᛘᚴᚬᛂᛏᚱ × ᛚᛁᛏ × ᚱᛂᛋᛆ × ᚢᚠᛏᛂᛦ × ᚬᚦᛁᚿᛏᛁᛋᚢ × ᚴᚢᚿᚢ × ᛋᛂᚿᚬ × ᚴᚢᛘᛒᚱ × ᚼᛁᚠᚱᚤᛆ × ᛏᛁᛚ × ᚼᛆᛋᚢᛁᛘᚢᚱᛅ × ᛁᚴᛁᛒᛂᛏᚱ × ᚦᚬᚿ ᛒᚤᛁᚱᛆᚦᚱ \\
ᚱᚬᚦᛒᛅᛚᛁᚱ × ᚱᛁᛋᛏᛁ × ᚱᚢᚿᛁ × ᚦᛁᛋᛆ × ᛋᛁᚴᛘᚢᚿᛏᛆᛦ × ᚢᛆᛦ ... ᛋᛂᛋᛏᛦ × ᚴᚢᚦ

%textrunic

\subsection{Transliteration}

\textbf{buąnti × kuþr × hulmkąetr × lit × resa × ufteʀ × ąþintisu × kunu × seną × kumbr × hifrya × til × hasuimura × ikibetr × þąn × byiraþr} \\
\textbf{rąþbalir × risti × runi × þisa × sikmuntaʀ × uaʀ ... sestʀ × kuþ}

\subsection{Text}

\bpg\bpa[0]%
\startsub Bóandi góðr Holm·gø̂tr lét rêsa øptiʀ \edtrans{Óðin-Dísu}{Weden-Dise}{\Bfootnote{This name is unusual for two reasons.  Firstly, the god \emph{Óðinn} ‘Weden’, is poorly attested as the first element (proterotheme) in a dithematic proper name (the only other North Germanic runic example is the Danish name \emph{Óðin·kǫ́rr}, \cite[171]{Peterson2007}).  Secondly, dithematic proper names are, in general, not allowed to be weak.
The best solution to this problem is to see \emph{Dísa} as the woman’s full name, and \emph{Óðin-} as a prefix rather than a true prototheme.  The prefixing of the god to the shortened form of her name (equivalent to something like “Church-Betty”) would then derive from her particular devotion to the god Weden, something corroborated by the lack of Christian symbols or language on the stone; cf.~\textcite[21]{Vikstrand2009}.}}, konu si̇́na.\endsub \\%TODO: Check Vikstrand2009
Kombr \alst{h}í-fręyja \hld\ til \edtrans{\alst{H}ǫs*u*-mýra}{Hazemoor}{\Afootnote{\emph{\textbf{hasm=uimura}} \textup{Vs 24}}\Bfootnote{The Old Norse gen.~pl.~form of the ancestor of Swedish \emph{Hassmyra}, a place name which remains in use to this day.
The sequence \textbf{m=ui} should be understood as the result of a mistake by the carver, since neither the mediæval form \emph{Hasumyr-} \parencite[179]{Fridell2013} nor the modern Swedish form of the name permit the sequence \emph{-mvi-} in this position.  \textcite{Fridell2013} most reasonably explains the error in the following way: after carving \textbf{has}, the carver left out the \textbf{u} and carved \textbf{hasm} along with the vertical stave of the \textbf{u}, which now appears as an \textbf{i}.
He thereafter corrected \textbf{hasm} to \textbf{hasu} by carving the sloping stave of the \textbf{u} on the vertical stave of the \textbf{m}, but since the vertical stave of the following incomplete \textbf{u} was too close, he left it unpainted and instead carved the following \textbf{m} rune.  When painted, \textbf{hasm=uimura} would thus have read \textbf{hasumura}.  Cf.~Fig.~1 in \textcite[178]{Fridell2013}.}} \\
êgi \alst{b}ętri, \hld\ \edtext{þø̂n}{\Bfootnote{The Old Swedish nom.~sg.~f. demonstrative pronoun, here used as a relative pronoun.  \emph{þø̂n}, which is also found in the eastern Old Swedish laws, is a monophthongised form of \emph{þaun} as found in Old Gutnish, which is probably originally derived by analogy from the neuter plural (Old West Norse \emph{þau}).  Old West Norse and the West-Geatish Law instead retain the older \emph{sú} (cf.~Gothic \emph{só}), though the latter also has \emph{þý} \parenciteshorthand[s.vv.~\textit{þæn, su}]{Schlyter1877}.}} \alst{b}ýi ráðr. \\
\alst{R}ø̂ð-Balliʀ \alst{r}ísti \hld\ \alst{r}u̇ni\emph{ʀ} þęssa\emph{ʀ}. \\
\alst{S}ig·mundaʀ vaʀ \edtext{\emph{Óðin-Dísa}}{\Afootnote{obliterated \textup{Vs 24}}\Bfootnote{A part of the band which would have fit approximately 8--10 runes is chipped off, but the space is perfectly compatible with the name \emph{Óðin-Dísa}, which has thus been inserted.}} \hld\ \alst{s}yst\emph{i}ʀ góð.\epa

\bpb The good farmer Holmgeat had this stone raised after Weden-Dise, his wife. \\
There will come to Hazemoor no better \\
housewife who rules the farmstead. \\
Red-Ball carved these runes. \\
To Syemund Weden-Dise was a good sister.\epb\epg

\sectionline\clearpage


\section{Ög 8 (Kälvesten)}

\begin{flushright}%
\textbf{Dating:} C9th

\textbf{Meter:} \Fornyrdislag
\end{flushright}%

\subsection{Introduction}

A remarkably old stone.  Both orthography and the shape of runes agree very closely with the Ög 136 Rök.

\subsection{Runes}

TODO: Runes%textrunic

\subsection{Transliteration}

§ A \textbf{stikuʀ (’) karþi kubl þ(a)\textasciicircum (u) aft auint sunu sin ’ sa fial austr} \\
§ B \textbf{miʀ aiuisli ’ uikikʀ faþi auk| |krimulfʀ}

\subsection{Text}

\bpg\bpa[0]%
\startsub Styggʀ gęrði kumbl þau apt Ęy·vind, sunu si̇nn.\endsub \\
Sá fiall \alst{au}str \hld\ meðʀ \alst{Ęi}v·ísli. \\
Víkingʀ fáði ok Grím·ulfʀ.\epa

\bpb Stig made these monuments after Eanwind, his son. \\
He fell in the east with Awyils. \\
Wiking and Grimwolf painted.\epb\epg

\sectionline\clearpage


\section{Ög 66}

\begin{flushright}%
\textbf{Dating:} C11th

\textbf{Meter:} \Fornyrdislag
\end{flushright}%

\subsection{Introduction}

The incorrect use of \textbf{h} (\textbf{halkuþan} \emph{*all-góðan}, \textbf{it} \emph{’êt} for \emph{*hêt}) indicates that the sound has been dropped

\subsection{Runes}

TODO: Runes%textrunic

\subsection{Transliteration}

\textbf{: ikiualtr : ristþi : stin : þisi : eftiʀ : stufialt : bruþur : sin : suain : halkuþan : sun sbialbuþa : it : inik : ant}

\subsection{Text}

\bpg\bpa[0]%
\startsub Ingi·valdr ręisti stęin þę\emph{nn}si ęptiʀ Styfj·ald, bróður si̇nn,\endsub \\
\alst{s}vęin *all-góðan, \hld\ \alst{s}un Spiall·buða. \\
’Êt inni’k ęnt.\epa

\bpb Ingold raised this stone after Stivold, his brother, \\
an all-good swain, Spellbid’s son. \\
I declare the promise fulfilled.\epb\epg

\sectionline\clearpage


\section{Ög 68}

\begin{flushright}%
\textbf{Dating:} C11th

\textbf{Meter:} \Fornyrdislag
\end{flushright}%

\subsection{Introduction}

TODO.  Wering may be identified with the man of the same name mentioned on the non-metrical Ög 111.

\subsection{Runes}

TODO: Runes%textrunic

\subsection{Transliteration}

\textbf{suina × karþi × bru × þesi × eftiʀ × ouint × bruþur × sin × han × uas × uesteʀ × tauþeʀ × i × uereks × (k)ai-i}

\subsection{Text}

\bpg\bpa[0]%
\startsub Svęina gęrði bró þęssi ęptiʀ Ęy·vind, bróður si̇nn.\endsub \\
Hann vas \alst{v}estʀ dauðʀ \hld\ i̇́ \alst{V}ę́rings \textbf{†kai-i†}.\epa

\bpb Swanie made this bridge after Eanwind, her brother. \\
He died in the west in the Wering’s ...\epb\epg

\sectionline\clearpage


\section{Ög 77}

\begin{flushright}%
\textbf{Dating:} C11th

\textbf{Meter:} \Fornyrdislag
\end{flushright}%

\subsection{Introduction}

TODO.

\subsection{Runes}

TODO: Runes%textrunic

\subsection{Transliteration}

\textbf{: tuna : sati : stin : þ-... {}-{}-tʀ : sin : uar : þurfast : uas han : man : mist : uniþik ·}

\subsection{Text}

\bpg\bpa[0]%
\startsub Tonna satti stên þ\emph{ęnna ęp}t\emph{i}ʀ si̇nn vęr Þȯr·fast.\endsub \\
\edtrans{Vas hann \edtext{\alst{m}an\emph{na}}{\Afootnote{\emph{\textbf{man} Ög 77}}} \hld\ \edtext{\alst{m}ęst\emph{r}}{\Afootnote{\emph{\textbf{mist} Ög 77}}} \edtext{ȯ·níðing\emph{ʀ}}{\Afootnote{\emph{\textbf{uniþik} Ög 77}}}%
}{He was of men the greatest un-nithing}{\Bfootnote{Found near-identically in Sm 5 and Sm 37.}} \epa

\bpb Tonne set this stone after her husband, Thurfast. \\
He was of men the greatest un-nithing.\epb\epg

\sectionline\clearpage


\section{Ög 81 (Högby)}

\begin{flushright}%
\textbf{Dating:} C11th

\textbf{Meter:} \Fornyrdislag
\end{flushright}%

\subsection{Introduction}

A stone standing at the former site of the mediaeval Högby church, into which it was once bricked in.  The stone is inscribed on 2 sides, A and B.

The orthography is notable, since old \emph{ęi} is consistently written as a monophthong, while old \emph{au} is still written as a diphthong.  This unusual variation is reflected in the standardized orthography.

Thurkettle, the carver, also signed Ög 82, Ög 84 and Ög 165 (Skänninge, see below), and is the likely carver of Ög 83 (TODO: cite Östergötlands runinskrifter).

\subsection{Runes}

TODO: Runes%textrunic

\subsection{Transliteration}

§ A · \textbf{þukir · resþi · stin · þansi · eftiʀ · asur · sen · muþur·bruþur · sin · iaʀ · eataþis · austr · i · krikum} · \\
§ B · \textbf{kuþr · karl · kuli · kat · fim · syni · feal · ą · furi · frukn · treks · asmutr · aitaþis · asur · austr · i krikum · uarþ · ą hulmi · halftan · tribin · kari · uarþ · at uti · ¶ auk · tauþr · bui · þurkil · rist · runaʀ ·}

\subsection{Text}

\bpg\bpa[0]%
\startsub Þor·gęr\emph{ðr} rêsþi stên þannsi ęptiʀ Ǫss·ur sên, móður-bróður si̇nn,\endsub \\
iaʀ \alst{ę}ndaðis \hld\ \alst{au}str i̇́ Grikkjum.\epa

\bpb Thurgerd raised this stone after Answer the slow, her mother’s brother \\
who met his end to the east in Greece.\epb\epg

\bvg\bva[][3]%
\alst{G}óðr karl \alst{G}ulli \hld\ \alst{g}at fimm syni; &
\alst{f}iall ȧ \alst{F}ýri \hld\ \alst{f}rǿkn dręng\emph{ʀ} Ǫ̇s·mundr, &
\alst{ę}ndaðis \alst{Ǫ}ss·urr \hld\ \alst{au}str i̇́ Grikkjum, &
varð ȧ \alst{H}olmi \hld\ \alst{H}alf·dan drepinn, &
Kári varð at \alst{U}ddi, \hld\ \alst{au}k dauðr Bói. &
Þȯr·kell rêst ru̇naʀ.\eva

\bvb The good churl Gull begat five sons. \\
On Fyre fell the brave dreng Osmund; \\
Answer met his end to the east in Greece; \\
upon Holm was Halfden slain; \\
Cheer was [slain] by Orde, also dead is Bo. \\
Thurkettle carved the runes.\evb\evg

\sectionline\clearpage


\section{Ög 136 (Rök)}

\begin{flushright}%
\textbf{Dating:} C11th

\textbf{Meter:} \Fornyrdislag
\end{flushright}%

\subsection{Introduction}

TODO.

Due to its great length, use of ciphers, unique content and damage in one part the Rök stone is certainly the most difficult runestone to interpret and explain.  A full-length treatment could easily take up a 300-page book and therefore the present is only an excerpt of the immediately introductory prose passage followed by the 4-line stanza.

\subsection{Runes}

TODO: Runes%textrunic

\subsection{Transliteration}

{[...]} \textbf{’ þat sakum ąna¶rt huaʀ fur niu altum ą n urþi fiaru ¶ miʀ hraiþkutum auk tu ¶ miʀ ąn ubsakaʀ ¶ raiþ| |þiaurikʀ hin þurmuþi stiliʀ ¶ flutna strąntu hraiþmaraʀ sitiʀ nu karuʀ ą} \\
§ B \textbf{kuta sinum skialti ub fatlaþʀ skati marika}

\subsection{Text}

\bpg\bpa[0]%
\startsub{[...]}
Þat sǫgum annart, hvaʀ fór níu aldum ȧ, ęn urði fiaru meðʀ Hręið-gutum, ok dó meðʀ ȧn of·sakaʀ.\endsub\epa

\bpb This I say secondly, who journeyed upon the nine waves but struck the shore amidst the Reth-Gots and died amidst them without guilt.\epb\epg


\bvg\bva[][2]%
Ręið \alst{Þ}iau·ríkʀ, \hld\ hinn \alst{þ}or-móði &
\alst{st}illiʀ flutna, \hld\ \alst{st}rǫndu Hręið-maraʀ; &
sitiʀ nú \alst{g}ǫruʀ \hld\ ȧ \alst{g}uta si̇́num, &
\alst{sk}ialdi of fatlaðr, \hld\ \alst{sk}ati Mę́ringa.\eva

\bvb Thedric rode, the daring-minded \\
controller of sailors, upon the Rethmare shore. \\
Now he sits ready upon his horse, \\
his shield strapped, the lord of Meerings.\evb\evg

\sectionline\clearpage


\section{Ög 155}

\begin{flushright}%
\textbf{Dating:} C11th

\textbf{Meter:} \Malahattr
\end{flushright}%

\subsection{Introduction}

TODO.

\subsection{Runes}

TODO: Runes%textrunic

\subsection{Transliteration}

\textbf{þurfriþ × risti × eftiʀ × askut × auk × kauta sunu × sina × stin × þasi × han × kuti × etaþis × i × ikuars × hilfniki ×}

\subsection{Text}

\bpg\bpa[0]%
\startsub Þȯr·fríðr rêsti ęptiʀ Ǫ̇s·gaut ok Gauta, sunu sina, stên þa\emph{nn}si.\endsub \\
Hann Gauti \alst{ę}ndaðis \hld\ i̇́ \alst{I}ng·vars hęlfningi.\epa

\bpb Thurfride raised after Osgeat and Geat, her sons, this stone. \\
Geat met his end in Ingwar’s ship-troop.\epb\epg

\sectionline\clearpage


\section{Ög 165 (Skänninge)}

\begin{flushright}%
\textbf{Dating:} C11th

\textbf{Meter:} \Fornyrdislag
\end{flushright}%

\subsection{Introduction}

A tall standing stone from Skänninge.  The carver, Thurkettle, is the same as that of Ög 81; see there.

\subsection{Runes}

TODO: Runes%textrunic

\subsection{Transliteration}

§ A \textbf{(·) þurun · risti · auk × þiʀ × bruþr × suniʀ :} \\
§ B \textbf{× tusta × iftiʀ · sin · faþur} \\
§ C \textbf{× þurkil×k rist × stin : þansi : aufti : tusta :}

\subsection{Text}

\bvg\bva[]%
\alst{Þ}ȯr·unn \edtrans{rêsti}{raised [the stone]}{\Bfootnote{The wording is so formulaic that the object can be left out without any injury to the sense.}} \hld\ auk \alst{þ}êʀ brǿðr, &
\alst{s}yniʀ Tosta, \hld\ ęptiʀ \alst{s}i̇nn fǫður. &
Þȯr·kell ’k rêst stên þannsi øpti Tosta.\eva

\bvb Thurun raised [the stone] along with her brothers \\
the sons of Toste, after their father. \\
I, Thurkettle, carved this stone after Toste.\evb\evg

\sectionline\clearpage


\section{Ög 212}

\begin{flushright}%
\textbf{Dating:} C11th

\textbf{Meter:} \Fornyrdislag
\end{flushright}%

\subsection{Introduction}

The stone is lost. TODO.

\subsection{Runes}

TODO: Runes%textrunic

\subsection{Transliteration}

\textbf{[tula : risti : stina : iftiʀ · buka : auk · bru þa:si : ai : buna · þn]}

\subsection{Text}

\bvg\bva[]%
Tóla rêsti stêna ęptiʀ Bukka &
auk \alst{b}ró þási \hld\ a\emph{t} \alst{b}ón\emph{d}a þ\emph{a}nn.\eva

\bvb Toole raised these stones after Bucke, \\
and this bridge after that husband.\evb\evg

\sectionline\clearpage


\section{Öl 1 (Karlevi)}

\begin{flushright}%
\textbf{Dating:} C10th--C11th

\textbf{Meter:} \Drottkvett
\end{flushright}%

\subsection{Introduction}

Although technically outside of the scope of the present edition by virtue of its being a Scaldic composition in the \Drottkvett, the Karlevi runestone is of such importance that its inclusion was felt imperative.

\subsection{Runes}

TODO: Runes%textrunic

\subsection{Transliteration}

\textbf{+ s-a... -{}-(s)- i(a)s · satr · aiftir · si(b)(a) · kuþa · sun · fultars · in hąns ·· liþi · sati · at · u · (t)ausaiþ?... +: fulkin : likr : hins : fulkþu : flaistr (:)· uisi · þat · maistar · taiþir : tulka · þruþar : traukr : i : þaimsi · huki · munat : raiþ:uiþur : raþa : ruk:starkr · i · tanmarku : -{}-ntils : iarmun··kruntar : urkrąntari : ląnti}

\subsection{Text}

\bpg\bpa[0]%
\startsub Stęinn sási ias sattr ęptir Sibba góða son Foldars, en hans liði satti at ø̂ \textbf{†(t)ausaiþ?...†}\endsub\epa

\bpb This stone is set after Sib the Good, son of Folder, but his retinue set it upon the island \dots\epb\epg


\bvg\bva[][2]%
\alst{F}olginn liggr hinn’s \alst{f}ylgðu, · ---\alst{f}lęstr vissi þat---męstar &
\alst{d}ę́ðir, \alst{d}olga Þrúðar · \alst{d}raugr, i̇́ þęimsi haugi; &
mun-at \alst{r}ęið-Viðurr \alst{r}áða, · \alst{r}óg-starkr, i̇́ Dan-mǫrku, &
\emph{\alst{Ę}}ndils \alst{jǫ}rmun-grundar, · \alst{ø}r·grandari, landi.\eva

\bvb Hidden lies the one whom the greatest deeds followed---most men knew that---the tree of the enmities of Thrith \ken{battles > warrior} within this mound. \\
No more upright strife-strong Wither of the chariot of Andle’s immense ground \ken{sea > ship > seafarer} will wield land in Denmark.\evb\evg

\sectionline\clearpage


\section{Öl 39 (Bägby)}

\begin{flushright}%
\textbf{Dating:} C11th

\textbf{Meter:} \Fornyrdislag
\end{flushright}%

\subsection{Introduction}

TODO.

\subsection{Runes}

TODO: Runes%textrunic

\subsection{Transliteration}

\textbf{× sua(i)(n) kiarþi × iftiʀ × sin × faþur × uikar × sun × ainiki × sialfʀ × raisti × stain ×}

\subsection{Text}

\bvg\bva[]%
Svęinn giarði ęptiʀ si̇nn fǫður, Vík·ar, &
\alst{s}unn ęin⸗igi \hld\ \alst{s}ialfr ręisti stęin.\eva

\bvb Swain made [this] after his father, Wicker; \\
the only son raised the stone by himself.\evb\evg

\sectionline\clearpage

\section{Index of proper nouns}

TODO.

\section{Index of carvers}

TODO.


% \bookStart{Younger Runic Poetry from Norway}

% \input{books/runic/Norway.tex}

\bookStart{Runic Poetry from Denmark and Scania}

%Runic poetry from Denmark and Scania.

\section{Introduction}

The signa used for Danish inscriptions vary.  According to the more common system all Danish inscriptions are numbered according to a unitary system based on the Danish corpus edition, which unlike the Swedish one is not divided into provinces.  The online database \emph{Danske Runeindskrifter} (‘Danish Runic Inscriptions’) does however follow a system identical to the Swedish one, and since those signa are easier to work with due to the ability immediately to identify the origin of any given inscription, it has been adopted here, although references to the DR signa are given in parentheses following the origin.

\subsection{Abbreviations of Danish regions}

\begin{itemize}
  \item Bh = Bornholm
  \item MJy = Midtjylland (Middle Jutland)
  \item NJy = Nordjylland (North Jutland)
  \item Sj = Sjælland (Zealand)
  \item SJy = Sønderjylland (South Jutland)
  \item SlB = Slesvig By (the town of Slesvig)
  \item Sk = Skåne (Scania)
  \item Syd = Lolland-Falster

\end{itemize}

\sectionline

\newpage


\section{Bh 33 (DR 379; Nylarsker-sten 1)}

\begin{flushright}%
\textbf{Dating:} C11th--12th

\textbf{Meter:} \Fornyrdislag
\end{flushright}%

\subsection{Introduction}

TODO.

\subsection{Transliteration}

\textbf{sasur · lit · resa · sten · eftiʀ · aluarþ · faþur · sin · truknaþi · han · uti · meþ · ala · sk¶ibara · etki · i · kristr · hab¶(i) ¶ siolu ¶ has ¶ sten · þesi · stai ¶ eftir}

\subsection{Text}

\bpg\bpa[0]%
\startsub Sassurr lét rêsa stên ęptiʀ Hall·varð, fǫður sinn.\endsub \\
Druknaði hann \alst{ú}ti \hld\ með \alst{a}lla skipara. \\
\textbf{†etki†} é Kristr hialpi siólu hans. \\
\alst{St}ênn þęssi \hld\ \alst{st}ái ęptiʀ.\epa

\bpb Saswer had this stone raised after Hallward, his father. \\
He drowned far out with all the shipmen. \\
\dots may Christ always help his soul. \\
May this stone afterwards stand.\epb\epg

\sectionline\newpage


\section{Bh 34 (DR 380; Nylarsker-sten 2)}

\begin{flushright}%
\textbf{Dating:} C11th--12th

\textbf{Meter:} \Fornyrdislag
\end{flushright}%

\subsection{Introduction}

TODO.

\subsection{Transliteration}

\textbf{kobu:suain : raisti : stain : þ(e)na : a(f)tir : bausa : sun : sin : tr(i)... ...n : þan : is : tribin : ua(r)þ : i : (u)(r)ostu : at : ut:la(n)(k)iu : kuþ : tr(u)tin : hi(a)lbi : hans : ont : auk : sata : mikial :}

\subsection{Text}

\bpg\bpa[0]%
\startsub Kǫ́pu-Svęinn ręisti stęin þęnna ęptir Bausa, sun sinn,\endsub \\
\alst{d}ręng góðan, \hld\ þann es \alst{d}repinn varð \\
i̇́ \alst{o}rrostu \hld\ at \alst{Ú}t-lęngju. \\
Guð dróttinn hialpi hans ǫnd auk \edtrans{santa Mikael}{Saint Michael}{\Bfootnote{Something of a patron saint on Bornholm, also invoked in Bh 1, Bh 4, and Bh 49.  Cf.~Syd 11 below.}}.\epa

\bpb Cloak-Swain raised this stone after Beaze, his son, \\
the good dreng who was slain \\
in battle by Utlengie. \\
Lord God and Saint Michael help his spirit.\epb\epg


\section{Bh 42 (DR 387; Vestermarie-sten 5)}

\begin{flushright}%
\textbf{Dating:} C11th--12th

\textbf{Meter:} \Ljodahattr
\end{flushright}%

\subsection{Introduction}

TODO.

\subsection{Transliteration}

\textbf{: asualdi : risti : stein : þinsa : iftʀ : alfar : bruþur : sin : drinr : koþr ¶ : trebin u:syni : auk : skogi : suek : saklausan :}

\subsection{Text}

\bpg\bpa[0]%
\startsub Ǫ̇s·valdi rísti stęin þęnnsa ęptiʀ Alf·ar, bróður sinn,\endsub \\
\alst{d}ręngr góðr \hld\ \alst{d}repinn u̇·sẏni, \\
\ind auk Skógi \alst{s}vęik \alst{s}ak-lausan.\epa

\bpb Oswald carved this stone after Elver, his brother; \\
the good dreng [was] shamefully slain \\
\ind and Skoog betrayed him, guiltless.\epb\epg

\sectionline\newpage


\section{MJy 37 (DR 94; Ålum-sten 1)}

\begin{flushright}%
\textbf{Dating:} C11th

\textbf{Meter:} \Fornyrdislag
\end{flushright}%

\subsection{Introduction}

TODO.

\subsection{Transliteration}

\textbf{tuli : (r)(i)s-(i) : stin : þasi : aft ¶ ikal:t : sun : sin : miuk (:) (k)¶(u)... ...k : þau : mun(u) ¶ mini : m-(r)gt : ¶ iuf (:) þirta :}

\subsection{Text}

\bpg\bpa[0]%
\startsub Tóli rêsþi stên þannsi apt Ingj·ald sun sinn, miok gó\emph{ðan dręn}g.\endsub \\
Þau \alst{m}unu \alst{m}inni \hld\ \alst{m}\emph{ę}rkt ê of \emph{b}irta.\epa

\bpb Toole raised this stone after Ingeld, his son, a very good dreng. \\
TODO.\epb\epg


\sectionline\newpage

\section{MJy 71 (DR 131; Års-sten)}

\begin{flushright}%
\textbf{Dating:} C11th

\textbf{Meter:} \Fornyrdislag
\end{flushright}%

\subsection{Introduction}

TODO.

\subsection{Transliteration}

§ A \textbf{: ąsur : sati : stin : þąnsi : aft : ual:tuka : trutin : ¶ × sin} \\
§ B \textbf{× stin : kuask : hirsi : stąnta : ląki : saʀ : ual:tuka × ¶ × uarþa : nafni}

\subsection{Text}

\bpg\bpa[0]%
Ǫnsurr satti stên þannsi apt Val-Tóka, dróttin sinn. \\
\alst{St}ênn kweðsk hérsi \hld\ \alst{st}anda lęngi; \\
sáʀ \alst{V}al-Tóka \hld\ \edtrans{\alst{v}arða}{cairn}{\Bfootnote{Acc. sg. of \emph{varði} ‘cairn or beacon raised as a landmark’.  The word still survives in Danish \emph{varde} ‘id.’, Faroese \emph{varði} ‘id.’}} nęfni.\epa

\bpb Answer set this stone after Wal-Took, his lord. \\
The stone proclaims it will stand here for long; \\
may it name Wal-Took’s cairn.\epb\epg

\sectionline\newpage


\section{MJy 79 (DR 68; Århus-sten 5)}

\begin{flushright}%
\textbf{Dating:} C10th--11th

\textbf{Meter:} \Ljodahattr
\end{flushright}%

\subsection{Introduction}

A well-preserved stone with a small piece struck off at the lower corner at the very beginning of the inscription.

\subsection{Transliteration}

§ A \textbf{(-)usti × auk × hufi × auk × þir × frebiurn × risþu × stin × þąnsi × eftiʀ × ¶ × ąsur × saksa × filaka × sin × harþa ×} \\
§ B \textbf{kuþan × trik × saʀ × tu × ¶ × mana × mest × uniþikʀ × ¶ saʀ × ati × skib × miþ × arną +}

\subsection{Text}

\bpg\bpa[0]%
\startsub \emph{T}osti auk Hófi þêʀ Frø̂·biorn rêsþu stên þannsi ęptiʀ Ǫnsur saxa, fé-laga sinn, harða góðan dręng.\endsub \\
Sáʀ dó \alst{m}anna \hld\ \alst{m}ęst u̇·níðingʀ; \\
\ind sáʀ \alst{á}tti skip með \alst{A}rna.\epa

\bpb Toste and Hove, they and Freabern raised this stone after Answer the seax, his partner, a very good dreng. \\
He died the greatest un-nithing of men; \\
\ind he owned a ship with Arnie.\epb\epg

\sectionline\newpage


\section{NJy 41 (DR 149; Ydby-sten)}

\begin{flushright}%
\textbf{Dating:} C10th--11th

\textbf{Meter:} \Fornyrdislag
\end{flushright}%

\subsection{Introduction}

The inscription is entirely poetic.  The stone is heavily damaged and divided into several parts but the inscription survives in older depictions.

\subsection{Transliteration}

§ A \textbf{þurkisl : [sati · au(k) (·) suniʀ (·) lifa · i] : s[t]aþ [· þa]¶:nsi :} \\
§ B \textbf{[sti](n) (÷) (u)[ftiʀ · lifa]}

\subsection{Text}

\bvg\bva[]%
Þȯr·gísl \alst{s}atti \hld\ ok \alst{s}yniʀ Lêfa &
i̇́ \alst{st}að þannsi \hld\ \alst{st}ên øptiʀ Lêfa. \eva

\bvb Thuryils set and the sons of Lofe \\
in this place the stone after Lofe.\evb\evg

\sectionline\newpage


\section{Sj 93 (DR 229; Sandby-sten 3)}

\begin{flushright}%
\textbf{Dating:} C11th

\textbf{Meter:} \Fornyrdislag
\end{flushright}%

\subsection{Introduction}

A heavily damaged stone, the verse can be tentatively reconstructed based on Syd 11.

\subsection{Transliteration}

§ A \textbf{sylfa : rest... ... ...i : sbalklusu : eyfti : susur : faþur ... ...(r)þi : bru : þisi : iki : ¶ þurils : brþur : sin :} \\
§ B \textbf{i mun · san... ¶ ...if · uitrik · susi · eʀ · uan · sil...}

\subsection{Text}

\bpg\bpa[0]%
\startsub Sylfa rêst\emph{i stên þannsi} i̇́ Spalk-lø̂su øpti Susur, fǫður \emph{sinn, ok gę}rði bró þęssi ę\emph{pt}i Þȯr·ils, br\emph{ó}ður sinn.\endsub \\
Ê mun \alst{s\emph{t}}an\emph{da, \hld\ með \alst{st}ênn á l}íf, \\
\alst{v}itring súsi, \hld\ eʀ \alst{v}ann Syl\emph{fa}. \epa

\bpb Silve raised this stone in Spalklease after Saswer, his father, and made this bridge after Thuryils, his brother. \\
Ever will stand while the stone has life \\
this memento which Silve accomplished.\epb\epg

\sectionline\newpage


\section{SJy 5 (DR 37; Egtved-sten)}

\begin{flushright}%
\textbf{Dating:} C10th--11th

\textbf{Meter:} Unclear
\end{flushright}%

\subsection{Introduction}

A fragmented stone; the remaining piece is easily read, but a large part is missing.

\subsection{Transliteration}

\textbf{... ... ...at ’ fai(n) [’] (t)u ÷ i suiu ’ raist ¶ ... ...uþiʀ ’ aft ’ bruþur ¶ stain ’ sasi ’ skarni ’ ...}

\subsection{Text}

\bvg\bva[]%
\dots\ at fáinn; dó i̇́ Svíu. &
\alst{R}ęist \emph{iak \alst{r}u̇naʀ,} &
\emph{\alst{b}r}óðiʀ apt \alst{b}róður; &
stęin sási skęrni \dots\eva

\bvb the painted; died in Sweden. \\
I carved the runes, \\
a brother after [my] brother; \\
may this stone TODO \dots\evb\evg

\sectionline\newpage


\section{SJy 14 (DR 40; Randbøl-sten)}

\begin{flushright}%
\textbf{Dating:} C10th

\textbf{Meter:} \Ljodahattr
\end{flushright}%

\subsection{Introduction}

A rare instance of a stone carved by a man for his wife.

\subsection{Transliteration}

\textbf{tufi ÷ bruti ÷ risþi ÷ stin ÷ þansi ÷ aft ÷ lika ÷ ¶ brutia ÷ þiʀ ÷ stafaʀ ÷ munu ÷ ¶ þurkuni ÷ miuk ÷ liki ÷ lifa ÷}

\subsection{Text}

\bpg\bpa[0]%
\startsub Tófi bryti rêsþi stên þannsi apt \edtrans{líka}{helpmate}{\Bfootnote{Or ‘equal, match’.  Although grammatically masculine (the nom. sg. is \emph{líki}) the word can be used for both genders, as here.}} brytja.\endsub \\
\alst{Þ}êʀ stafaʀ \hld\ munu \alst{Þ}ȯr·gunni \\
\ind miok \alst{l}ęngi \edtrans{\alst{l}ifa}{live}{\Bfootnote{Stones and staves ‘living’ is a common expression.}}.\epa

\bpb Tove the bailiff raised this stone after the bailiff’s helpmate. \\
These staves will for Thurguth \\
\ind for a very long time live.\epb\epg

\sectionline\newpage


\section{Sk 5 (DR 263; Skabersjö-fibula)}

\begin{flushright}%
\textbf{Dating:} C8th

\textbf{Meter:} \Fornyrdislag
\end{flushright}%

\subsection{Introduction}

TODO.

\subsection{Transliteration}

§ P \textbf{    ... ʀʀʀʀʀʀʀʀʀʀʀʀʀʀʀʀ -aþi tuk fauka fiaʀ sis in a iak as(u) (þ)ui launat ¶ ... auab-k sua fakat...} \\
§ Q \textbf{    ... (ʀ)(ʀ)ʀʀʀʀʀʀʀʀʀʀʀʀʀʀ (r)aþi (t)(u)k fauka fiaʀ sis in a iak asa þui launat ¶ ... auab-k sua fakat...}

\subsection{Text}

\bpg\bpa[0]%
TODO: text\epa

\bpb TODO: trans\epb\epg

\sectionline\newpage


\section{Sk 42 (DR 337; Valleberga-sten)}

\begin{flushright}%
\textbf{Dating:} C11th

\textbf{Meter:} \Fornyrdislag
\end{flushright}%

\subsection{Introduction}

A stone.  TODO.

\subsection{Transliteration}

§ A \textbf{: suin : auk : þurgutr : kiaurþu : kubl : þisi ¶ eftiʀ : mana ¶ auk · suina} \\
§ B \textbf{kuþ : hialbi : siaul : þeʀa : uel : ian : þeʀ : likia : i : luntunum}

\subsection{Text}

\bpg\bpa[0]%
\startsub Svênn auk Þȯr·gø̂tr giǫrðu kumbl þęssi ęptiʀ Manna auk Svenna. Guð hialpi siǫ́l þêʀa vęl,\endsub \\
ian þêʀ \alst{l}iggja \hld\ i̇́ \alst{L}ondúnum.\epa

\bpb Swayn and Thurgeat made these monuments after Manny and Swenny.  God help their soul well, \\
but they lie in London.\epb\epg

\sectionline\newpage


\section{Sk 51 (DR 279; Sjörup-sten)}

\begin{flushright}%
\textbf{Dating:} C11th

\textbf{Meter:} \Ljodahattr
\end{flushright}%

\subsection{Introduction}

TODO.

The incorrect use of \textbf{h} in \textbf{huptiʀ} \emph{øptiʀ} and the lack of it in \textbf{an} \emph{’ann} (= ON \emph{hann} ‘he’) \textbf{a(f)þi} \emph{’afði} (= ON \emph{hafði} ‘had’) indicates the common East Norse loss of that phoneme.

\subsection{Transliteration}

\textbf{[+ sa]ksi : sati : st[in] : þasi : huftiʀ : ą[s]biurn : (s)in : fil(a)gą ’ ¶ (t)u-a[s : sun :] ¶ saʀ : flu : aki : a[t :] ub:sal(u)m : an : ua : maþ : an : u¶abn : a(f)þi ’}

\subsection{Text}

\bpg\bpa[0]%
\startsub Saxi satti stên þa\emph{nn}si øptiʀ Ǫ̇s·biorn, sinn fé-laga, Tó\emph{k}a* sun.\endsub \\
Sáʀ fló \alst{ê}gi \hld\ at \alst{U}pp-sǫlum, \\
\ind ęn \alst{v}á, með ’ann \alst{v}ǫ́pn ’afði.\epa

\bpb Saxe set this stone after Osbern, his partner, Took’s son. \\
He did not flee at Upsele \\
\ind but fought as long as he had weapons.\epb\epg

\sectionline\newpage


\section{Sk 80 (DR 295; Hällestad-sten 1)}

\begin{flushright}%
\textbf{Dating:} C11th

\textbf{Meter:} \Fornyrdislag
\end{flushright}%

\subsection{Introduction}

TODO.

\subsection{Transliteration}

§ A \textbf{: askil : sati : stin : þansi : ift[iʀ] ¶ : tuka : kurms : sun : saʀ : hulan : ¶ trutin : saʀ : flu : aigi : at : ub:¶:salum} \\
§ B \textbf{satu : trikaʀ : iftiʀ : sin : bruþr ¶ stin : ą : biarki : stuþan : runum : þiʀ :} \\
§ C \textbf{(k)(u)(r)(m)(s) (:) (t)(u)(k)(a) : kiku : (n)(i)(s)(t)[iʀ]}

\subsection{Text}

\bpg\bpa[0]%
\startsub Ǫ̇s·kęll satti stên þannsi ęptiʀ Tóka, Gorms sun, séʀ hollan dróttin;\endsub \\
sáʀ fló \alst{ęi}gi \hld\ at \alst{U}pp-sǫlum. \\
\alst{S}ǫttu dręngaʀ \hld\ ęptiʀ \alst{s}inn bróður \\
\alst{st}ên ȧ biargi \hld\ \alst{st}uððan ru̇num; \\
þêʀ \alst{G}orms Tóka \hld\ \alst{g}ingu nę́stiʀ. \epa

\bpb Oskettle set this stone after Took, Gorm’s son, to him a trusty lord; \\
he did not flee at Upsele. \\
The drengs set after their brother \\
a stone on the hill propped up by runes; \\
they marched closest to Gorm’s Took.\epb\epg

\sectionline\newpage


\section{Sk 81 (DR 296; Hällestad-sten 2)}

\begin{flushright}%
\textbf{Dating:} C11th

\textbf{Meter:} \Fornyrdislag
\end{flushright}%

\subsection{Introduction}

TODO.

\subsection{Transliteration}

\textbf{: ąskautr : ristþi : stin : þansi (:) ¶ (:) (i)ftiʀ : airu : brþur : sin : ian : ¶ : saʀ : uas : him:þiki : tuka : nu : ¶ : skal : statą : stin : ą : biarki :}

\subsection{Text}

\bpg\bpa[0]%
\startsub Ǫ̇s·gautr rêsþi stên þannsi ęptiʀ Ęiru, bróður sinn, ian sáʀ vas hêm-þęgi Tóka.\endsub
Nú skal \alst{st}anda \hld\ \alst{st}ên ȧ biargi.\epa

\bpb Osgeat raised this stone after Ore, his brother, and he was Took’s retainer. \\
Now shall stand the stone on the hill.\epb\epg

\sectionline\newpage


\section{SlB 3 (DR Schl3; Slesvig-pind)}

\begin{flushright}%
\textbf{Dating:} C11th

\textbf{Meter:} \Ljodahattr
\end{flushright}%

\subsection{Introduction}

A rune-inscribed stick containing a puerile curse.

\subsection{Transliteration}

§ A \textbf{runaʀ ÷ iag ÷ risti ÷ a ÷ rikiata ÷ tre sua} \\
§ B \textbf{reþ ÷ saʀ ÷ riki ÷ mo(g)ʀ ÷ asiʀ ÷ a ÷ artagum} \\
§ C \textbf{hulaʀ ÷ auk ÷ bulaʀ ÷ meli ÷ þeʀ} \\
§ D \textbf{ars ÷ sum ÷ magi}

\subsection{Text}

\bvg\bva[]%
\alst{R}u̇naʀ iak \alst{r}ísti \hld\ ȧ \alst{r}íkjanda tré &
\ind svá \alst{r}éð sáʀ \alst{r}íki mǫgr; &
\alst{ę̇}siʀ ȧ \alst{á}r-dǫgum \hld\ hullaʀ ok bullaʀ &
\ind \alst{m}ę́li þéʀ ars sum \alst{m}agi.\eva

\bvb Runes I carve on the powerful wood; \\
\ind so did the powerful boy counsel: \\
The Eese in days of yore, hurlies and burlies, \\
\ind may thy arse speak like thy belly!\evb\evg

\sectionline\newpage


\section{Syd 11 (DR 212; Tillitse-sten)}

\begin{flushright}%
\textbf{Dating:} C11th

\textbf{Meter:} \Fornyrdislag
\end{flushright}%

\subsection{Introduction}

The stone actually has three sides, but § C is by a different carver and dedicated to a different person; as such it ought to be treated as a separate inscription and not included.

\subsection{Transliteration}

§ A \textbf{eskil : sulka : sun : let : res(a) ¶ sten : þena : eft : sialfan ¶ sik · e mun · stanta · meþ · sten ¶ lifiʀ · uitrint · su · iaʀ · uan · eskil} \\
§ B \textbf{kristr · hialbi · siol · hans ¶ aok · santa · migael}

\subsection{Text}

\bpg\bpa[0]%
\startsub Ę̇s·kil, Súlka sun, lét rêsa stên þęnna ępt sialfan sik.\endsub \\
É mun \alst{st}anda, \hld\ með \alst{st}ênn lifiʀ, \\
\alst{v}itrind sú \hld\ iaʀ \alst{v}ann Ę̇s·kil. \\
Kristr hialpi siól hans auk santa Mikael.\epa

\bpb Oskettle, Sulk’s son, had this stone raised after himself. \\
Ever will stand while the stone lives \\
this memento which Oskettle accomplished. \\
Christ and Saint Michael help his soul.\epb\epg

\sectionline


% \bookStart{British Runic Poetry}

% The inscriptions on stone monuments are all highly formulaic; they typically conclude with the sentence \emph{ge·biddaþ þǽr sâwle} ‘pray for the soul’.

Reference: Lilla Kopár (2015), Heroes on the Fringes of the Anglo-Saxon Poetic Corpus: Vernacular Memorial Inscriptions on Stone Sculpture

OE-GB-26 (Falstone)
\emph{s}ęttæ æftær Hrǿþberht
bécun æf\emph{ter} Eomæ \hld\ ge·bidæd þe\emph{r} saule

OE-GB-29 (Great Urswick)
Túnwini sęttæ \hld\ æfter Torohtrǽdæ &
békun æfter his bæurnæ \hld\ ge·bidæs þer saulæ &
lyl þis w\emph{rât.}

OE-GB-82 (Thornhill 2)
Gilswi̇́þ a·rærde \hld\ æft\emph{er} Berhtswi̇́þe
békun ǫn bergi \hld\ ge·bidda(þ) þær saule

OE-GB-108 (Honington Tweezers)
Þec blǿtsig\emph{e} \hld\ bilwit fæddæ\emph{r}
ǫnd werca ge·hwelc \hld\ hefæn ǫnd ecl\emph{as}.

(Dewsbury)
[\dots]rhtæ
bécun æfter beornæ. \hld\ Gi·biddad dær saule.

